\section{Manifesto}
\subsection{Introdução}

O que eu quero abordar agora, é de certa forma, bastante sensível, por isso eu presumo que será um longo texto. Eu sou uma pessoa que não é indígena e que está querendo apoiar a luta indígena. Um obstáculo que se coloca diante de mim é ser impossível que a leitura que eu faça desta causa, não seja uma leitura externa, uma vez que eu me encontro fora deste movimento. O risco que isso acarreta, é que mesmo que não seja minha intenção, eu tenha uma má compreensão dos povos originários e os retrate de forma desfigurada e até mesmo desrespeitosa.

Estou fazendo um grande esforço para evitar cometer esses equívocos. Por anos eu tenho estudado política, economia, arte e cultura. Todo este meu estudo se deu por autores, que assim como eu, também não são indígenas. Então é este conhecimento que eu atualmente carrego comigo, e me é por agora, impossível realizar qualquer análise sem estar sob esta influência.

O que eu posso fazer, e acredito que seja interessante, é buscar construir uma ponte entre este conhecimento que tenho acumulado e a luta indígena que apenas começo a conhecer. Tenho ‘boas intenções’, mas sei que boas intenções não são o suficiente. A posição que eu busco assumir para realizar este diálogo da melhor forma possível, é a de realizar um esforço para que a luta indígena não seja deformada para caber em projetos políticos externos a si. Pelo contrário, eu busco mobilizar o conhecimento que eu carrego de forma que possa ser visto como um complemento a causa indígena, uma ferramenta adicional que talvez possa ser útil para aqueles que de fato vivem essa luta.

Não quero, de modo algum, buscar dar a entender que os povos originários e sua luta precisam encontrar validação em teorias e ideologias que eu trago. Pelo contrário, o público alvo que eu viso com este texto é não tanto os povos originários, que muito provavelmente não precisam ouvir o que eu tenho a falar, mas sim pessoas que como eu não são indígenas, mas queiram compor a luta anticapitalista.

Acredito que nós sim precisamos ouvi-los. Minha intenção aqui, é então construir uma ponte, ainda que imperfeita e altamente criticável, mas que permita que pessoas que se encontrem na mesma situação que eu consiga dar os primeiros para pensar sobre a luta dos povos originários, ainda que carreguem com si uma visão de mundo diferente.

Eu penso que talvez o ideal fosse que nós conseguíssemos deixar para trás toda nossa concepção de mundo para entender os povos originários não a partir de nossas lentes, mas a partir das próprias lentes dos próprios povos originários. Mas eu não sei o quanto é realisticamente possível alcançar essa meta, por mais nobre que seja. Sendo assim, me parece que um diálogo que busque articular de forma honesta aquela bagagem cultural que trazemos do ‘mundo ocidental’ com a luta dos povos originários no qual começamos a nos educar tão recentemente, parece ser frutífero.

Eu estou tentando escolher as melhores palavras, pois eu não quero incentivar uma falsa segregação na sociedade, uma separação artificial entre ‘nós’ e ‘eles’. Porém eu acredito que na forma com que a sociedade se estrutura hoje, existam hoje recortes na sociedade no qual determinados grupos possuem interesses que são próprios deste grupo. Assim faz sentido reconhecermos uma separação na sociedade entre indígenas e não indígenas.


\subsection{Marxismo}

Por questão de honestidade, eu devo declarar que a bagagem teórica que eu carrego se insere principalmente na corrente do marxismo. Eu quero explicar um pouco melhor, de maneira geral, como eu vejo essa corrente política, e por motivos de transparência, quero expor também os conflitos teóricos potenciais que surgem entre o marxismo e a luta indígena. Uma falha gravíssima nesse momento, é a falta de leitura de autores marxistas e indígenas, pretendo corrigir esta falha em um texto futuro. Mas no momento isto se deve ao fato de que sou, pode-se dizer assim, sou um bebê no que se refere a luta indígena e me encontro na leitura dos seus primeiros autores.

Mas começamos então com a questão mais óbvia possível: o que é o marxismo? A verdade é que não há um consenso sobre o que é marxismo. Há diversas vertentes que podem ser mencionadas, e cada vertente propõe uma própria definição sobre o que é marxismo, e evidentemente, a partir desta definição, ela interpreta que a sua própria vertente é possivelmente o ‘verdadeiro’ marxismo.

Ingo Elbe (“Entre Marx, marxismo e marxismos: leituras da teoria de Marx”) vai discutir a existência de três marxismos que divergem por uma questão filosófica: Marxismo Tradicional, Marxismo Ocidental e Nova Leitura de Marx. Dentro da vertente tradicional ainda podemos encontrar uma divisão entre \textit{stalinistas} e \textit{trotskistas}, que tem como um dos principais pontos de divergência a reivindicação das experiências socialistas reais e as implicações que decorrem disto. Se quisermos prosseguir com mais subdivisões, dentro do \textit{stalinismo} (marxismo-leninismo) podemos encontrar mais um leque de diferentes vertentes como maoismo, \textit{hoxhaismo}, etc.

E para piorar, mesmo quando duas vertentes concordam em algum nível (por exemplo: o que define marxismo é o método) elas podem discordar em outro nível (por exemplo: qual exatamente é este método). Eu particularmente não defendo que exista nenhum ‘verdadeiro’ marxismo. O que me interessa é identificar qual vertente melhor interpreta e modificar a realidade.

\begin{quote}
Em matéria de marxismo, a ortodoxia se refere antes e exclusivamente a um método.

(Georg Lukács em História e Consciência de Classe)
\end{quote}

Se entendermos marxismo como uma simples continuidade do trabalho de Marx, todas as vertentes que partem de Marx, em maior ou menor grau, se prestam a esse papel e tornam-se igualmente verdadeiras.

\begin{quote}
A tradição de todas as gerações mortas oprime como um pesadelo o cérebro dos vivos. E justamente quando parecem empenhados em revolucionar-se a si e às coisas, em criar algo que jamais existiu, precisamente nesses períodos de crise revolucionária, os homens conjuram ansiosamente em seu auxilio os espíritos do passado, tomando-lhes emprestado os nomes, os gritos de guerra e as roupagens, a fim de apresentar e nessa linguagem emprestada.

(Karl Marx em O 18 Brumário de Louis Bonaparte)
\end{quote}

Agora se adicionarmos a exigência de que seja coerente com os escritos originais, precisamos refletir sobre mais uma coisa: Todo o escrito de Marx é coerente entre si? Se não,a vertente deve ser coerente com que parte do trabalho de Marx?

A pergunta é mais difícil do que parece. Mas eu tendo a pensar que precisamos reconhecer a existência de incoerências entre diferentes obras do Marx. Marx morreu com mais de 60 anos e por toda sua vida adulta estudou e produziu trabalhos intelectuais. É de se esperar que ao longo do tempo algumas ideias tenham sofrido alteração.

\begin{quote}
Existem ao menos seis exposições da forma-valor, todas diferentes.

(A “nova leitura de marx”: um mapeamento de suas premissas e desenvolvimentos)
\end{quote}

Ainda há uma mudança no foco de estudo de Marx que se reflete no foco das suas obras publicadas, partindo de obras mais puramente filosóficas como “A Sagrada Família” e se movendo na direção de obras que passam a ter a economia política como tema central, como o “O Capital”. É evidente que o quanto de filosofia persiste em suas obras posteriores é também assunto de polêmicas, nem todos concordariam que “O Capital” não é um livro essencialmente de filosofia. Esta discussão atinge tal nível que alguns autores distinguem em Marx a existência de um “Marx esotérico” e um “Marx exotérico”, outros discutem um corte epistemológico entre o “jovem Marx” e o “velho Marx”. E até mesmo o próprio Marx é acusado de simplificar suas formulações teóricas.

Um relato da trajetória de estudos do Marx é dado pelo próprio no prefácio de Para a Crítica da Economia Política.

\begin{quote}
O meu estudo universitário foi o da jurisprudência, o qual no entanto só prossegui como disciplina subordinada a par de filosofia e história. No ano de 1842-43, como redactor da Rheinische Zeitung, vi-me pela primeira vez, perplexo, perante a dificuldade de ter também de dizer alguma coisa sobre o que se designa por interesses materiais. Os debates do Landtag Renano sobre roubo de lenha e parcelamento da propriedade fundiária, a polémica oficial que Herr von Schaper, então Oberprásident da província renana, abriu com a Rheinische Zeitung sobre a situação dos camponeses do Mosela, por fim as discussões sobre livre-cambismo e tarifas alfandegárias proteccionistas deram-me os primeiros motivos para que me ocupasse com questões económicas.

(Karl Marx em Para a Crítica da Economia Política)
\end{quote}

Sobre o corte epistemológico em Marx é possível encontrarmos esta discussão sendo desenvolvida por Althusser.

\begin{quote}
Existem diversos debates no interior do marxismo, mas poucos foram capazes de gerar tantas polêmicas como as que sucederam após a publicação das obras Por Marx e Ler O capital de Louis Althusser (1918-1990). E aqui me refiro, especificamente, ao debate que envolve a tese da divisão do pensamento de Marx em duas etapas: a da juventude, fortemente influenciado por Hegel e Feuerbach; e a da maturidade, em que há uma ruptura de Marx com a influência desses dois filósofos.

(Althusser e as duas etapas do pensamento de Marx: o conceito de “sobredeterminação”)
\end{quote}

O leitor pode pode retomar o artigo “A ‘nova leitura de Marx’: um mapeamento de suas premissas e desenvolvimentos” para ler as críticas sobre as diferenças entre as duas primeiras edições de O Capital.

\begin{quote}
Sobre as diferenças no primeiro capítulo da primeira e segunda edição de O Capital: “A partir da segunda edição (1872), a análise da forma-valor também é “popularizada”, uma vez que o conteúdo da Seção I passa a ser o mesmo da versão anteriormente oferecida no anexo. Como já destacado, com isso se perde a exposição da “forma genérica” e da “forma IV” presentes na primeira edição.

(A “nova leitura de marx”: um mapeamento de suas premissas e desenvolvimentos)
\end{quote}

E também é possível encontrar a distinção entre o Marx esotérico e exotérico no debate sobre a reconstrução do método marxiano no artigo de Ingo Elbe:

\begin{quote}
O debate da reconstrução agora aplicaria a distinção esotérica/exotérica à própria obra de Marx.

(Entre Marx, marxismo e marxismos: leituras da teoria de Marx)
\end{quote}

Ou seja, entendo que Marx possui um diverso conjunto de obras que permite que você escolha diferentes recortes e desenvolva uma dada teoria que seja consistente com este recorte escolhido. Este é um dos motivos por que não considero que este seja uma boa meta a busca pelo ‘verdadeiro’ marxismo. Ou indo além, não considero sequer existir um verdadeiro marxismo. Diversas vertentes buscam desenvolver diferentes momentos, ênfases e obras de Marx. Mas nenhuma é o próprio Marx continuando sua obra, logo, outro critério deve ser adotado.

Defendo aqui a ideia de abandonar o personalismo e a própria figura do Marx para visar a formulação do marxismo que melhor descreva a realidade. Partindo de uma perspectiva científica, adoto a postura de que o marxismo deve ter mais compromisso com a adequada descrição do que pode ser observado, do que respeito com as ideias de qualquer indivíduo do passado.

Evidentemente essa correta descrição vai precisar de alguns pressupostos corretos e é aqui que reside a herança de Marx. De certa forma, tenho uma formulação bastante flexível de marxismo. Entendo então, que por marxista, quero dizer que defendo que a correta descrição da realidade é obtida quando adotamos a concepção de mundo defendida por Marx e Engels, popularmente conhecida como materialismo dialético.

\subsection{Materialismo}

Por materialismo dialético eu quero dizer, em poucas palavras, que esta é a concepção de mundo que encontramos exposta na obra “Anti-Dühring”.

\begin{quote}
… o ‘sistema’ do Senhor Dühring, de que este livro é uma crítica, estende-se a domínios teóricos muito vastos: tive de segui-lo por toda Parte e opor às suas concepções as minhas. Assim, a crítica negativa resultou positiva; a polêmica transformou-se em exposição mais ou menos coerente do método dialético e da ideologia comunista defendida por Marx e por mim, numa série de domínios bastante vastos …

…Uma observação de passagem: tendo sido criada por Marx, e em menor escala por mim, a concepção exposta neste livro, não conviria que eu a publicasse à revelia do meu amigo. Li-lhe o manuscrito inteiro antes da impressão; e o décimo capítulo da Parte segunda, consagrada à economia política (Sobre a história crítica) foi escrito por Marx. Infelizmente, eu o tive de resumir por motivos extrínsecos. Era, aliás, hábito nosso ajudarmo-nos mutuamente na especialização de cada um.”

(Friedrich Engels em Anti-Dühring).
\end{quote}

Mais do que um conjunto de leis rígidas ou um método, materialismo dialético é para mim uma concepção de mundo pautada em alguns pressupostos gerais sobre como é este mundo.

\begin{quote}
Uma concepção de mundo abrange a totalidade das opiniões e concepções de um indivíduo ou sociedade sobre o mundo, sobre nós mesmos como seres humanos e sobre a vida e a posição dos seres humanos no mundo.

(Curriculum Of The Basic Principles Of Marxism-leninism)
\end{quote}

É a partir desta obra que gosto de introduzir noções de materialismo dialético e dialética materialista.

\begin{quote}
\begin{itemize}
\item O \textbf{Materialismo Dialético} é uma compreensão científica da matéria, da consciência e da relação entre ambas. Ele é usado para entender o mundo através do estudo dessas relações.

\item A \textbf{Dialética Materialista} é uma ciência que estuda as leis gerais do movimento, da mudança e do desenvolvimento da natureza, da sociedade e do pensamento humano.
\end{itemize}
(Curriculum Of The Basic Principles Of Marxism-leninism)
\end{quote}

Ou seja, por materialismo dialético temos uma uma concepção materialista do mundo que considera uma relação dialética entre os entes que constitui esse mundo. Esta concepção de mundo, por materialista, quer dizer que vai se basear na ideia da existência de um mundo que existe objetivamente independente de qualquer consciência, portanto, anterior à mente, e consequentemente regido por leis naturais e impessoais.

Ao chamar a dialética materialista por ciência, não queremos que a palavra “ciência” aqui seja entendida no sentido moderno mais restrito, como uma disciplina acadêmica de universidade. Mas sim como um conjunto de conhecimentos sobre um dado assunto e/ou uma área de interesse de investigação. Assim, por dialética materialista temos um conjunto de pressupostos e a investigação destes pressupostos que nos orientaram sobre como as entidades que compõem o mundo se relacionam e se transformam. Ela diz respeito sobre nosso interesse em entender como o mundo evolui.

A dialética materialista vai pressupor um mundo em constante transformação e totalmente conectado, em outras palavras, não há ente no mundo nem isolado e nem que permaneça inalterada.

Materialismo dialético e dialética materialista formam então uma só teoria e concepção de mundo que deve também englobar tanto como concebemos o mundo tanto quanto sobre o nosso lugar nele. Não devem vê-los como duas coisas separadas e unidas mecanicamente: materialismo e dialética. Temos uma concepção de mundo materialista que engloba também uma perspectiva dialética de como este mundo se transforma.

Estas definições são bastante genéricas e flexíveis. Não são nada mais do que alguns pressupostos básicos sobre a natureza e transformação do mundo. Por marxismo, ainda faço uma adição de um posicionamento político bem definido em defesa das massas. Isto é algo que não surge necessariamente apenas de uma correta compreensão do mundo, porém é uma parte intrínseca do movimento político que se iniciou com Marx. De forma grosseira, entendo então o marxismo tanto quanto um movimento científico em busca da adequada descrição da realidade pautado em alguns pressupostos materialistas e dialéticos, quanto um movimento político de transformação dessa realidade em prol das massas, necessariamente.

Esta forma particular de conceber o marxismo não é livre de polêmicas. Tanto à inflexibilidade com que defendo Engels e sua obra “Anti-Dühring” na base do marxismo, quanto a flexibilidade com que defino o marxismo quanto às questões mais específicas são questões passíveis de polêmicas.

Marx e Engels trabalharam na construção de uma crítica direcionada principalmente à três grandes movimentos de sua época: a filosofia alemã, a economia política inglesa e o socialismo francês. Consequentemente da crítica negativa a estes três movimentos nasce uma exposição positiva do que poderíamos chamar de materialismo dialético, economia marxista e socialismo científico. Há muitos debates se existe uma economia marxista ou apenas uma crítica à economia política. Da mesma forma defendo que podemos levantar uma discussão se existe uma filosofia marxista ou apenas uma crítica à filosofia. Mas também vejo estas discussões como questões secundárias.

\begin{quote}
A tarefa da filosofia (que depois será também alvo de crítica de Marx e será substituída por teoria, entre outros termos ocasionais, como “socialismo científico”, “ciência” e, após ele, será chamado marxismo ou teoria)...

(Nildo Viana em Introdução à crítica da filosofia do direito de Hegel, o manifesto inaugural do materialismo histórico)
\end{quote}

Mais importante que isso, entendo que uma consequência fundamental desta forma de definir o marxismo é que quebra-se o vínculo com Hegel. Não há necessariamente uma negação de Hegel, mas certamente há um rebaixamento da importância de Hegel. Hegel passa a ser visto como uma influência filosófica em Marx, mas não necessariamente um componente crucial do marxismo. O marxismo assim exposto, torna-se independente de Hegel e sua dialética.

Gramsci, nos Cadernos do Cárcere, nos diz que é necessário entender que a filosofia da práxis nasceu de “aforismos critérios práticos por um mero acaso” e nos explica que Marx dedicou sua força intelectual de forma sistemática a outros problemas, principalmente de natureza econômica. Sua concepção de mundo aparece então apenas de forma implícita nestes aforismos.

\begin{quote}
Uma das causas do erro pelo qual se vai em busca de uma filosofia geral que esteja na base da filosofia da práxis e se nega implicitamente a esta uma originalidade de conteúdo e de método parece que reside no seguinte: que se faz confusão entre a cultura filosófica pessoal do fundador da filosofia da práxis, ou seja, entre as correntes filosóficas e os grandes filósofos pelos quais ele se interessou fortemente em sua juventude e cuja linguagem reproduz frequentemente… os elementos de spinozismo, de feuerbachismo, de hegelismo, de materialismo francês etc. não são de modo algum partes essenciais da filosofia da práxis

(Cadernos do Cárcere)
\end{quote}

Não tenho como negar que há uma influência hegeliana em Marx durante toda a sua vida, e que esta influência vai deixar marcas em todas suas obras. Mas argumento que conforme à teoria marxista se desenvolve, já em Marx e Engels, essa influência vai sendo engolida por práticas, que ao meu ver, são mais científicas que filosóficas. Vejo ecos deste fenômeno na obra de Engels.

\begin{quote}
Logo que descobrirmos – e afinal de contas ninguém mais do que Hegel nos ajudou a descobri-lo – que, assim colocada, a tarefa da filosofia se reduz a pretender que um filósofo isolado realize aquilo que somente a humanidade em seu conjunto poderá realizar, em seu desenvolvimento progressivo – assim que descobrirmos isso a filosofia, no sentido tradicional da palavra, chega a seu fim. Já não interessa a “verdade absoluta”, inatingível por este caminho e inacessível ao único indivíduo, e o que se procura são as verdades relativas, adquiridas através das ciências positivas e da generalização de seus resultados por meio do pensamento dialético.\textbf{ A filosofia, em seu conjunto, termina com Hegel}; por um lado, porque em seu sistema se resume, da maneira mais grandiosa, todo o desenvolvimento filosófico; por outro lado, porque este filósofo nos indica, ainda que inconscientemente, a saída desse labirinto dos sistemas para o conhecimento positivo e real do mundo…

...Esta interpretação põe fim à filosofia no campo da história, exatamente da mesma forma que a concepção dialética da natureza torna a filosofia da natureza tão desnecessária quanto impossível. Agora, já não se trata de tirar do cérebro as conexões entre as coisas, mas de descobri-las nos próprios fatos. Expulsa da natureza e da história, só resta à filosofia um único refúgio: o reino do pensamento puro, no que dele ainda está de pé: a doutrina das leis do próprio processo do pensamento, a lógica e a dialética.

(Ludwig Feuerbach e o Fim da Filosofia Clássica Alemã)
\end{quote}

Mas também na obra de Marx.

\begin{quote}
É preciso “deixar a filosofia de lado”, é preciso desembarcar dela e dedicar-se como um homem comum ao estudo da realidade, tarefa para a qual existe uma gigantesca quantidade de material literário … A relação entre filosofia e estudo do mundo real corresponde à relação entre onanismo e amor sexual.

(Karl Marx em A Ideologia Alemã)
\end{quote}

E eu não me encontro totalmente sozinho nesta posição polêmica. Além do corte epistemológico de Althusser e das declarações de Gramsci, podemos citar um marxista mais moderno: Paul Cockshott. Além de questionar o fato de uma fração dos marxistas verem como vital o estudo de Hegel para o marxismo devido a sua influência, mas não terem a mesma opinião acerca de economistas como David Ricardo e Adam Smith, que também foram fortes influências em Marx no campo econômico, de forma ainda mais agressiva, Cockshott argumenta que vê no estudo de Hegel um potencial maior de confundir do que de esclarecer.

\begin{quote}
Taimur, você certamente é um palestrante muito bom e claro, mas não acho que a palestra tenha me tornado mais simpático à ideia de estudantes contemporâneos estudarem Hegel. Longe de ser uma ajuda, creio que seria um obstáculo ao seu desenvolvimento intelectual…

...Temos tantas ferramentas para analisar o mundo, desenvolvidas nos últimos 200 anos, que retroceder à década de 1820 seria um terrível passo retrógrado.

(Paul Cockshott em A comment on a lecture about Hegel)
\end{quote}

A ideia central da crítica é relativamente simples: Hegel tem o mérito de ser um pensador avançado para o seu tempo e uma influência positiva em Marx. Contribuindo de forma inegável para o desenvolvimento do que hoje temos por marxismo. Mas seus métodos (e consequentemente o seu estudo) não constituem mais necessariamente a forma mais eficiente de entender a realidade. Pautando o objetivo do marxismo como sendo a melhor investigação e compreensão da realidade, resta ao marxista moderno buscar as melhores ferramentas disponíveis para entender o mundo, e esta ferramenta não é, necessariamente, a dialética hegeliana.

\subsection{Dialética}

Devido a posição central que a dialética ocupa na polêmica que envolve Marx, Engels e Hegel, eu acredito ser interessante uma sessão cobrindo este assunto, ainda que rapidamente.

É comum a ideia de que há um método em Marx no sentido hegeliano, onde dialética assume de forma mais literal o papel de um método de apreensão da realidade (neste caso, normalmente estritamente da realidade social). A partir disso, algumas pessoas então vão se dedicar a procurar a dialética hegeliana oculta em o “O Capital”. E de tudo isso ainda podemos iniciar uma nova discussão se esta dialética hegeliana é parte vital ou não do marxismo. Afinal “O Capital” e a obra marxiana pode conter a dialética na concepção hegeliana, e ainda sim pode ser que isto não constitua o núcleo do marxismo.

Eu reconheço que em “O Capital” há uma tentativa por Marx de usar a dialética mais no sentido de Hegel do que Engels, especialmente no primeiro capítulo, Marx é quem diz isso no prefácio. O que se pode questionar é a profundidade e a relevância deste esforço. Este é um assunto polêmico, pois em seus escritos o que geralmente encontramos tende a ser mais uma crítica à filosofia do que uma descrição do seu próprio sistema filosófico.

\begin{quote}
Confessei-me, portanto, abertamente discípulo daquele grande pensador e coqueteei mesmo aqui e ali no capítulo sobre a teoria do valor com o modo de expressão que lhe é peculiar.

(Karl Marx em O Capital)
\end{quote}

Acredito que os manuscritos do Grundrisse sejam o mais perto que chegamos de um trabalho com este tom. Porém deve-se lembrar que este mesmo projeto foi abandonado por Marx. O resultado disso é que torna-se necessário procurar ao longo das obras publicadas e alguns manuscritos abandonados o método dialético, considerando que ele existe e está oculto

Já Engels, diferente de Marx, escreveu de forma mais propositiva sobre filosofia, por exemplo no “Anti-Dühring”. Engels se afasta mais de Hegel que Marx se afastou. É mais fácil encontrar resquícios da dialética de Hegel no “O Capital” do que no “Anti-Dühring”. Mas novamente, pode ser interessante recordarmos que esta obra foi e revisada por Marx antes da publicação, foi decidido em comum acordo entre ambos que Engels a escreveria, Marx contribuiu escrevendo um capítulo, e a recomendou em cartas. Mas diante de toda esta situação surge uma tensão sobre a importância de Hegel e da fidelidade à dialética hegeliana no marxismo.

\begin{quote}
… e se o Sr. Most falhou em notar que há muito a ser aprendido com as exposições positivas de Engels [no “Anti-Dühring”], não apenas por trabalhadores comuns e até mesmo ex-trabalhadores como ele, que se supõem capazes de conhecer tudo e se pronunciar sobre tudo no menor tempo possível, mas até mesmo por pessoas cientificamente educadas, então eu só posso ter pena dele por sua falta de julgamento.

(Karl Marx em carta para Wilhelm Bracke de 11/04/1877)
\end{quote}

Minha posição é de que “O Capital” não é bom devido ao uso da dialética hegeliana por Marx, mas apesar dela. Assim, não entendo que a aplicação da dialética hegeliana seja a principal contribuição do Marx, e consequentemente não é isto que vai definir o marxismo. Tendo a pensar ser incorreto acreditar que Marx fez suas descobertas \textbf{por usar} a dialética hegeliana. Ela foi útil no amadurecimento de Marx, mas não acredito que as descobertas de Marx devem-se estritamente ao fato de tentar usar a dialética hegeliana.

Entendo que temos mais uma organização da exposição das idéias seguindo um modelo típico dialético, do que realmente descobertas sendo feitas devido a aplicação deste método de investigação dialético. Penso bastante sobre como o primeiro capítulo é o que Marx diz explicitamente que queria render homenagem ao Hegel, e este também é reescrito já na segunda edição pelo próprio Marx. Isto faz com que aqueles com um interesse mais voltado pelo aspecto filosófico hegeliano do marxismo critiquem a própria edição de Marx por verem nesta edição uma “popularização” da exposição dialética.

Dessa forma, me parece que a dialética hegeliana no “O Capital”, acaba ganhando uma sobrevida por uma questão estética de como os resultados foram apresentados, não sendo igualmente claro que tenha se provado ser o método último de investigação da realidade.

É popular encontrarmos nesta discussão a distinção entre método de exposição e método de investigação. Assim, o método de exposição apresentado no “O Capital” é de relativamente fácil observação pois é diretamente acessível pela própria exposição das ideias no livro. Mas o método de investigação é onde reside os problemas e os estudiosos sofrem para tentar “reconstruí-lo” investigando “método oculto” no livro.

Marx comenta a existência de um “método oculto” em “Para a Crítica da Economia Política”. Por isso eu entendo a existência desta discussão acerca da centralidade de Hegel em Marx. Minha crítica está na avaliação de que esta reconstrução e aplicação seja necessária.

\begin{quote}
Marx reconhece que a elaboração do método de apresentação de seus estudos foi facilitada pela releitura de algumas passagens da Lógica de Hegel na segunda [carta] Marx diz que sua publicação anterior – Para a crítica da economia política (1859) – continha argumentos que foram “popularizados”, algo que também acontecerá alguns anos depois na análise da teoria do valor , além de confidenciar que “o método está muito mais escondido”

(A “nova leitura de Marx”: um mapeamento de suas Premissas e desenvolvimentos)
\end{quote}

Entendo que a dialética hegeliana é incapaz de entregar tudo que promete a seus defensores, assim como acho provável que durante a própria investigação do capitalismo, Marx tenha visto ser necessário ir além da proposta de dialética hegeliana.

Precisamos então discutir “qual” é a dialética que eu considero importante. Eu defendo que seja o conjunto de pressupostos gerais como apresentado por Engels, ou seja, princípios de como o mundo se transforma e evolui. Na minha visão, esta é uma formulação mais interessante da dialética. Dessa forma, ela funciona mais como um guia para a investigação, como um conjunto de princípios norteadores, mais do que um método em si. Além de ser também menos nebuloso e mais flexível.

Eu não acredito que seja coincidência que até hoje os estudiosos debatam e adotem posições tão contrárias quanto a correta interpretação da dialética hegeliana. Um exemplo, é a acusação moderna de que interpretar a dialética como uma tríade, é uma interpretação superficial. Ainda que seja verdade, mas essa interpretação esteve presente desde os primeiros estudiosos de Hegel, não é apenas uma opinião de ignorantes que surgiu apenas agora. Debate-se até mesmo, se o próprio conjunto de obras de Hegel utiliza sua dialética.

\begin{quote}
Outras seções da filosofia de Hegel não se encaixam no exemplo triádico e clássico de Ser-Nada-Devir, como observaram até mesmo intérpretes que apoiaram a leitura tradicional da dialética de Hegel.

(Dialética Hegeliana na Enciclopédia de Filosofia de Stanford)
\end{quote}

Ainda que eu reconheça este resquício de Hegel nas obras marxianas, me parece claro que todos grandes movimentos políticos derivados da obra marxista, todo o movimento político marxista que obteve êxito em algum grau, todas as revoluções socialistas que se inspiraram nas obras de Marx, se basearam em uma teoria no qual não teve especial carinho pela dialética hegeliana.

O que acredito que levanta um debate saudável sobre quais são os aspectos da teoria marxista que através da prática se provaram válidos e fundamentais. Sei que há um esforço de alguns intelectuais em reconstruir algumas destas teorias sob a luz da dialética hegeliana, mas este é um esforço posterior tanto ao desenvolvimento da teoria quanto da prática revolucionária atrelada a esta teoria. Um contra-argumento possível é o de que quando estes movimentos falharam, falharam exatamente por essa deficiência teórica. Pessoalmente acho essa resposta simplista e preguiçosa. Eu tenho uma outra forma de abordar esta questão que irei expor mais adiante.

Entendo possíveis acusações que devo receber por “deturpar” o marxismo, pois é verdade que defendo uma identificação como marxista que de certa forma, é bem simples e flexível. 

\begin{quote}
Se você é capaz de tremer de indignação a cada vez que se comete uma injustiça no mundo, então somos companheiros.

(Che Guevara)
\end{quote}

Tenho por marxista, uma questão de princípios gerais que moldam nossa concepção do mundo e sobre como nos posicionamos neste mundo. Defender os pressupostos materialistas de como o mundo é (matéria antes da mente, etc) em conjunto com os pressupostos dialéticos de como esse mundo material evolui (tudo está em constante transformação, etc), associado com uma questão política sobre como nos posicionamos e atuamos politicamente no mundo (em defesa da classe dos trabalhadores e pela superação do capitalismo).

Assim, marxismo deixa de ser sobre dominar um método específico, não é sobre dialética hegeliana, mas passa a ser sobre uma concepção de mundo, sobre como vemos o mundo e como interagimos com este mundo. Acusam personalidades como Lenin, Che Guevara e Ho Chi Minh de não entenderem a dialética hegeliana, o que de certa forma eu não tenho estudo suficiente para discordar. Porém minha posição é de que talvez o fato de que “não entenderam”, ainda que verdadeira, não seja importante.

Talvez também, apenas não quiseram usar ela. Mas mais que isso, possivelmente não precisaram usar ela. Nem na sua prática revolucionária e nem na sua atividade teórica. Ao mesmo tempo também concordo que Marx teve a intenção de a usar, que isso era importante para Marx. Então a minha posição é de reconhecer a influência da dialética hegeliana em Marx sem estender o status de necessária ao marxismo. Conscientemente eu divirjo de Marx. Ainda que tenho meus motivos para pensar que não completamente.

Pois ao mesmo tempo que entendo a importância que a dialética hegeliana teve para Marx, ao analisar o conjunto das obras escritas em conjunto e suas correspondências trocadas com Engels, também me parece que não era um problema para Marx que Engels tivesse um posicionamento diferente sobre esta questão.

Ficamos então, de certa forma, em uma encruzilhada: ou então toda a teoria produzida por marxistas-leninistas é errada porque não usou devidamente a dialética, incluindo até mesmo partes das obras de Engels que foram produzidas enquanto era contemporâneo e foram revisada por Marx. Ou estes autores deram sorte de não usar dialética hegeliana corretamente e ainda sim produzir algo útil. O que me parece uma posição frágil

Também podemos considerar que estas obras poderiam se tornar ainda melhores caso se tornassem mais fiéis a Hegel. Uma posição intermediária: ainda que não totalmente erradas, estas teorias e práticas se beneficiariam se fossem formuladas de forma a serem mais corretas de uma perspectiva hegeliana. Porém, enquanto incorretos, estes conhecimentos obtiveram sucesso ao serem suportes teóricos para diversas revoluções, ao mesmo tempo de que não tenho registro de nenhuma revolução que triunfou ao fazer uso de alguma teoria revisada para se tornar mais fiel à Hegel. O que coloca em dúvida a veracidade dessa posição.

Por fim, talvez a ideia mais simples seja mesmo a de que a dialética hegeliana não seja necessária. Esta é a posição que eu adoto.

\subsection{Método}

Quando eu ouvi a primeira vez eu concordei com a afirmação de que o que definia o marxismo era o método. Na época eu pensei que por método queríamos falar sobre o materialismo dialético e isto era uma questão simples.

Mas hoje acho que essa discussão demanda um pouco mais de cuidado. Pois eu noto que muitas vezes quando se usa esta afirmação, ela diz respeito estritamente sobre a dialética hegeliana ser o método marxista. Porém, como tenho discutido, a diferença da dialética hegeliana para o que entendo por dialética é mais profunda do que eu tinha pensado inicialmente, inclusive, sendo justo questionar se aquilo que chamo de dialética ainda consiste de fato de um método.

Marxistas-leninistas - tradição no qual eu costumo me inserir - também chamam a dialética como método, mesmo que o que entendem por dialética seja reconhecidamente diferente do que é entendido por dialética pelos hegelianos. Eu tenho um certo incômodo com esta prática. Pois reflito se para que isso fosse possível foi necessário realizarmos uma espécie de ressignificação da palavra. Suponho que isto tenha acontecido pois queríamos herdar o prestígio de que a dialética hegeliana gozava para o que nós agora estávamos chamando de dialética, ainda que sejam coisas diferentes.

Acredito que um interesse semelhante nos fez chamar o materialismo dialético de “filosofia” marxista. Para manter o prestígio de dizer que o marxismo é uma filosofia, afinal, filosofia sempre gozou de prestígio na sociedade. Para que isso fosse possível de maneira análoga, foi necessário ressignificar o que é filosofia. Assim filosofia, no marxismo foi definido exatamente da mesma forma que Engels definiu a teoria que desenvolveu com Marx no prefácio do “Anti-Dühring”, como uma concepção de mundo e um método de investigação.

\begin{quote}
Filosofia é então a unidade de uma concepção de mundo e metodologia.

(Dialectical Materialism)
\end{quote}

Partindo para exemplos mais pontuais, quando eu vou estudar física, eu entendo que não vou aplicar um “método” dialético da mesma forma que se pensaria caso este método fosse a dialética hegeliana. A influência da dialética engelsiana na minha pesquisa vai se dar no fato de que o método de investigação adotado precisará ser coerente com um conjunto de princípios que eu entendo por dialética. 

Podemos citar como exemplo o Big Bang. Por eu ter esta concepção marxista do mundo, eu sou incentivado a investigar teorias que busquem explicar a origem do Big Bang considerando a existência de um universo anterior ao Big Bang, pois por princípio, me vejo querendo rejeitar a ideia do Big Bang como o começo de tudo. Isto decorre do fato de que o materialismo dialético propõe um universo que esteja infinitamente em transformação, tanto no passado quanto no futuro, o que entendo que contraria uma ideia de começo ou de fim.

Também é por este conjunto de princípios, que por exemplo, parece errado imaginar que as espécies sempre foram da forma que as vemos hoje, ou que o sistema solar sempre existiu da forma que observamos hoje. Essa imutabilidade das coisas viola os pressupostos da dialética que compõem minha concepção do mundo.

A minha concepção de dialética soa então significativamente diferente da adotada pelos hegelianos, e eu acredito ter menos motivo para chamar isto de método. A forma que entendo em Engels, dialética diz mais respeito a um conjunto de pressupostos de como o mundo evoluiu. Defendo que esta é a posição mais interessante de adotar para buscar interpretar e modificar o mundo, de forma que ainda que o Marx tenha de fato um método oculto em sua obra, a reconstrução deste método se torna algo não apenas secundário, mas também dispensável.

Além disso, deslocando a definição de marxismo de método para uma concepção de mundo, reforço novamente que devemos incluir a defesa da classe proletária nesta concepção de mundo, uma vez que a concepção de mundo inclui não apenas como a gente vê o mundo, mas também como nos posicionamos diante dele. Uma postura que não surge naturalmente no marxismo se o entendermos apenas por um método.

Para finalizar esta parte filosófica, na prática, a forma com que concebo o materialismo dialético é então como uma concepção de mundo. Que vai incluir tanto pressupostos sobre a relação entre a matéria e a consciência, quanto sobre como o mundo se transforma. Frisando sempre que a distinção entre materialismo e dialética, como se materialismo dialético fosse uma combinação de duas coisas distintas, é errada. Talvez precisaríamos de um novo nome, ou então adotar a sugestão de Gramsci e chamar de filosofia da práxis.

O que eu posso conceber é retomar a distinção entre materialismo dialético (como concepção de mundo) e dialética materialista (como ciência). Mas aqui há uma sutileza: as ‘leis da dialética’ conforme enunciadas por Engels, também estão incorporadas no materialismo dialético. Este conjunto de princípios sobre como o mundo se transforma e se conecta é um pilar fundamental da concepção de mundo do materialismo dialético. 

Portanto, por dialética materialista, não estamos nos referindo a uma ‘parte dialética’ que se combina com a ‘parte materialista’ para formar o materialismo dialético. Mas estamos nos referindo à ciência, no sentido mais geral usado por Engels, como a investigação destas leis.

\begin{quote}
A palavra “ciência”, e, por extensão, “científico”, no marxismo-leninismo, possui um significado específico … cabe observar que a expressão inglesa “scientific socialism” deriva do uso que Engels faz da expressão alemã “wissenschaftlich sozialismus”.

“Wissenschaft” é uma palavra que pode ser traduzida diretamente como “arte do conhecimento” em alemão, e esse termo abrange um conceito muito mais amplo e geral do que a palavra “ciência” como é geralmente usada em inglês … Assim, o “socialismo científico” pode ser entendido como um conjunto de teorias que analisa e interpreta o mundo natural para desenvolver um corpo de conhecimento, o qual deve ser constantemente testado contra a realidade, com o objetivo de transformar o mundo para alcançar o socialismo por meio da liderança do proletariado.

(Curriculum Of The Basic Principles Of Marxism-leninism)
\end{quote}

Acredito que seja válido comentar que ainda que eu não conceba o materialismo dialético como um método, mas sim como uma concepção de mundo, isso não implica que marxistas não podem trabalhar no desenvolvimento de um método de investigação da realidade derivado dessa concepção comunista do mundo. 

O que estou afirmando é apenas que eu não coloco nenhum método como sendo um pilar do marxismo, assim como que este método não é o materialismo dialético em si, mas seria um produto do materialismo dialético e da dialética materialista. De toda forma, o que nos interessa é que a dialética engelsiana não é nem a mesma de Hegel, e possivelmente nem um método propriamente dito. E isto não deveria ser um problema.


\subsection{Filosofia}

Eu gostaria ainda de completar a polêmica iniciada em textos anteriores.Agora estamos prontos para dispensar a filosofia. Ou pelo menos, dispensar a leitura mais tradicional do papel da filosofia.

Praticamente desde que me interessei por marxismo eu tenho lido sobre a filosofia. Isso decorreu do fato de que eu sempre fui crítico e cético quanto à filosofia. Porém eu gosto de entender o que quero criticar, logo, senti a necessidade de entender melhor a filosofia. Claro que não significa que me tornei um especialista no assunto, mas aprendi um pouco no percurso.

Uma das primeiras coisas que tive a oportunidade de aprender foi de que não há um consenso sobre o que é filosofia\cite{glauber1}. Cada pensador que ganhou o direito de ser reconhecido como filósofo definiu a filosofia de uma forma diferente. E isso de certa forma fornece uma certa flexibilidade para que a filosofia seja defendida de forma abstrata. Em um momento filosofia é 'pensar de forma crítica' (e ninguém vai se opor à esta tarefa), mas em outro momento não existe filosofia chinesa pois a filosofia é definida de forma que necessariamente nasceu na Grécia antiga. 

Dificilmente a mesma pessoa que faz as duas alegações Pois implicaria em dizer que os chineses não eram capazes de pensar de forma crítica. Eu sou leigo, mas me parece que em grande parte a definição mais regular de filosofia na academia é aquela que é dada pela disciplina de história da filosofia colocando sua origem necessariamente no mundo grego antigo. É isto que vamos encontrar na coleção de “História Da Filosofia” de Giovanni Reale que é constantemente recomendada.

Este projeto de história da filosofia que trata ela como um ‘produto exclusivo do gênio grego’, me parece um projeto que nasceu intrinsecamente como um projeto eurocêntrico, desde Hegel que é essencialmente o fundador da história da filosofia\cite{usphegel}. E normalmente, quando eu teço críticas a ‘filosofia’, minha crítica se refere principalmente a esta tradição com origem em platão. Tradição esta que exige que todo novo filósofo para que mereça este título, seja capaz de ser inserido nesta grande árvore genealógica com as raízes em Platão, isto é, que de alguma forma, este filósofo dialogue direta ou indiretamente com platão.

Vale aqui dizer, que não sou o primeiro a abordar esta questão, pensadores da escola chamada ``teoria crítica'' distinguiram a ``teoria tradicional'' da ``teoria crítica'' para demarcar uma mudança no modo de produção dentro da filosofia, correspondendo de certa forma a transição efetuada por Marx na forma de conceber a produção de conhecimento\cite{glauber2}.  Lembrando também que o próprio Marx eventualmente para de se referir ao que esta produzindo como ``filosofia'' e passa a adotar outros termos, como ``teoria''\cite{nildo}

Outra descoberta que tive foi de que não há também consenso sobre o que é a filosofia marxista conforme exposto anteriormente. Porém, para que o materialismo dialético na formulação marxista-leninista recebesse o título de filosofia, foi necessário adotar uma definição de filosofia que fosse adequada. Pode ser lido no “Dialectical Materialism” do Spirkin, por exemplo, que filosofia é a unidade de uma concepção de mundo e um método de investigação, parafraseando o próprio Engels quando descreve a teoria que desenvolveu junto de Marx como uma concepção de mundo comunista e método de investigação dialético; 

O fato é que o materialismo dialético na concepção marxista leninista, é significativamente diferente da filosofia como é tradicionalmente concebida. Algo que o próprio Spirkin reconhece e o leva a classificar o marxismo como uma filosofia científica, e mais precisamente, como a primeira filosofia verdadeiramente científica. 

Um título que talvez não faça muito sentido da forma que a filosofia é tradicionalmente concebida. Essa ruptura com a filosofia clássica fica evidente por exemplo em Engels quando escreve o “Ludwig Feuerbach e o fim da filosofia clássica alemã” e o “Anti-Dühring”. Não por coincidência, meus textos preferidos sobre o assunto, e acredito que não por coincidência também, os textos recomendados por Lenin no seu famoso texto ``As Três Fontes e as Três partes Constitutivas do Marxismo''. Aqui por honestidade e para destacar a influência do Engels na formação do Marxismo, podemos observar que os dois principais textos recomendados por Lenin para abordar a filosofia marxista, são escritos por Engels, e não por Marx.

Mas como o próprio Engels coloca, isso não implica que necessariamente tudo da filosofia clássica tenha se tornado obsoleto e que a própria tarefa de filosofar não possua mais nenhuma utilidade. Eu defendo que tem a ver mais como uma redução de escopo e ambição. É nesse sentido que eu li e apreciei o “Tractatus Logico-Philosophicus” de Wittgenstein e principalmente “Investigações Filosóficas” do mesmo autor. 

Entendo hoje, que a filosofia, em sua formulação mais tradicional e abrangente, isto é, a filosofia que não é marxista, mas também não se restringe à tradição iniciada em Platão, mas sim uma reflexão especulativa geral sobre tudo, não deve ser uma teoria que visa explicar o mundo. Mas sim que a filosofia deve ser encarada como uma atividade de esclarecimento sobre o uso da linguagem: ela dissolve pseudoproblemas ao revelar quando estamos usando a linguagem além dos limites do sentido.

Me parece que Wittgenstein, de certa forma, assim como Engels, não elimina a filosofia por completo, mas altera a ideia sobre qual é o papel da filosofia. Ao invés de vermos na filosofia o papel de questionar como o mundo é, transferimos esta responsabilidade para as ciências e permitimos que a filosofia tenha por atividade questionar como usamos a linguagem.

A filosofia se torna então uma prática de esclarecimento da linguagem. Eu tendo a pensar que a filosofia na prática, que eu não iria discordar da utilidade, é exatamente isso. Quando se pergunta "Deus existe?" A gente acaba tendo que discutir o que queremos dizer com "Deus". E talvez até mesmo o que quer dizer com "existir". Se tivermos tudo esclarecido, podemos então buscar a comprovação ou não de sua existência através da ciência. 

Esta é a discussão que entendo ser por natureza filosófica. Lembrando novamente que tendo a ver o marxismo não como filosofia, mas como uma teoria. Dessa forma também cabe então a ciência produzir conhecimento, e não a filosofia. A filosofia, pela forma que opera, pelos limites da lógica, não é capaz de produzir conhecimento, pois ela não se baseia na observação e extração de 'verdade' nenhuma do mundo real. 

Acredito que minha posição é coerente com as ideias expostas por Engels:

\begin{quote}
Por toda a parte, não se trata mais de congeminar conexões na cabeça, mas de as descobrir nos factos. Para a filosofia desalojada da Natureza e da história, fica ainda então apenas o reino do pensamento puro, na medida em que ainda resta: a doutrina das leis do próprio processo do pensar, a lógica e dialéctica.

(Engels em Ludwig Feuerbach e o fim da filosofia clássica alemã)
\end{quote}

\subsection{Ciência}

Agora que já ataquei a filosofia falando que para mim materialismo não é filosofia, e ela encontrou seu auge e fim com Hegel e que já reduzi a filosofia clássica à função de resolver os problemas de linguagem, eu queria discutir a relação entre ciência e filosofia,e de certa forma, oferecer uma posição de negar que a filosofia seja a mãe da ciência. De forma simples, já vou resumir meu argumento. A ciência atual e a filosofia atual ambas se originam em uma mesma tradição mais antiga, que pode ser chamada tanto como “filosofia universal” quanto como “ciência universal” e não é nem a filosofia atual e nem a ciência atual.

E não apenas isso, a ciência moderna deve mais ao conhecimento prático e empírico acumulado por séculos pelas classes baixas, do que às abstrações frequentemente especulativas e erradas dos filósofos que pertenciam à classe dominante. Esse apagamento histórico de um lado e exaltação além da realidade do outro é um fruto direto da luta de classes.

\begin{quote}
Basta dizer que, antes de 1600, antes do "nascimento da ciência e da filosofia modernas", a colheita de frutos da árvore da filosofia era bastante escassa…

...Por outro lado, os antigos gregos e romanos e os antigos egípcios demonstraram uma arte e uma habilidade maravilhosas na construção e mesmo na engenharia mecânica. As suas abóbadas e arcos, que suportavam cargas tão pesadas, foram construídos com tal precisão e força que indicam um conhecimento altamente sofisticado da mecânica prática. Seja como for, o conhecimento destes antigos construtores e engenheiros não era científico, mas tecnológico. Era puramente empírico, não se baseava em princípios filosóficos.

(The Origin of the separation between science and philosophy)
\end{quote}

Começamos então com estes dados históricos. Platão, que está na raiz da filosofia ocidental moderna, teria nascido uns 4 séculos antes de cristo, o que nos dá uma janela de tempo de 2 mil anos de existência da filosofia sem que desse muitos frutos. Se formos voltar a Tales de Mileto, precisamos adicionar mais 2 séculos de improdutividade da filosofia.

Só isto já deveria ser o suficiente para questionar.  E este período de escassa produção corresponde exatamente ao período anterior ao nascimento da ciência moderna. Isto é, a filosofia só foi capaz de produzir frutos quando a ciência moderna surgiu. Posição semelhante à de Einstein:

\begin{quote}
Admiramos a Grécia antiga porque fez nascer a ciência ocidental… Mas para atingir uma ciência que descreva a realidade, ainda faltava uma segunda base fundamental que, até Kepler e Galileu, foi ignorada por todos os filósofos. Porque o pensamento lógico, por si mesmo, não pode oferecer nenhum conhecimento tirado do mundo da experiência. Ora, todo o conhecimento da realidade vem da experiência e a ela se refere. Por este fato, conhecimentos, deduzidos por via puramente lógica, seriam diante da realidade estritamente vazios. Desse modo Galileu, graças ao conhecimento empírico, e sobretudo por ter se batido violentamente para impô-lo, tornou-se o pai da física moderna e provavelmente de todas as ciências da natureza em geral. Se, portanto, a experiência inaugura, descreve e propõe uma síntese da realidade…

(Einstein em Como vejo o mundo)
\end{quote}

Acredito que esta analogia sobre a colheita da árvore da filosofia venha da afirmação de Descartes de que a ‘utilidade’ da filosofia está nos seus frutos. Evidentemente isso não significa que as civilizações do passado não detinham conhecimento, pelo contrário.

Porém, se pegarmos o conhecimento do mundo antigo, que era um conhecimento menos geral e mais técnico voltado ao desenvolvimento de tecnologias, isto existiu de forma independente da filosofia. Frequentemente se fala isso sobre outras nações como Índia e China, mas muitos tentam colocar civilizações ocidentais como ‘superiores’ pois supostamente não tinham ‘apenas tecnologia, tinham filosofia’. Porém, como podemos ver desde o surgimento da filosofia, até o nascimento da ciência, ela não produziu grandes frutos, nem mesmo no mundo ocidental.

\begin{quote}
O grande fosso que existiu durante a Antiguidade e a Idade Média entre "princípios" e "saber-fazer", entre "ciência" e "tecnologia", estava intimamente ligado à estrutura social. O "saber-fazer" do artesão e do engenheiro era adquirido através do trabalho físico e não do puro esforço intelectual. Era o conhecimento empírico do homem de baixo estatuto social, do artesão e até do escravo. O conhecimento filosófico era cultivado e promovido pelos homens de estatuto social elevado, nomeadamente os padres e os estadistas. Os estratos "inferiores" coleccionavam "factos"; os superiores avançavam princípios. O contacto entre os dois tipos de conhecimento era desencorajado pelos costumes sociais. Se um homem de estatuto social elevado tentasse aplicar a sua "filosofia" ou "ciência" ao trabalho do artesão, era condenado ao ostracismo.

(The Origin of the separation between science and philosophy)
\end{quote}

Temos então uma separação social, uma questão de classes na valorização do conhecimento produzido pelos trabalhadores de classes baixas e pelos intelectuais da classe alta. Ou seja, não é apenas que o conhecimento empírico não se baseia na filosofia, mas o conhecimento empírico estava fatalmente separado do conhecimento filosófico. Não apenas eram áreas distintas, como as pessoas que praticavam eram pessoas distintas. 

E ainda mais, eram pessoas que não tinham contato entre si pois pertenciam a diferentes classes. Se alguém de classe alta fosse ‘pego’ trabalhando com um conhecimento empírico, seria punido socialmente. Então aqui a gente tem ainda que a evolução destas diferentes áreas do conhecimento ocorreu de forma independente, por pessoas diferentes, de classes sociais diferentes. 

Evidentemente eu não estou negando que em torno do século XV, quando surgiu o que hoje chamamos de ciência, esteve relacionado com o que hoje chamamos de filosofia. Devido a mudanças sociais na época, membros da classe alta com mais tempo livre para abstrair questões mais práticas, treinados na filosofia, tiveram contato e se interessaram então por este conhecimento técnico, criando as condições necessárias para o surgimento de em uma ciência mais geral e abstrata. 

Se quisermos, podemos localizar aqui então a contribuição da filosofia. Porém, ainda vejo esta hipótese com certo ceticismo. Neste processo, o que eu percebo é que com a transformação do conhecimento técnico em ciência, a separação que existia entre filosofia e conhecimento técnico, se transforma na separação entre filosofia e ciência. E é preciso notar que descrevendo a história da ciência desta forma, temos um conhecimento empírico que se transforma em ciência, beneficiada do contato com a filosofia, mas está mesmo permanece tanto separada do conhecimento técnico no passado, quanto da filosofia no presente.

Ainda sim, a história do desenvolvimento da ciência moderna está atrelada também às necessidades do modo de produção moderno. Se reconhece que o conhecimento técnico nasceu em resposta às necessidades materiais, por exemplo, a matemática como necessidade de calcular áreas e controlar o imposto, mas a ciência moderna também responde a condições materiais específicas. O desenvolvimento da termodinâmica não é apenas causa da revolução industrial, mas também consequência.

O desenvolvimento da ciência também responde a uma questão social e de necessidade. Ainda que uma vez praticada predominantemente por membros da classe dominante que tinham menos necessidade material, a ciência pode perseguir uma abstração maior e menos imediata que o conhecimento técnico praticado predominantemente pelas classes baixas. Porém, da mesma forma que a astronomia surge por aqueles que se interessavam em astrologia, ou a química por quem se interessava por alquimia, e nunca vamos ouvir alguém defender a astrologia como algo sério por ser a mãe da astronomia, questiono o rótulo da filosofia como ‘mãe da ciência’. Na verdade, diria que as ciências se desenvolveram apesar destes sistemas de pensamento.

Dessa filosofia pré-1600, considerando então uma espécie de “ciência universal”, chamada de filosofia, o autor diz que surge então a “filosofia própria” e a “ciência própria”. Se trocarmos então o título de “filosofia” por “ciência universal”, poderíamos dizer que a ciência é a mãe da filosofia? Acredito que não seria uma ideia bem recebida.

\begin{quote}
Kant afirmou sem rodeios que os factos observáveis do mundo físico são completa e satisfatoriamente descritos e sistematizados pela "ciência propriamente dita"; a "filosofia propriamente dita" nunca nos poderá dizer nada sobre eles. A função da filosofia é erigir uma superestrutura sobre a ciência, não para fazer afirmações sobre factos físicos, mas antes para atribuir "valor" a certos tipos de ação humana... A investigação que procura afirmações verdadeiras para além do domínio da "ciência propriamente dita" e tenta obter essas afirmações a partir de uma "interpretação filosófica da ciência" é chamada "investigação metafísica". O resultado deste tipo de investigação é um ramo do pensamento sistemático conhecido como "metafísica", cuja função é fornecer-nos regras para a conduta humana, para nos mostrar "o bem".

(The Origin of the separation between science and philosophy)
\end{quote}

Ou seja, a filosofia conforme descrita aqui trata daquilo que não são os fatos do mundo físico. Kant então ‘proíbe’ a filosofia de fazer afirmações sobre o mundo físico. Se nos lembrarmos da árvore do Descartes, a metafísica era a raiz da árvore. Ainda que o termo metafísica não tenha um significado constante ao longo da história, tendo a pensar que aquilo que tradicionalmente é concebido como filosofia atualmente é em grande o que era então classificado como metafísica, a raiz da árvore.

E do tronco em diante, temos a física, ou de forma mais geral, as ciências. Antes de concluir, acredito que é interessante acompanharmos mais um trecho:

\begin{quote}
Segundo Duhem, em suma, Platão distingue três graus de astronomia: observacional, geométrica e teológica. A conceção platónica da astronomia torna clara a natureza da divisão dentro da "filosofia" ou da "ciência". A astronomia observacional e geométrica constituem, em conjunto, um instrumento que permite prever as posições observáveis dos astros, mas a "astronomia teológica" não é necessária nem sequer útil para esse efeito. O que ela pode fazer é dar-nos uma imagem do universo que será considerada uma réplica da sociedade humana ideal; ela pode ajudar a moldar a conduta humana.

(The Origin of the separation between science and philosophy)
\end{quote}

E se ainda há dúvidas sobre o caráter teológico e a distinção explícita entre ciência e filosofia desde Platão, podemos consultar as próprias palavras de Platão:

\begin{quote}
Existem três graus de conhecimento. O grau mais baixo é o conhecimento pela observação sensorial. . . . O grau supremo é o conhecimento pelo intelecto puro... Entre o primeiro (mais baixo) e o terceiro (supremo) grau de conhecimento existe um tipo de raciocínio misto e híbrido que ocupa o grau intermediário (segundo). O conhecimento nascido desse raciocínio intermediário é o conhecimento geométrico. A esses três graus de conhecimento correspondem três graus de astronomia.

A percepção sensorial é responsável pela astronomia da observação. Esse tipo de astronomia persegue as curvas complicadas descritas pelas estrelas... e não pode fornecer nenhuma relação comensurável ao aritmético, nem nenhuma figura definida ao geômetra... Por meio do raciocínio geométrico, a mente produz uma astronomia capaz de figuras precisas e relações constantes. Essa “verdadeira astronomia” substitui os caminhos erráticos que a astronomia observacional atribuía às estrelas por órbitas simples e constantes; ela produz as aparências complicadas e variáveis que são conhecimento falso. ... É a aproximação do conhecimento pelo intelecto puro, que revela a terceira e suprema astronomia, a astronomia teológica. ... Na constância dos movimentos celestes, ela vê uma prova da existência de espíritos divinos que estão unidos aos corpos celestes.

(The Origin of the separation between science and philosophy)
\end{quote}

Podemos perceber que os dois primeiros graus de conhecimento, os dois primeiros graus de astronomia são exatamente o que corresponde a ciência. O primeiro degrau corresponde ao conhecimento empírico, à experimentação. À coleta de dados e observação da natureza através dos sentidos. O segundo degrau corresponde à generalização e abstração dos dados observados, à construção de modelos gerais, na física em particular, à construção das leis matemáticas que generalizam os casos observados.

E deixamos então para a filosofia o último grau. O conhecimento superior na ideia dos filósofos. Mas novamente o que a filosofia aqui está discutindo, já não é mais o mundo físico, mas discutindo mais a questão do ideal humano, algo que visa moldar a conduta humana. Neste sentido, eu vejo uma coerência na definição do que é filosofia, de Platão a Kant. 

Neste debate, uma definição interessante entre ciência e conhecimento técnico, é que o conhecimento técnico nos permite lidar com circunstâncias que nos deparamos antes, e o objetivo do pensamento científico é aplicar o passado a novas circunstâncias. Isto é, uma generalização do conhecimento produzido no passado para que consigamos ter uma capacidade de predição  sobre novas coisas do futuro. Mas novamente, isto conecta o desenvolvimento da ciência como a evolução natural do conhecimento técnico do passado.

Evidentemente a história da humanidade é complicada e certamente a ciência não tem apenas uma origem, mas meu esforço aqui é resgatar esta conexão esquecida.  O artigo também cita um possível caminho do conhecimento, que parte da tecnologia, passa pela ciência e culmina na filosofia, considerando-a uma unificação no conhecimento. É uma possibilidade, mas isto torna a filosofia uma espécie de ciência universal, e em grande parte, retoma nossas definições anteriores da filosofia como uma concepção de mundo. Neste caso, uma concepção de mundo moldada pelo conhecimento científico.

Segundo Philipp, a primeira grande cisão entre filosofia e ciência moderna ocorre com a astronomia. O filósofo e astrônomo, até então a mesma pessoa, teria se visto obrigado a escolher uma das duas opções:

\begin{itemize}
    \item Manter os princípios “inteligíveis” e conformar-se à falta de concordância com os dados observados.
    \item Desistir da compreensão do mundo apenas pela razão e ficar satisfeito com uma hipótese “ininteligível” que concorda com os fatos.
\end{itemize}
É evidente que a filosofia opta pelo primeiro caminho, priorizando os princípios, e a ciência pelo segundo buscando concordar com os fatos.  Neste contexto, “inteligível” não significa simplesmente compreensível no sentido comum, mas aquilo que pode ser deduzido da razão pura como necessário e evidente, obedecendo a princípios racionais considerados universais e autoexplicativos. Um mundo inteligível é aquele que “faz sentido por si”, independentemente da experiência, como no ideal clássico da filosofia antiga. 

Já o “ininteligível” não designa algo absurdo ou irracional, mas uma hipótese que não pode ser justificada por princípios racionais evidentes, sendo aceita apenas por sua concordância com os fatos observados.  Como exemplo, um princípio inteligível no sentido clássico é a ideia de que as órbitas dos planetas deveriam ser círculos perfeitos, já que o círculo era concebido como a forma geométrica perfeita e, portanto, a mais adequada para descrever o movimento dos corpos celestes. Em contraste, a hipótese ininteligível e apoiada pelos dados, é a de que as órbitas planetárias possuem excentricidades, isto é, são elípticas. Essa hipótese concorda com os dados observacionais, mas não decorre de nenhum princípio racional autoevidente sobre perfeição, harmonia ou necessidade.

A partir disso, Philipp destaca dois critérios para checar a validade de uma declaração:

\begin{itemize}
    \item Pode ser derivado logicamente de uma declaração auto-evidente;
    \item Pode ser tratado como uma hipótese que podemos derivar consequências e testar com fatos.
\end{itemize}
Evidentemente a primeira é um critério “filosófico” e o segundo “científico”. 

\subsection{Socialismo}

Já comentamos sobre a filosofia marxista (\textit{vulgo} materialismo dialético), um próximo tema importante é o socialismo marxista (\textit{vulgo} socialismo científico). Socialismo científico pode ser entendido como nada mais sendo do que a investigação e a prática acerca da superação do capitalismo. Consequentemente, a construção de uma sociedade pós-capitalista (\textit{vulgo} comunista).

Colocado dessa forma, o socialismo, como um dos componentes do marxismo não deve ser entendido de forma simplista como um mero modelo de estado, nem como sendo apenas uma etapa de transição. De fato, Marx nunca falou sobre socialismo neste sentido, o que há é uma distinção do comunismo em fase inferior e fase superior, assim como menções à ditadura do proletariado, mas nenhuma destas etapas recebe na obra marxiana o nome do socialismo.

\begin{quote}
Entre a sociedade capitalista e a sociedade comunista medeia o período da transformação revolucionária da primeira na segunda. A este período corresponde também um período político de transição, cujo Estado não pode ser outro senão a ditadura revolucionária do proletariado…

...Do que se trata aqui [- a fase inferior do comunismo -] não é de uma sociedade comunista que se desenvolveu sobre sua própria base, mas de uma que acaba de sair precisamente da sociedade capitalista e que, portanto, apresenta ainda em todos os seus aspectos, no econômico, no moral e no intelectual\~ o selo da velha sociedade de cujas entranhas procede…Estes defeitos, porém, são inevitáveis na primeira fase da sociedade comunista…

...Na fase superior da sociedade comunista, quando houver desaparecido a subordinação escravizadora dos indivíduos à divisão do trabalho e, com ela, o contraste entre o trabalho intelectual e o trabalho manual; … só então será possível ultrapassar-se totalmente o estreito horizonte do direito burguês e a sociedade poderá inscrever em suas bandeiras: De cada qual, segundo sua capacidade; a cada qual, segundo suas necessidades.

(Marx em Crítica ao programa de Gotha)
\end{quote}
Já em Engels, quando se fala de socialismo científico em “Do Socialismo Utópico ao Socialismo Científico”, refere-se ao socialismo como uma ciência, no sentido amplo como citado anteriormente. Como um conjunto de conhecimentos e práticas relacionadas à investigação acerca da superação do capitalismo. Nunca como um modelo de sociedade.

É neste ponto que se inserem algumas das maiores contribuições de Lênin. Sua obra aprofunda a teoria marxista ao formular uma concepção de partido revolucionário que atua como instrumento consciente da transição. Também gostaria neste assunto de destacar a contribuição do Paul Cockshott em “\textit{Towards a New Socialism}” onde ousa propor um possível modelo de sociedade comunista.

\begin{quote}
Em termos de instituições democráticas e mecanismos de planejamento eficientes, devemos reconhecer que os problemas que emergiram no caso soviético refletem certas fragilidades do marxismo clássico… Marx, Engels e Lenin foram muito mais incisivos em suas críticas ao capitalismo do que em suas teorizações positivas sobre a sociedade socialista…

...Quanto aos mecanismos de planejamento, Marx e Engels apresentaram algumas sugestões interessantes, mas estas nunca foram desenvolvidas além do nível de generalidades bastante vagas.

(Towards a New Socialism)
\end{quote}

Esta proposta recebe o nome de cibercomunismo por alguns autores, e nada mais é do que a proposta de um modelo de comunismo onde utiliza-se os mais recentes avanços científicos e computacionais no desenvolvimento e gestão desta sociedade pós-capitalista. Não ficando restrita ao uso da computação apenas para cálculos econômicos, mas transformando a tecnologia em uma ferramenta de efetiva participação popular nas decisões políticas e econômicas desta sociedade, isto é, colocando nas mãos do trabalhador o poder de decisão sobre a produção e distribuição de bens.

\begin{quote}
A primeira coisa a esclarecer sobre a natureza deste projeto é que o planejamento comunista nada mais significa do que a regulação consciente, racional e democrática da economia…

...Nesse sentido, o cibercomunismo está longe de ser uma “economia algorítmica” ou um “governo por algoritmos”. Pelo contrário, trata-se de um tipo de economia onde … os processos de avaliação e decisão individual, bem como os de deliberação democrática, permeiam toda a operação econômica: desde a escolha pessoal da profissão ou dos meios de consumo até a definição democrática dos objetivos de desenvolvimento econômico e social, passando pelo controle social do investimento, pela promoção de talentos individuais e pela implementação de novos projetos produtivos. Não se trata, portanto, de um “supercomputador” ou de uma “inteligência artificial central” que tente “adivinhar os desejos e preferências dos indivíduos”, “prever o futuro” ou “programar a vida social”. Muito pelo contrário: o que se propõe aqui é justamente ampliar o horizonte da autonomia individual e da participação democrática com o auxílio dos avanços científico-técnicos, dentro de uma estrutura de igualdade social e plenas liberdades individuais.

(Economic Calculation, Complexity, and Cyber-communism: Bad News for the Austrian School)
\end{quote}

\subsection{Eleições}

Sobre a superação do capitalismo e a construção do capitalismo, um tópico que não pode ser ignorado é a revolução. O que é a revolução comunista e em que condições ela se dá. Aqui eu quero esboçar apenas um rascunho sobre o que penso acerca da situação que nos encontramos no Brasil.

Em primeiro lugar eu gostaria de deixar claro que por revolução eu estou falando em uma mudança nas relações sociais de produção, que consequentemente conduz a uma mudança no modo de produção. Isto se dá quando as forças de produção se desenvolvem de tal modo que as atuais relações de produção passam a ser incompatíveis com as forças de produção, servindo então como uma amarra de seu potencial. Neste momento apenas uma revolução nas relações de produção é capaz de libertar as forças de produção para o seu máximo potencial através da implementação de novas relações de produção mais compatíveis com o atual nível das forças de produção, configurando então um novo modo de produção.

De forma mais clara:

\begin{quote}
\begin{itemize}
\item \textbf{Modo de produção}: O método de produção dos bens e serviços necessários à vida (seja para saúde, alimentação, moradia ou necessidades como educação, ciência, nutrição, etc.). O Modo de Produção é a unidade das forças produtivas e das relações de produção. 

\item \textbf{Forças de produção:} As forças produtivas são a unidade de todo o trabalho, os instrumentos de produção (ex.: edifícios e máquinas) e os sujeitos de produção (ex.: matérias-primas. 

\item \textbf{Relações de produção:} As relações materiais objetivas que existem em qualquer sociedade, independentemente da consciência humana, formadas entre todas as pessoas no processo de produção social, troca e distribuição de riqueza material.
\end{itemize}
(Glossário de termos do marxists.org)
\end{quote}
Em segundo lugar, eu acredito que a gente vive uma situação particularmente diferente de todos países que eu conheço que sofreram revolução. Próximo a época de revolução, estes países viveram uma situação de intensa instabilidade, e não por conta dos comunistas. Ou seja, a instabilidade e o próprio cenário revolucionário que estes países se encontravam não foi produto dos comunistas, o que os comunistas fazem é guiar esta revolução.

Por exemplo, a China estava ocupada militarmente pelo Japão, já Cuba vivia sob a ditadura de Fulgêncio Batista, etc… Entendo que isto cria verdadeiramente uma sensação de que “não há nada a perder”. De certa forma, talvez a Rússia seja o país mais próximo do Brasil, já que ela não era colonizada. Mas ainda sim, o império russo estava envolvido em uma Guerra Mundial e sequer encontrava-se ainda com um modo de produção capitalista. Ou seja, não vivia sob um “capitalismo democrático” como o Brasil.

Mesmo se tentássemos ignorar os efeitos da Primeira Guerra Mundial, devemos lembrar que a assembleia constituinte formada entre a revolução de fevereiro (revolução burguesa) e a revolução de outubro (revolução socialista) constava com mais de 20\% dos deputados eleitos pertencendo ao partido bolchevique. Isto é, em um primeiro momento houve uma revolução burguesa, e no segundo momento, anterior à revolução socialista, os comunistas gozavam de cada vez mais popularidade. No momento não temos nenhum dos dois cenários no Brasil. Nem uma possível revolução e nem um apoio expressivo aos partidos marxistas.

Hoje o Brasil é relativamente estável politicamente. Além disso, o capitalismo aprendeu a ceder algumas migalhas em prol da estabilidade. Como resultado, a classe média cresceu e se tornou uma das principais forças políticas, chegando a rivalizar com a classe baixa em termos de tamanho (a depender da métrica utilizada) e possuindo mais poder político. Se a pequena burguesia já era uma classe vacilante na época de Marx, ela possivelmente é ainda mais nos dias de hoje devido ao poder econômico e político que ganhou. Dessa forma, eu vejo que a situação do Brasil hoje em dia se aproxima mais de países como Índia e Chile, do que Cuba e China.

\begin{quote}
A ascensão do capitalismo a partir do longo século XVI está associada a uma queda nos salários para níveis abaixo da subsistência, uma deterioração na estatura humana e um aumento na mortalidade prematura…

...Onde houve progresso, melhorias significativas no bem-estar humano começaram apenas por volta do século XX. Esses avanços coincidem com a ascensão de movimentos políticos anticoloniais e socialistas.

(Capitalism and extreme poverty: A global analysis of real wages, human height, and mortality since the long 16th century)
\end{quote}

Então como eu vejo a revolução sendo possível? O caminho que consigo visualizar é através da proposta de levar aos limites a democracia burguesa. Eu não vejo a possibilidade de convencer a população a expropriar os capitalistas enquanto as pessoas não defenderem sequer a existência de um salário mínimo, eu não vejo como vamos convencer a planificar a economia se as pessoas não defendem nem mesmo a existência de um imposto progressivo, então eu acredito que ainda existe um grande trabalho a ser feito dentro dos limites da democracia burguesa.

Eleições não são as nossas metas, mas servem como um termômetro sobre o quanto o público está comprando nosso projeto de sociedade. Se as pessoas não se sentem convencidas sequer a votar em um candidato comunista, quais as chances de apoiarem uma revolução comunista? Não vou dizer que sejam nulas, mas acredito que sejam baixas. Se vivemos em uma situação que é percebida como estável por grande parte da população, que mesmo em uma situação difícil conseguem ter a perspectiva de ter um planejamento para o futuro, se torna uma tarefa mais árdua convencer este trabalhador de que ele realmente não tem nada a perder com uma mudança radical de toda organização econômica e social da sociedade.

Então como lidar com uma população que é fortemente povoada por uma classe média com poder político e vacilante? Eu vejo que devemos seguir com certa precaução o exemplo do Salvador Allende, usar as vias legais tanto quanto possível. Se pudermos avançar em reformas, não há porque nos negarmos. Isso serve de propaganda das ideias comunistas de duas formas. Primeiro, quando os projetos são aprovados e executados serve como uma demonstração que nosso projeto de sociedade é de fato superior. E segundo, quando reprovado serve como uma de denúncia sobre a perseguição que as ideias comunistas sofrem no capitalismo. Nada é mais mentiroso de que “é fácil ser comunista no capitalismo”.

\begin{quote}
Sob as condições de que a classe dominante em um país capitalista esteja disposta a permitir que o Partido Comunista opere legalmente, o referido partido não rejeita a oportunidade. Afinal, o principal interesse de tal partido é elevar o nível de consciência do proletariado e de outras pessoas e organizá-las. Reformas também podem ser conquistadas de tempos em tempos.

(Jose Maria Sison em Basic Principles of Marxism-Leninism: A Primer)
\end{quote}

Eu não acredito que seja possível atingir o comunismo por esta via. Penso que a exemplo do próprio caso chileno, a burguesia tende a reagir violentamente. Mas é talvez neste momento que surge a instabilidade necessária para que a revolução se concretize e as relações de produção sejam modificadas. Um momento de instabilidade, onde a população sinta-se obrigada a escolher entre avançar com um projeto radical ou recuar para uma situação inferior, pode ser o momento determinante para que o capitalismo seja superado.

Quando até mesmo a classe média sentir que o seu projeto de sociedade está sendo atacado, este talvez seja o momento em que as condições necessárias para a revolução se concretizar se apresentam. Isto tudo sem abrir mão das possíveis reformas que surjam diante do caminho.

É importante notar que não tenho por meta uma nem uma revolução armada e nem um cenário de instabilidade. Não tenho por objetivo a violência de nenhuma forma. No melhor dos cenários o comunismo é atingido via reformas legais. Porém, no pior dos cenários, a população precisa reagir à resposta violenta da burguesia. É neste processo que eu tendo a ver a classe média como um importante aliado dos proletários em países capitalistas com uma democracia liberal.

Minha principal divergência com partidos sociais democratas, mesmo que adotem o discurso de que estão fazendo “tudo que é possível no momento”, se dá principalmente pelo horizonte reformista apresentado no discurso. Neste processo, eu defendo que o partido comunista tenha abertamente a intenção de superação do capitalismo, isto é, um horizonte revolucionário. As propostas devem seguir esse caminho, mesmo quando sabemos que serão rejeitadas, a rejeição deve servir de propaganda. O discurso deve se manter transparente, o objetivo não é tornar o capitalismo mais dócil, mas superá-lo. Ainda que, enquanto não superado, seja útil promover também projetos reformistas que garantam um alívio na classe trabalhadora.

Através de uma postura propositiva que inclua projetos de políticas públicas realistas, a atuação política das organizações comunistas devem permitir que as pessoas vislumbrem uma nova possibilidade de vivência e percebam que é possível ir além do capitalismo. De fato, nada disso é totalmente novo. Apenas desejo ressaltar a importância de utilizarmos os meios legais até onde for possível para promover o desenvolvimento da qualidade de vida da classe e das forças de produção, até atingirmos um ponto que uma mudança nas relações de produção torna-se uma possibilidade concreta. E isto tudo sem cair na inocência de acreditar que a burguesia não vai reagir.

Isso também implica que a questão central é a disputa da consciência de classe. Apesar de isto não ser novidade, talvez ainda estejamos negligenciando a importância desta disputa. Hoje em dia, as forças militares do estado são mais fortes do que em qualquer outro momento da história. O poder bélico estatal não tem adversários à altura que não seja outro poder militar estatal. Isso fica evidente, por exemplo, quando acompanhamos o genocídio que Israel comete na Palestina. Não há resistência armada civil capaz de parar Israel. Então mais do que nunca, é necessário o apoio da população, é a população que deve fazer resistência contra o ataque da burguesia. População, do qual, inclusive as forças policiais e militares fazem parte.

\begin{quote}
… os rebeldes Bolcheviques encontraram pouca resistência e foram capazes de penetrar no edifício sem muita dificuldade e tomá-lo. Na sua maioria, a revolta em Petrogrado ocorreu sem derramamento de sangue e numa atmosfera de normalidade geral na capital. Estima-se que apenas cinco marinheiros e um soldado entre os atacantes foram mortos no assalto e que os defensores não sofreram baixas fatais.

(Wikipédia)
\end{quote}

Esta necessidade de que a revolução seja feita pela massa de trabalhadores deve ficar muito clara. Deve estar claro para todas as pessoas que quem faz a revolução é, e sempre foi a classe trabalhadora. O povo. O partido comunista é apenas o guia da revolução. Mas não é o partido que faz a revolução. Não há como um pequeno punhado de marxistas fazerem uma revolução, ela só pode ser feita pela classe trabalhadora. 

O partido por estar inserido na classe trabalhadora deve se tornar capaz de reconhecer seus objetivos e anseios, e assim conseguir apresentar propostas efetivas, que são tomadas pela classe trabalhadora como um guia. Então torna-se nítido que a tarefa mais importante de qualquer organização comunista hoje, é estar inserida nos trabalhadores e disputar a consciência de classe dos trabalhadores.

\begin{quote}
Para triunfar, a revolução deve ser dirigida pela classe operária … A moralidade revolucionária consiste em unir-se às massas, em ter confiança nelas, em compreendê-las, em escutar atentamente suas opiniões. Por suas palavras e seus atos, os membros do Partido, da União da Juventude trabalhadora e os quadros ganham a confiança, a admiração e o afeto do povo, realizam a união estreita das massas em torno do Partido, organizam-nas, fazem a propaganda entre elas, mobilizam-nas, para aplicar com entusiasmo a política e as resoluções do Partido.

(Ho Chi Minh em Da Moralidade Revolucionária)
\end{quote}

São os trabalhadores que podem e devem fazer a revolução, acredito que nunca é demais enfatizar isso nos tempos atuais. Os trabalhadores devem desenvolver uma consciência acerca do poder que a própria classe possui, uma consciência acerca das contradições do capitalismo, e assim devem se tornar conscientes da necessidade de superação do capitalismo. Apenas através desse processo de educação que entendo que a classe trabalhadora seja capaz de aderir ao projeto político comunista e termine então por defender a revolução comunista.

Quando se discute sobre a via legal como uma ferramenta para a construção do comunismo, torna-se impossível não discutir o significado da eleição de candidatos que se declaram socialistas ou comunistas. Para além da crítica mencionada anteriormente devido a ausência do horizonte revolucionário de muitos candidatos, e consequentemente um discurso rebaixado, eu acredito que temos algo a aprender.

Pensando nos acontecimentos recentes, podemos refletir sobre a eleição nos EUA, especificamente para prefeito de Nova York, onde algumas coisas me chamam a atenção. Primeiro eu quero deixar claro que não o considero revolucionário, comunista, nem sequer socialista. Seu posicionamento político não é algo que possa ser comparado nem ao meu posicionamento político e nem copiado por comunistas/marxistas. 

Grande parte da campanha de Zohran Mamdani foi baseada em diferenciar o que ele entende de socialismo do que ele entende de comunismo, deixando claro que não é comunista. Ao mesmo tempo que se vende como uma campanha socialista, também promove uma campanha que se apoia fortemente no anticomunismo para traçar essa linha de separação entre ambos. Uma consequência desta prática, é que apaga-se a separação do capitalismo e socialismo, alterando o próprio sentido de socialismo. Isto faz com que a luta-anticapitalista sob esta bandeira do socialismo deixe de ser sobre a socialização dos meios de produção, e se torne meramente um capitalismo com menor desigualdade através de algumas políticas de redistribuição de renda.

\begin{quote}
Em vez do lema conservador de: “Um salário justo por uma jornada de trabalho justa!”, deverá inscrever na sua bandeira esta divisa revolucionária: “Abolição do sistema de trabalho assalariado!”.

(Karl Marx em Salário, Preço e Lucro)
\end{quote}

Mas dito isso, eu ainda acredito que há o que aprender desta experiência. Afinal é a maior cidade do mundo, e me parece evidente que a principal preferência burguesa nestas eleições saiu derrotada. De forma alguma eu estou negando o poder e influência da ideologia dominante, até porque quem venceu não foge totalmente dela, de certa forma, venceu dentro dos limites da defesa desta mesma ideologia.

Mas ainda sim, demonstra como as eleições são complexas, e sob certas condições, mesmo uma campanha com menos dinheiro pode derrotar o principal candidato da burguesia. Evidentemente a disponibilidade de recursos financeiros e a postura adotada do candidato diante da classe burguesa ainda é muito importante e claramente apontam tendências sobre quem serão os favoritos, mas sobra espaço para um certo grau de manobra e indecisão.

Além disso, este resultado tem como consequência reforçar o sentimento de que há espaço para melhora dentro do sistema legal no capitalismo, o que reforça também o sentimento de que há algo a perder em uma revolução. Evidentemente pode-se alegar que ele venceu porque tem apoio da pequena-burguesia, ou da classe média no geral. Longe de negar isto, é isto que estou querendo colocar em discussão: A influência e o poder político da pequena burguesia. Necessitamos olhar para ela, traçar estratégias para e com ela. Anteriormente, mencionei sobre a classe média ser a principal aliada, de forma mais precisa, a pequena-burguesia me parece ser atualmente uma das principais aliada do proletariado em um país como o Brasil.

Eu reconheço aqui que posso estar em uma posição privilegiada por eu mesmo ser classe média. Mas entendo que atualmente a gente precisa discutir o apoio da pequena burguesia e da classe média em geral. O que novamente não é novidade, esta questão sempre esteve presente nos textos marxistas, apenas reforça a necessidade de discutirmos isto mais claramente.

Concluindo, eu entendo que estas eleições recentes não foram uma derrota da burguesia pelo proletariado, o candidato eleito não pode trazer grandes mudanças estruturais pois independente das ideias que pessoalmente traz, sua atuação está limitada pelo próprio funcionamento da estrutura capitalista. Mas eu vejo como uma clara derrota da grande burguesia para a pequena burguesia, e isto é significativo. Evidentemente que a burguesia pode impedi-lo de implementar todas as propostas fazendo com que o episódio se converta em uma derrota para a pequena-burguesia (e possivelmente para toda classe trabalhadora). Mas esta é outra discussão.

O que eu queria chamar a atenção é que neste episódio, neste contexto das eleições, a grande burguesia sai derrotada. E isso não pode ser negligenciado. Como a pequena burguesia não tem aspirações revolucionárias, seus representantes também não terão. Mas ainda penso que esse episódio é algo no qual podemos tirar algumas lições.

\subsection{Revolução}

Muito se fala sobre a necessidade da disputa pela consciência de classe. Comentei que as eleições não devem ser completamente ignoradas. Mas exatamente de que forma se disputa a consciência de classe? Aqui abrimos um novo tópico em nossa discussão acerca da cultura. Eu defendo, e sempre defendi que a arte, e a cultura como um todo, possui um uso político que não pode ser negligenciado. Mas apesar disso, na prática, me parece que o movimento comunista não tem dado a devida atenção à este tema.

Acredito que em grande parte isto ocorreu devido ao fato de que movimentos marxistas anteriores viveram situações excepcionais que permitiram se darem o luxo de prestar menos atenção à essa questão cultural e ainda assim obterem revoluções de sucesso. Porém isto só foi possível em contextos no qual as sociedades já se encontravam em uma situação de flagrante instabilidade.

Eu tenho bastante convicção da necessidade da construção de uma contracultura, isto é, uma cultura que se oponha a cultura hegemônica capitalista. Exatamente por não termos um ambiente de instabilidade prestes a estourar uma revolução. Sendo assim, a prioridade atual não é sobre vencer uma revolução armada, mas a tarefa principal que nos resta é disputar a consciência dos trabalhadores através da construção de uma contracultura.

Devo chamar a atenção ao fato de que esta disputa de consciência política dos trabalhadores não se restringe a vídeos educacionais ou campanhas políticas diretas. Por melhores que sejam os professores envolvidos e por mais importante que estas tarefas sejam. Para mim, a prioridade do dia é toda a construção de uma cultura que incorpore essas críticas de forma natural. Isto é, a criação de uma cultura comunista, pensando na cultura como um conjunto de signos, valores e práticas. De forma que aqueles indivíduos que estejam imersos nessa cultura comunista, produzem e reproduzem naturalmente esta consciência política comunista que tanto buscamos difundir.

Como um exemplo vindo do lado oposto do espectro político, então neste caso não temos exatamente uma contracultura, mas mais uma subcultura capitalista mesmo, podemos pensar no movimento bolsonarista. Sua construção, principalmente através das redes de disseminação de notícias falsas em 2018, construiu uma cultura. Há signos compartilhados pelos bolsonaristas (exemplo: termos ‘petralha’ e ‘luladrão’), há valores compartilhados por bolsonaristas (exemplo: anticomunismo) e há também práticas compartilhadas por bolsonaristas (exemplo: a prática de compartilhar as falsas notícias que recebem). Foi construída uma subcultura bolsonarista em torno do candidato Bolsonaro.

Claro que a agitação e propaganda na sua forma mais direta, mais explicitamente política é importante, não estou negando isso. O que estou buscando discutir é que em tempos de ‘paz’, é insuficiente. É preciso reconhecer que há diferenças na construção do movimento revolucionário se vivemos em uma sociedade instável como a China ocupada pelo Japão Imperial ou se vivemos em um capitalismo estável como o Brasil. 

Não acredito ser coincidência que Gramsci, que ficou famoso por discutir exatamente a cultura, tenha vivido no contexto de uma Itália capitalista relativamente estável. Se pensarmos no contexto da China e dos locais onde outras revoluções socialistas ocorreram, a violência, a revolta armada era iminente, com ou sem organização comunista. Isto é, a revolta armada de certa forma não é produto dos partidos comunistas, eles apenas atuam na orientação desta revolta.

Mas o fato é que atualmente não vivemos nessa situação de revolta iminente, então precisamos adotar uma estratégia adequada. Se a ‘guerra de assalto’ (tomar todo território inimigo de uma só vez) não nos é possível, precisamos desenvolver uma estratégia adequada para uma ‘guerra de posição’ (a tomada de posição a posição do território inimigo).

Assim sendo, vejo a disputa pela consciência de classe como uma disputa cultural. E claro, uma disputa política, uma vez que não existe cultura apolítica. Temos uma disputa por valores, por afeto, pelo imaginário do trabalhador. Quando a pessoa comum, isto é, quem não é militante, se depara com um discurso ou um projeto político, ela reage a isso em grande parte baseado na cultura que carrega. Como ela vai receber este discurso depende exatamente da cultura no qual está inserida. Não que esta pessoa seja desprovida de raciocínio crítico e da capacidade de refletir sobre aquilo que vai contra suas crenças atuais, mas a concepção de mundo que o indivíduo carrega afeta sua recepção inicial a novas ideias.

E não se constrói uma cultura apenas com agitação e propaganda, apenas com política direta, é preciso arte. Arte no seu duplo papel de ensinar e divertir.

Lênin ensinava que a revolução socialista não é um único ato nem uma só batalha, mas toda uma época de agudos conflitos de classes, uma longa sucessão de batalhas em todas as frentes, isto é, em torno de todos os problemas da economia e da política. A fim de se colocar em condições de vencer a burguesia e de ganhar a batalha, \textbf{o proletariado deve travar, sempre e em toda parte, uma luta consequente e revolucionária} pela democracia.

\begin{quote}
“Quanto mais vigorosa é a força mecânica da reação, tanto mais se afrouxa a ligação com as massas, e tanto mais se inscreve na ordem do dia a tarefa da preparação da consciência das massas (e \textbf{não a tarefa da ação direta}); tanto mais se inscreve na ordem do dia a utilização dos caminhos da propaganda e da agitação criados pelo velho poder (e não o ímpeto direto das massas contra esse mesmo velho poder).” (Op. cit., pág. 144)

(Lenin: vida e obra)
\end{quote}

Ou seja, nem tudo se reduz à “guerra de assalto” como diz Gramsci, ou a “ação direta” como diz Lênin, a revolução não é construída unicamente através da revolta armada ou ações diretas. A luta deve ser travada em todas as frentes, e isto para mim, implica necessariamente que é também, e talvez principalmente, uma disputa pela hegemonia cultural. A tarefa da agitação e propaganda de disputar a consciência política, pode ser colocada de forma mais geral como a tarefa de construir uma cultura que promova e incorpore a consciência política que desejamos promover entre os trabalhadores.

Para ressaltar a importância da cultura, podemos recorrer a Bakhtin, que nos propõe uma relação direta entre a vida social e nossa própria atividade consciente individual. Para Bakhtin,a vida social não é uma mera exteriorização da vida interna, na verdade, a vida individual interna é a vida social interiorizada. Ou seja, a consciência política que queremos promover é internalizada no indivíduo a partir do momento que o indivíduo  desenvolve sua vida social dentro do contexto de uma dada cultura que promova esta consciência. 

\begin{quote}
A consciência individual não é o arquiteto dessa superestrutura ideológica, mas apenas um inquilino do edifício social dos signos ideológicos…A consciência individual é um fato sócio ideológico.

(Mikhail Bakhtin em marxismo e Filosofia da Linguagem)
\end{quote}

Em outras palavras, para mim, disputar a consciência política dos trabalhadores, é, exatamente, construir uma contracultura comunista adequada. O momento atual da nossa revolução, é então uma guerra de posição contra a hegemonia cultural capitalista.

\subsection{Cultura}

Pela centralidade com o qual eu estou colocando a cultura em discussão, pode-se correr o risco de achar que estou abrindo mão da centralidade da base econômica na análise da sociedade. Eu gostaria de tranquilizar o leitor que esta não é minha intenção. E é sobre isto que quero discutir agora. A cultura não é um mero reflexo da base econômica, ao mesmo tempo que tem uma importância vital na sustentação do capitalismo ela não deixa de ser determinada pela base econômica no qual esta sociedade se sustenta.

Minha intenção, é então, defender a seguinte afirmação:

\textit{A cultura não é um reflexo secundário do modo de produção, mas parte constitutiva da estrutura social na qual as relações de produção existem; ainda assim, ela permanece historicamente determinada (limitada) pelo modo de produção.}

Primeiro, eu gostaria de fazer uma definição:

\textbf{Cultura}: Um conjunto historicamente produzido de signos, valores e práticas que mediam toda atividade social, organizando a experiência humana. Lembrando que toda atividade social é mediada por signos que nunca são neutros.

Apesar de já termos expressado anteriormente, eu gostaria de enfatizar novamente o significado de alguns termos centrais no marxismo:

\begin{quote}
\begin{itemize}
\item \textbf{Modo de produção}: O método de produção das necessidades básicas da vida (seja para saúde, alimentação, moradia ou necessidades como educação, ciência, cuidados infantis, etc.). O Modo de Produção é a unidade das forças produtivas e das relações de produção. A produção começa com o desenvolvimento de seu aspecto determinante – as forças produtivas – que, uma vez atingido um certo nível, entram em conflito com as relações de produção dentro das quais se desenvolveram… O conflito entre as forças produtivas e as relações de produção é a base econômica da revolução social.

\item \textbf{Forças de Produção}: A união dos meios de produção e do trabalho. Os meios de produção são as ferramentas (instrumentos) e a matéria-prima (material) que você usa para criar algo.

\item \textbf{Relações de produção}: As relações materiais objetivas que existem em qualquer sociedade independentemente da consciência humana, formadas entre todas as pessoas no processo de produção social, troca e distribuição da riqueza material. A produção não é possível sem relações de produção – os seres humanos não podem produzir fora de uma estrutura social, seja uma nação ou uma família – as relações de produção existem para todos os produtores.
\end{itemize}
(Glossário de Termos)
\end{quote}

Se unirmos todos esses conceitos, já podemos tirar algumas conclusões interessantes.

Toda proposta revolucionária do marxismo consiste no seu objetivo em promover a superação do modo de produção capitalista. Uma vez que o que define se o modo de produção é capitalista ou comunista são as relações de produção, a revolução socialista não é outra coisa senão uma mudança nas relações de produção. E isto é independente da existência de uma revolução armada ou não. Assim sendo, as relações de produção ocupam uma posição central no debate.

Uma vez entendido isso, para avançar na reflexão podemos sintetizar três ideias já expostas:

a) Os seres humanos não podem produzir fora de uma estrutura social.

b) Toda atividade social é mediada por signos não neutros.

c) Definimos a cultura como um conjunto de signos historicamente produzidos.

Com isso em mente, devemos perceber que a cultura não é um mero reflexo do modo de produção, mas ela é parte da estrutura social no qual as relações de produção - e consequentemente o modo de produção - existem. Nem por isso a cultura deve ser vista como autônoma, ela é determinada pelo modo de produção. Devemos entender que por “determinada”, o que queremos dizer é que o modo de produção impõe um limite nas possibilidades de desenvolvimento cultural que podem emergir nesta sociedade. Isto ficará mais claro com alguns exemplos posteriores.

Também não podemos esquecer que nem toda relação social é uma relação de produção. Relações de produção possuem um significado mais estrito. No glossário de onde as definições foram retiradas podemos ler ainda que “\textit{a base das relações de produção é a propriedade dos meios de produção}”. É importante não cairmos em um ‘culturalismo’ buscando combater o economicismo. Devemos lembrar que o modo de produção ainda determina (limita) a cultura de forma que a revolução socialista não pode ser realizada se perdermos de vista da posição determinante do modo de produção.

Deve-se lembrar também que o modo de produção é a unidade das forças e relações de produção e esta unidade é dialética, isto é, estes elementos não existem de forma separada e independente. Não é uma mera união mecânica de duas coisas separadas. As relações de produção não existem independente das forças de produção, assim como também a expressão cultural de uma sociedade não é independente do seu modo de produção. Isto deve ficar claro, pois esta discussão insere a cultura no debate marxista, mas não retira o modo de produção de uma posição central deste mesmo debate.

O que estou buscando deixar claro é que a cultura é mais que um reflexo passivo do modo de produção como pregado pelo economicismo. E também é mais que um mero fenômeno secundário, como muitas vezes acaba sendo expresso mesmo por alguns críticos do economicismo que recorrem à clássica divisão entre infraestrutura e superestrutura. A cultura pode ser vista como algo que atravessa o modo de produção uma vez que relações de produção também são relações sociais. Não há um modo de produção que exista a parte de qualquer cultura.

Isto é, a estrutura social no qual as relações de produção existem possuem uma dimensão cultural no qual o modo de produção se insere. Como toda atividade social é mediada por signos que nunca são neutros, e estes são produzidos socialmente dentro de uma dada cultura hegemônica, a cultura pode ser entendida como um dos mecanismos responsáveis por manter o modo de produção capitalista estável de forma “democrática”.

E aqui consiste uma das principais características do capitalismo, o senso de liberdade e escolha democrática construído pela cultura hegemônica. Sob certos limites, podemos afirmar que grande parte dos trabalhadores aceitam se submeter à exploração capitalista voluntariamente, pois esta exploração foi (através da cultura hegemônica) naturalizada.

A opressão em outros modos de produção era mais direta, haviam guardas fisicamente coagindo os escravizados para produzirem coisas para o senhor escravista. Era uma opressão que acontecia diretamente por meios militares. A opressão era física, visível, explícita. No feudalismo pode-se argumentar que existia ainda uma questão cultural em justificar qual era a família real. Há um aspecto cultural na determinação de porque era aquela família a real e não outra, mas uma vez que se aceita isso, a dominação se dá novamente de forma bastante direta. É tal família que detém o poder político e econômico, é esta família que possui o sangue real, uma família camponesa não pode se tornar uma família real.

Esta separação das duas esferas (político-jurídico) e econômica, de modo que há também uma modificação nos meios empregados para a exploração (de meios extra-econômicos para meios econômicos) está na origem do próprio capitalismo. Ellen Wood discute esta etapa de transição em seu livro. Comentando as ideias de Perry Anderson, Wood comenta:

\begin{quote}
O estado absolutista era essencialmente feudal, insiste ele, porque representou o deslocamento para cima e a centralização dos poderes coercitivos político-jurídicos dos senhores feudais, separando-os da exploração econômica. Dito de outra maneira, o Estado absolutista separou os dois momentos da exploração - o processo de extorsão do excedente de um lado, e o poder coercitivo que o sustentava de outro.

(A Origem do Capitalismo)
\end{quote}

Lembrando que esta distinção entre os diferentes meios de coerção adotado pelas sociedades pré-capitalistas e capitalistas também está presente desde Marx:

\begin{quote}
No evolver da produção capitalista desenvolve-se uma classe de trabalhadores que, por educação, tradição e hábito, reconhece as exigências desse modo de produção como leis naturais e evidentes por si mesmas … a coerção muda exercida pelas relações econômicas sela o domínio do capitalista sobre o trabalhador. A violência extraeconômica, direta, continua, é claro, a ser empregada, mas apenas excepcionalmente.

(Karl Marx em O Capital 1)
\end{quote}

Podemos comparar esta situação imposta pelo capitalismo, com uma situação híbrida, onde o capitalismo ainda não se tornou o modo de produção dominante.

\begin{quote}
Basta, aqui, uma simples alusão a formas híbridas, em que o mais-valor não se extrai do produtor por coerção direta e que tampouco apresentam a subordinação formal do produtor ao capital. Nesses casos, o capital ainda não se apoderou diretamente do processo de trabalho

(Karl Marx em O Capital 1)
\end{quote}

Fica então evidente um claro contraste entre a coerção direta de formas pré-capitalistas e a coerção indireta no capitalismo.  Merece uma atenção especial o fato de que a classe de trabalhadores aceita o modo de produção através da “\textit{educação, tradição e hábito}”. Esta ideia reaparece no mesmo livro,  quando é feito um comentário sobre as contribuições de Thompson.

\begin{quote}
Em sua obra o desenvolvimento do capitalismo ganha vida não apenas como um processo de proletarização, particularmente em seu clássico \textit{A formação da classe trabalhadora inglesa} (1963), mas também como um confronto vivo entre os princípios do mercado, as práticas e valores alternativos…

(A Origem do capitalismo)
\end{quote}

Ou seja, para que o capitalismo possa surgir, temos um confronto vivo entre valores e práticas. Se lembrarmos que definimos cultura como um conjunto de signos, valores e práticas, podemos presumir que foi necessário a formação de uma nova cultura. Para complementar esta discussão, podemos pegar um exemplo para abordar a questão dos signos, um elemento de nossa definição de cultura que está ausente nesta citação. Enquanto discorre sobre o desenvolvimento do capitalismo Wood chama a atenção do leitor para o fato de que a palavra “melhorar” em sua acepção original significava literalmente “\textit{fazer alguma coisa com vistas ao lucro monetário}”, então propõe a seguinte reflexão:

\begin{quote}
… essa palavra foi adquirindo um significado mais geral, no sentido como a conhecemos hoje (e seria interessante pensar nas implicações de uma cultura em que a palavra correspondente a “tornar melhor” enraíza-se no termo que corresponde a lucro monetário).

(A Origem do capitalismo)
\end{quote}

Há também uma breve discussão análoga sobre as mudanças no significado da palavra “produtor” com o advento do capitalismo. Assim, não deve nos surpreender que Thompson, que segundo Wood novamente, teria dito que um dos momentos axiais da formação do capitalismo seja o “\textit{o processo em que foram forjados um novo proletariado e uma nova cultura da classe trabalhadora.}” 

Retomando então a comparação entre o modo de produção capitalista e os modos pré-capitalistas, minha impressão é de que nestes sistemas sempre foi explícita a separação entre quem são os opressores e quem são os oprimidos. Se você estivesse na base da pirâmide social, e convencesse todos ao seu redor que as coisas deveriam mudar, ainda sim seria irrelevante. Claro, não seria irrelevante se esta mobilização resultasse em uma revolta, mas isso implica que não se podia nutrir a esperança de que fosse possível mudar as coisas “dentro do sistema”. Essas possibilidades não estavam previstas nem como uma mera possibilidade teórica. Há uma imposição explícita imposta de cima para baixo, a mobilidade social é abertamente inflexível. E entendo que isso opera um pouco diferente no capitalismo.

No capitalismo atual, a opressão ocorre de forma mais indireta e mediada. Ela se dá pela desigualdade entre quem necessita de dinheiro e quem detém o dinheiro. Mas essa dinâmica é revestida de uma aparência neutra. Não há um exército que force o produtor a entregar sua produção para o seu senhor, mas nasce uma imposição análoga devido a atuação das “forças econômicas”, da necessidade que o trabalhador tem de vender a sua força de trabalho. Porém agora, em teoria, qualquer um pode se tornar explorador, é principalmente a estrutura social que mantém essa distinção entre explorados e exploradores. A atividade social no capitalismo é mediada também pelo dinheiro, e em teoria, qualquer um pode ter dinheiro.

Essa exploração nasce então de “comum acordo”, nos submetemos a um trabalho assalariado de forma “livre”. E não é só isso, aparentemente o “poder” foi deslocado da família real ou do escravagista para o “estado”, que agora adquire ares de uma entidade abstrata que não discrimina os cidadãos. Na prática, o estado é uma ferramenta de opressão sendo utilizada classe dominante, mas em teoria, de acordo com a constituição, podemos eleger qualquer um para ocupar uma posição no estado e este representa os interesses de toda sociedade. Não há uma limitação teórica para a escolha dos representantes.

Isto se diferencia dos sistemas anteriores. Agora as limitações precisam ser impostas de forma mais implícita. Em teoria poderíamos eleger um congresso inteiro de comunistas. Em teoria, se eu convencer todos os explorados iguais a mim a eleger apenas candidatos que defendam nossos interesses, o estado poderia ser usado também  para atender nossos interesses.

Tudo isto é uma possibilidade que não existia em sistemas anteriores. Não conheço nenhum sistema no feudalismo que permitisse pacificamente mudanças tão grandes. Evidentemente na prática não é tão simples de levar isso às últimas consequências, mas um dos motivos de existir essa possibilidade teórica no capitalismo, e apesar disso ela nunca se concretizar, para mim, é a existência da questão cultural. Questão que é determinada pela base material e econômica da sociedade, deve-se ressaltar.

O sistema permite que as pessoas escolham quem vai “deter o poder” através do estado, mas ao mesmo tempo, a classe dominante mantém um controle cultural sobre os gostos, valores e afetos da população. Ou seja, as pessoas são livres para votar em quem vão representá-las, mas seus gostos são construídos seguindo os interesses da classe dominante. Um movimento que se dá através da cultura.

Não basta reconhecer que o estado representa o interesse da classe dominante, mas devemos discutir como se sustenta. Pois a não ser que defendamos “fraude nas urnas”, é inegável que milhões de pessoas votam e participam do processo eleitoral. Podemos alegar que se elas elegerem alguém radical de verdade a burguesia iria derrubar com um golpe, podemos lembrar o Allende no Chile, mas na prática nunca chegamos a isso. As pessoas sequer elegem os candidatos radicais, e devemos entender como isso ocorre.

É esta reflexão, que para mim, resulta na constatação de que um dos grandes focos de poder no capitalismo, se dá na dimensão da cultura. Isto é, na construção dos nossos valores e gostos, na construção do nosso imaginário. Uma das grandes ferramentas de controle e organização do capitalismo é sua capacidade de criar uma cultura que perpetua e justifica a exploração. A força militar surge de forma secundária principalmente em momentos de crise, mas não é a principal força atuante no cotidiano para a maior parte da população.

A indústria no capitalismo não se reduz a “indústria material”, aquela que tipicamente pensamos quando pensamos em indústria, ou seja, aquela indústria que produz bens materiais. Mas a indústria capitalista inclui, e não com menor importância, a indústria cultural, isto é, a indústria que produz  cultura. Como exemplo desta indústria, podemos lembrar da indústria de Hollywood estado-unidense exportando uma cultura atrelada ao “\textit{American Way of Life}”, ou então o fenômeno global de popularização da cultura sul-coreana chamada de onda coreana “\textit{Hallyu”}, também poderíamos citar o Japão e o \textit{soft-power} construído a partir da indústria de animes e mangás. Nenhuma destas indústrias existe sem uma base material e econômica que torne elas possíveis, e nenhuma delas existe também sem uma participação ativa do estado que age em prol do interesse das classes dominantes.

Para ficar mais claro, podemos trabalharmos com uma definição de cultura hegemônica:

\begin{quote}
\textbf{Hegemonia cultural}: refere-se ao domínio de um grupo cultural sobre outros dentro de uma sociedade. Esse domínio não é imposto apenas por meios políticos ou econômicos, mas é estabelecido e mantido por meio de instituições culturais, como a mídia, a educação e a religião. A hegemonia cultural molda percepções, normas e valores, muitas vezes levando à marginalização de perspectivas e estilos de vida alternativos.

(adaptado de Oxford Review)
\end{quote}

Ou seja, a classe dominante produz a cultura hegemônica (dominante). Esta cultura abrange e expande conforme necessário um conjunto de signos, valores e práticas que justificam e sustentam a manutenção desta mesma classe no poder. Não foi ao acaso que Marx, na Ideologia Alemã declarou que “\textit{as ideias da classe dominante são as ideias dominantes}”. Esta dominação é construída através da dimensão cultural da sociedade, e isto é construído também materialmente, uma vez que a cultura tem sua dimensão material, não é um mero campo de ideias descolado da realidade material. Isto é, a cultura precisa se manifestar também de forma material.

Se quisermos manter a distinção de infraestrutura e superestrutura podemos recuperar a citação original de Marx:

\begin{quote}
A totalidade destas relações de produção forma a estrutura económica da sociedade, a base real sobre a qual se ergue uma superestrutura jurídica e política, e à qual correspondem determinadas formas da consciência social. O modo de produção da vida material é que condiciona o processo da vida social, política e espiritual.

(Karl Marx em Para a Crítica da Economia Política)
\end{quote}

Mas é importante aqui perceber que por superestrutura, não é qualquer coisa que não corresponda a infraestrutura, mas a estrutura jurídica e política que organiza e defende a base econômica. Se quiser expandir um pouco a ideia de superestrutura, podemos lembrar por exemplo da discussão de Stálin sobre a linguagem não ser uma superestrutura:

\begin{quote}
A INFRA-ESTRUTURA é o regime econômico da sociedade numa etapa determinada de seu desenvolvimento. A super-estrutura são as opiniões políticas, jurídicas, religiosas, artísticas, filosóficas da sociedade e as instituições políticas, jurídicas e outras que lhes correspondem.

Toda infra-estrutura tem sua superestrutura correspondente…Se a infra-estrutura se transforma e desaparece, ela acarreta a transformação e o desaparecimento de sua super-estrutura; se nasce uma infra-estrutura nova, ela acarreta o nascimento da super-estrutura que lhe corresponde.

(Stálin em Sobre o marxismo na Linguística)
\end{quote}
Ainda que seja uma concepção mais polêmica, de maneira nenhuma a cultura como um todo se confunde como uma mera superestrutura mesmo nesta concepção. O que entendo, é que podemos pensar que a cultura hegemônica em particular age como uma superestrutura auxiliando na manutenção do modo de produção. Mas ainda devo chamar a atenção que mesmo a cultura hegemônica não pode ser vista como um mero reflexo ou produto passivo da base econômica, uma vez, que como foi dito, ela atravessa o próprio modo de produção.

Assim sendo, as próprias relações de produção que ocupam uma posição especial na estrutura social não são exteriores a esta estrutura, não podem existir fora de qualquer cultura. Ainda que costumeiramente ressaltamos o aspecto econômico das relações de produção e do modo de produção, estes também possuem um aspecto social e cultural intrínsecos.

As relações de produção são parte da vida social no qual a sociedade se organiza. Assim sendo, a cultura influencia como construímos a própria vida social no qual a sociedade se organiza e se mantém.

Não é de fato uma discussão inovadora. Reconhecidamente, uma das grandes contribuições de Marx para o debate da economia política foi reconhecer o aspecto social do trabalho. O próprio “valor” conforme exposto por Marx, só existe em uma relação de troca, e não como uma propriedade intrínseca do bem produzido.

Podemos trabalhar com alguns exemplos grosseiros. Ainda que o capitalismo imponha determinadas relações de produção, a postura que adotamos diante destas relações é menos definida. O senso comum promovido pela classe dominante que detém a hegemonia cultural visa promover uma percepção da realidade social que justifique a exploração. Isto pode ser expresso por exemplo, nas palavras (signos) escolhidas para abordar determinados assuntos.

Optar por chamar um trabalhador de assalariado ou ‘associado’ transmite diferentes ideias da posição que este trabalhador ocupa na sociedade. Enquanto o assalariado transmite uma ideia mais fria, ‘associado’ pode ser uma escolha de palavras mais adequada dentro de um discurso que visa convencer o funcionário a ‘vestir a camisa da empresa’ e aceitar de forma mais passiva a exploração no qual é submetido.

Às vezes a cultura hegemônica se revela através de discursos mais diretos. Não é incomum nos depararmos com a exposição da ideia de que os trabalhadores deveriam ser gratos aos burgueses por terem um trabalho. Reproduzindo assim uma visão de mundo onde os burgueses não apenas não são culpados pela exploração que promovem, mas que os explorados lhe devem gratidão. 

Em ambos exemplos discutidos, a relação de produção é uma só, mas a forma com que percebemos esta relação não é. A cultura hegemônica capitalista e uma contracultura comunista vão reunir um conjunto de signos, valores e práticas que vão levar a diferentes práticas e percepções sobre um mesmo evento. Evidentemente, enquanto um visa tornar os trabalhadores passivos e aceitarem o sistema atual, o outro visa tornar o trabalhador um agente ativo na revolução deste sistema.

Talvez um caso mais prático e atual que ilustra minha posição se dá nos próprios movimentos sociais. Podemos citar como exemplos os movimentos feminista e antirracista.

\begin{quote}
Não há luta contra o racismo sem luta contra o capitalismo. (Mamadou Ba para Esquerda.net)

A luta anticapitalista não é possível sem a questão racial. (Dennis Oliveira para SINPRO-DF)
\end{quote}

Podemos dizer, sem prejuízo de diminuirmos a importância da questão, que a superação do racismo como um problema estrutural demanda uma mudança cultural. Novamente, é preciso reforçar que cultura possui uma dimensão material, não é uma mera questão de ideias. Portanto, dizer que é uma mudança cultural não é defender que deve-se apenas mudar as ideias, mas inclui necessariamente mudanças materiais, como por exemplo mudanças políticas e jurídicas, ou de forma mais ampla, mudanças em instituições que também são responsáveis por produzir cultura.

É preciso sempre ter em mente que a cultura é produzida por seres humanos, portanto não é algo que paira sobre a humanidade, uma mudança cultural demanda necessariamente uma mudança material e social na sociedade onde vivem estes homens que reproduzem esta cultura. Podemos constatar então que o movimento antirracista esta ciente de que a mudança cultural que a luta antirracista exige, não está desconectada de uma mudança econômica de superação do capitalismo.

Nem é nem apenas uma mudança econômica onde a mudança cultural será um mero reflexo (ou seja, a ideia de que superar o capitalismo levaria automaticamente a superar o racismo), e nem apenas uma mudança cultural que pode ser feita independentemente de uma mudança econômica (ou seja, que o racismo pode ser superado sem a superação do capitalismo). Na prática então, o movimento antirracista já reconhece tudo que eu estou tentando argumentar.

Podemos ver então como o modo de produção capitalista determina a cultura hegemônica racista. Isso não implica que o modo de produção comunista automaticamente supera o racismo, mas ele expande os limites de modo que abre a possibilidade de que o racismo seja superado, algo que não é possível no capitalismo.

\textit{As relações sociais de produção, embora materialmente determinadas, são culturalmente mediadas. Logo, a base material é produzida pela atividade econômica e limita a cultura; entretanto, sua manutenção exige a produção contínua de valores ideais, significados e afetos que são produzidos culturalmente.}

\textbf{Moralidade}

Devo ressaltar, mais uma vez, que a questão cultural para mim ocupa uma posição central na superação do capitalismo. Isto nos conduz a algumas questões, como por exemplo, como produzir uma cultura comunista? É possível começar a produzir esta cultura quando ainda estamos em uma sociedade capitalista?

Produzir uma “economia comunista” no sentido de elaborar apenas um plano técnico de produção levando em conta a tecnologia, matéria prima e ferramentas disponíveis, considerando que todo mundo já concorda com o comunismo, não é o maior desafio, é uma questão meramente técnica. Podemos lembrar novamente do “Towards a new Socialism”, ou então do “Ciber-comunismo: planificación económica, computadoras y democracia”, ambos do mesmo autor.

Não que seja simples, mas é uma questão técnica mais fácil de se trabalhar e não é o principal desafio para superar o capitalismo. A ciência e tecnologia necessárias já existem. Agora como produzir uma cultura no qual as pessoas estejam comprometidas em construir e manter esse modo de produção comunista? Não são questões separadas, mas talvez, muitas vezes estejamos negligenciando essa necessidade.

Porém, de nenhuma forma eu mereço o mérito de ter sido o primeiro a trazer estas questões para debate. Alguns autores já têm pensado nisso de uma forma ou outra (podemos citar Terry Eagleton, Raymond Williams, Antonio Gramsci, Bertolt Brecht, Augusto Boal, etc), além disto, líderes comunistas sentiram na prática a necessidade de abordar o aspecto cultural na superação do capitalismo, e é isto que eu gostaria de comentar rapidamente.

Podemos começar lembrando do fato de que a China realizou uma revolução cultural onde visava literalmente “revolucionar” a cultura, construindo e consolidando uma “nova cultura proletária” em detrimento de uma “antiga cultura burguesa”. Discute-se o sucesso desse esforço e os métodos que foram empregados, mas dificilmente comunistas discordam da necessidade de realizar uma revolução cultural. Há quem se posicione que um dos problemas da URSS que levou ao fim do bloco soviético tenha sido exatamente a ausência de uma revolução cultural mais contundente. Ou seja, um dos possíveis motivos para a dissolução do bloco soviético pode consistir exatamente no aspecto cultural.

Isso reforça toda nossa discussão, não podemos ver a cultura como um mero reflexo ou algo secundário ao modo de produção, uma vez que aparentemente, sem uma cultura hegemônica adequada, um modo de produção pode falhar em se estabelecer. A experiência tem nos mostrado que não é possível esperar realizar primeiro uma revolução política e econômica, para futuramente realizar uma revolução cultural de forma separada, todas estas dimensões (política, econômica e cultural) precisam ser modificadas de forma conjunta.

E isto não é, de fato, algo que passou despercebido pelo movimento comunista. A construção desta cultura pode, e deve, começar no próprio interior dos partidos e da organização proletária em geral. Lenin defende a utilidade de sindicatos e greves como uma experiência que permite aos trabalhadores aprenderem a se organizar e lutar na prática, isto é, se me permitem chamar assim: desenvolver uma cultura organizativa. Mas acredito que podemos expandir isso para outros aspectos culturais.

\begin{quote}
Assim, as greves ensinam os operários a unirem-se; as greves fazem-nos ver que somente unidos podem agüentar a luita contra os capitalistas; as greves ensinam os operários a pensarem na luita de toda a classe patronal e contra o governo autocrático e policial. Exatamente por isso, os socialistas chamam as greves de "escola de guerra", escola em que os operários aprendem a desfechar a guerra contra seus inimigos, pola emancipaçom de todo o povo e de todos os trabalhadores do jugo dos funcionários e do jugo do Capital.

(Lenin em Sobre  as Greves)
\end{quote}

Ho Chi Minh (Da Moralidade Revolucionária) por exemplo, fala da necessidade dos membros do partido abandonarem o individualismo, que estes deveriam servir como modelos para os trabalhadores, isto é, na prática, a construção de uma nova cultura de valores coletivos em detrimento de uma cultura individualista. Enver Hoxha (A Luta Ideológica e a Educação do Homem Novo) vai falar da necessidade de reforçar a “consciência comunista” das massas trabalhadoras, o que é isto se não a construção de uma nova cultura? Che Guevara (Socialismo e o Homem em Cuba) vai dizer que precisamos criar o Homem do século XXI. Kim Il Sung (Sobre a Eliminação do Dogmatismo, do Formalismo e o Estabelecimento do Juche no Trabalho Ideológico) vai discutir sobre a necessidade de educar o povo e inspirar neste povo um amor ardente ao seu lugar de origem, ideia que está nos pilares da Ideia Juche. Certamente poderíamos seguir citando outros revolucionários, mas para concluir, quero deixar apenas uma última citação:

\begin{quote}
Nos últimos vinte anos, para onde quer que essa nova força cultural tenha levado os seus ataques, verificou-se uma grande revolução, tanto no conteúdo ideológico como na forma (por exemplo na linguagem escrita). A sua influência é tão grande e o seu impacto tão poderoso que se afigura invencível onde quer que chega. A mobilização a que se procedeu ultrapassou a de qualquer outro período na China. Lu Sun foi o maior e mais corajoso porta-bandeira dessa nova força cultural. Comandante em chefe da revolução cultural chinesa, ele não foi apenas um grande homem de letras, foi também um grande pensador e um grande revolucionário. Lu Sun foi um homem de coluna vertebral tesa, sem sombra de servilismo nem obsequiosidade — qualidade inestimável dos povos das colónias e semi-colónias. Na frente cultural, Lu Sun foi o representante da grande maioria da nação, o herói nacional mais correto, mais bravo, mais firme, mais fiel e mais ardente, um herói que abriu brecha e arrasou a cidadela do inimigo, um herói sem igual. A via de Lu Sun foi justamente a via da nova cultura nacional da China.

(Mao em Sobre a Democracia Nova)
\end{quote}

Como isso deve ser feito exatamente é uma questão em aberto, porém não há discussão sobre que isso precisa ser feito. E certamente precisa ir além do ambiente interno do partido, neste sentido devo destacar o Lu Sun sendo elogiado por Mao. Ou trabalhando com um exemplo mais nacional, lembrar do Augusto Boal e da sua posição de que se o teatro não é revolucionário, ele é pelo menos um ensaio para a revolução. Brecht também é conhecido por defender uma posição onde vê no teatro, e na arte no geral, uma ferramenta política a ser utilizada pelo movimento comunista.

Evidentemente, há uma limitação pela situação material que vivemos. Não estou dizendo que a partir de hoje só deveríamos consumir arte comunista, mas somos influenciados por tudo que consumimos e vivemos deve ser unicamente “arte comunista”, apenas que na medida do possível uma revolução que vise a superação do capitalismo, deve visar também a construção de um ambiente cultural no qual determinados valores são difundidos ou combatidos de acordo com o interesse revolucionário. Toda arte sempre é, e sempre foi política. Não é apenas sobre construir a tecnologia necessária para que seja possível um novo modo de produção e distribuição de bens, mas também sobre a construção de novos valores e afetos, de uma nova concepção de mundo comunista, e consequentemente, de um novo modo de vida comunista.

\begin{quote}
Os reacionários – especialmente os fascistas – procuram evitar que essa situação seja devidamente compreendida, procuram escamotear a significação política global da arte e até procuram estetizar a política… Essa estetização da política (com o embelezamento fascista da guerra) exige uma réplica: “A resposta do comunismo é a de politizar a arte”… o problema estético surge quando eles se perguntam como a arte deve se politizar.

(Bertolt Brecht em Os Marxistas e a Arte)
\end{quote}

\subsection{Teorias}

Eu tenho exposto a minha perspectiva de de uma teoria da cultura materialista, desenvolvida a partir da minha compreensão do materialismo dialético, buscando não importar conceitos externos. A forma com que eu abordo tem inspiração principalmente em nomes como Bakhtin, Gramsci, Raymond Williams, Brecht e Boal, mas esta não é a forma mais popular de abordar a cultura adotada pelos marxistas, então eu gostaria de fazer um rápido comentário sobre isto. 

Em um primeiro momento, eu concebi duas grandes abordagens diferentes da cultura no interior do marxismo. A primeira abordagem que eu gostaria de comentar, é aquela que me parece ser a mais conhecida. Esta é a abordagem que tipicamente é lembrada quando se fala em cultura sob uma perspectiva marxista. Esta abordagem mistura psicanálise como marxismo, e se não estou enganado surgiu principalmente com a escola de Frankfurt.

Suponho eu que ela surge como necessidade de discutir a questão cultural na sociedade uma vez que a prática militante mostrou uma crescente importância da luta que se dá nesta dimensão em um cenário onde não há um horizonte revolucionário próximo.  Neste contexto , em uma Europa onde o Freud é um nome que estava em alta, me parece que os marxistas viram uma oportunidade em recorrer a psicanálise para complementar o marxismo e expandir a discussão para a esfera da cultura.

Já vi mais de um autor colocar Freud e Marx como dois grandes pensadores que 'abalaram' as crenças da humanidade em diferentes áreas. Mas esta não é a única alternativa possível de realizar esta discussão. E na minha percepção, não é sequer a mais adequada.

Evidentemente, eu não discordo da necessidade de abordar a cultura. E menos ainda que a discussão mais tradicional do marxismo não esgotou o assunto, mas eu discordo da abordagem adotada por este freudo-marxismo. Evidentemente, há uma necessidade de maior estudo para elaborar uma crítica mais elaborada desta abordagem da cultura, mas me parece que a maior diferença entre as duas abordagens se concentra em onde a teoria da cultura desenvolvida por cada abordagem coloca a ênfase.

Enquanto a psique individual parece ser o foco do Freud e consequentemente daqueles que adotam a psicanálise como elemento constituinte da teoria da cultura, Gramsci e os outros marxistas vão colocar a ênfase na vida social. Nesta segunda abordagem eu chamarei de materialismo histórico cultural, ou apenas materialismo cultural para facilitar.

Me parece que esta primeira abordagem, agrupada sobre o nome de freudo-marxismo tende a recorrer a psique interna para explicar o comportamento social exibido pela sociedade. Muitas vezes, quando leio seus autores, me parece que a cultura e a vida social de um dado grupo é vista sempre como a expressão de uma atividade interna. O risco desta perspectiva é que inclusive os conflitos sociais passem a ser psicologizados e serem vistos como frutos de um conflito interno.

Ainda que a atividade interna da psique humana não seja vista como totalmente independente, me parece que nesta abordagem ela goza de uma autonomia exagerada. Caímos em uma postura psicologizante da vida social. 

Por outro lado, na segunda abordagem, tende-se a colocar a ênfase na vida social. Nesta forma de abordar a cultura, assumimos que internalizamos as práticas sociais e a própria consciência é fruto da internalização da vida social, o oposto da ideia de que a vida social é predominantemente a expressão da vida interna. Ou seja, enquanto para uma abordagem a vida social nasce da nossa consciência (freudo-marxismo), para outra abordagem a nossa consciência nasce da vida social (materialismo histórico). 

É claro que em algum grau as duas coisas acontecem e se influenciam, a vida social e a vida interna estão relacionadas. Mas as diferentes ênfases levam à diferentes análises, ou analisando a cultura partindo do indivíduo em direção à vida social, ou analisando a consciência individual partindo da cultura. 

Até onde me consta, Freud ataca a ideia de uma consciência totalmente racional, ele introduz a ideia do inconsciente e merece créditos por reconhecer que seres humanos não são soberanos sob seus pensamentos, que existe uma atividade da psique interna que escapa do nosso controle. Mas tenho a sensação que esta abordagem tende a descrever um inconsciente mais autônomo do que deveria ser em relação ao contexto social no qual o indivíduo está inserido.

Já a outra postura tem como objetivo analisar a cultura partindo diretamente da própria vida social mediada por signos. Partindo agora da vida social, avança em direção ao indivíduo, vendo a vida interna deste como fruto da internalização da vida social no qual participa. Não precisamos negar o inconsciente, mas aqui os signos mobilizados no inconsciente tem uma origem externa, o que coloca o inconsciente em uma posição muito mais dependente da vida social que no caso anterior.

Evidentemente, esta separação do marxismo em duas teorias da cultura é uma simplificação, há quase tantas teorias da cultura quanto marxistas. Mas acredito que meu posicionamento também fica ilustrado quando entendemos um pouco a discussão entre Adorno (um dos principais nomes da Escola de Frankfurt), Lukács e Brecht (no qual me inspiro para minha própria teoria da cultura. Lukács aqui não se insere exatamente em nenhuma das duas abordagens que comentei, por um lado ele não adere à tese do freudo-marxismo, mas por outro lado eu não considero que ele também tenha tentado desenvolver uma teoria da cultura a partir do materialismo dialético.

Assim como a Escola de Frankfurt, ele recorre a teorias preexistentes, porém ao invés da psicanálise, ele recorre a Hegel. Podemos ler:

\begin{quote}
Os três [Adorno, Lukács e Brecht] concordavam que a arte podia e devia ser um meio de compreender a realidade histórica. Mas Lukács e Adorno foram além, atribuindo à arte uma capacidade cognitiva intrínseca — mais precisamente, a formas específicas de arte. Ao fazerem isso, foram levados a elaborar versões marxistas de ideologias artísticas preexistentes. Para Lukács, assim como para Aristóteles e a tradição subsequente do pensamento estético realista, a arte era propriamente "a imitação de uma ação"; "ação" deveria ser interpretada à luz do materialismo histórico, mas "imitação" permanecia o propósito inquestionável de toda arte válida. Os ensaios de Adorno não eram tanto uma defesa marxista do modernismo, mas a expressão de um marxismo distintamente modernista: suas posições eram, mutatis mutandis, as da própria ideologia modernista.

As escolhas de Brecht foram produto de uma concepção diferente do papel da política na estética. O "realismo", como ele o definia, era um fim político e ideológico cujos meios formais eram variáveis, de acordo com os ditames do tempo e do lugar… O resultado foi uma estética que, pelo menos em sua concepção, estava muito mais atenta às valências mutáveis da forma do que os estudos de Lukács sobre o passado literário tradicional ou as reivindicações de Adorno em favor da alta vanguarda.

(Aesthetics and Politics)
\end{quote}

Evidentemente este é um recorte para jogar luz na discussão e demonstrar a posição que me interessa. Neste mesmo livro podemos encontrar críticas à exposição teórica do Brecht, como por exemplo, não ter sido tão completa e sistematizada como Lukács. Mas por hora, não é esta minha discussão.

Pode-se dizer que minha rejeição tanto a Adorno quanto a Lukács se dá exatamente pelo que eu entendo de ser uma espécie de reciclagem de ideologias artísticas preexistentes. Entendo que a intenção dos seus teóricos era que estas ideologias fossem ‘absorvidas’ pelo marxismo e re-elaboradas de forma marxista. Porém, eu entendo que terminaram por ganharem apenas uma aparência materialista sem verdadeiramente romper com a raíz idealista. 

Esta tentativa de salvar teorias anteriores ao marxista, que muitas vezes nasce do esforço do pensador de aproveitar as teorias no qual ele se dedicou por anos a estudar antes de se tornar marxista, gera alguns problemas sob a perspectiva do materialismo dialético. 

Por exemplo, é comum encontrarmos em discussões sobre Lukács que a ‘imitação’, ou seja o realismo, permanece um propósito inquestionável. Mas em que sentido “imitação”? Eu entendo que não é um realismo vulgar, uma simples cópia da natureza, porém, entendo que apesar de todos seus esforços, Lukács nunca conseguiu esclarecer satisfatoriamente essa questão. Frequentemente deixando a sensação que este conceito é manobrado de forma bastante flexível para corresponder à uma hierarquização das obras de arte baseada em seus próprios gostos pessoais.

Já Brecht trata a arte de forma mais explicitamente política. Retira dela uma certa ‘especificidade’ da arte, ela é vista de forma mais utilitária. Para mim, isto também corresponde a ver a arte de forma mais materialista, Ela é destituída de sua aura idealista, que muitas vezes se aproxima de até mesmo um certo caráter religioso. Eu já tenho defendido que sequer existe Arte, e o que chamamos de diferentes tipos de artes, são apenas diferentes manifestações da nossa vida social, diferentes tecnologias de mediação da vida social que recebem esse rótulo em determinadas épocas e sociedades. É seguindo este mesmo viés, que eu gosto muito também do trabalho do Boal.

Assim sendo, caminhando em direção à conclusão desta sessão, reconhecendo o peso que o nome do Lukács ocupou na discussão da estética marxista, podemos expandir então para dizer que há três principais abordagens quanto à teoria da cultura sob uma perspectiva marxista.

Enquanto a economia política e filosofia da práxis foi discutida diretamente principalmente por Marx e Engels respectivamente,  a cultura permaneceu sem nenhuma elaboração direta por nenhum dos dois principais nomes do marxismo. Essa ausência levou marxistas posteriores a buscar, fora do próprio marxismo, categorias capazes de explicar a cultura.

Lukács responde a esse problema recorrendo à tradição estética hegeliana. Da teoria da cultura elaborada por  Lukács temos por exemplo a defesa do realismo como forma  privilegiada em uma discussão que não consegue se expandir de forma satisfatória além de alguns gêneros narrativos. Lukács herda de Hegel a ideia de que a obra de arte é tanto capaz quanto obrigada a representar a totalidade social de maneira realista, sendo esta a distinção entre uma grande obra de arte e uma obra passageiro. O resultado é uma estética marxista que permanece em grande medida idealista e subordinada ao gosto pessoal que se esconde por trás de uma hierarquização das obras de arte que tenta ser impessoal.

Já a Escola de Frankfurt responde ao mesmo vazio por outra via: a incorporação da psicanálise freudiana. Nesta abordagem temos como principal característica a visão da cultura como um “sintoma” social, isto é, como expressão da vida interior dos membros dessa sociedade. A crítica ao capitalismo é deslocada então para o plano da subjetividade decorrente de uma psicologização da crítica cultural. Essa prática adotada afasta a teoria da cultura de um viés mais materialista.

E por fim, em contraste com essas duas soluções, autores como Brecht, Bakhtin e Gramsci recusam a construção de uma teoria estética autônoma e deslocam a discussão para campos mais amplos da prática social. De brecht temos uma perspectiva materialista da própria prática artística. A centralidade da recepção e uso social da obra é colocada em foco e a arte é vista como intervenção. Ainda que não elabore um sistema teórico, a arte é vista sob uma perspectiva materialista de forma prática .

Bakhtin é responsável por investigar a linguagem sob uma perspectiva materialista. Vê então ela como uma prática social, importante dizer, que uma prática material. Assim podemos ver a arte e a cultura como um todo como uma arena para conflitos ideológicos a partir do momento que Bakhtin nos leva a pensar em ambos a partir das vida social dos membros da sociedade. 

E por fim, Gramsci é talvez o primeiro grande pensador marxista a realizar esta discussão falando de forma explícita em “cultura”. Sua discussão sobre a guerra de posição (em contrapartida à guerra de assalto) e a importância que o papel que a cultura hegemônica desempenha na manutenção do capitalismo colocou a cultura oficialmente em uma posição de ser vista como um terreno de luta legítimo para os militantes comunistas.

É inspirado principalmente nestes três autores, e posteriormente também por Boal, que eu entendo que a alternativa mais promissora para o desenvolvimento de uma teoria materialista da cultura (um materialismo cultural conforme Raymond Williams coloca) não reside na tentativa de construir uma estética marxista de modo análogo aos desenvolvimentos de filosofias da arte do passado, mas sim em compreender a produção cultural como uma prática histórica situada, atravessada por conflitos sociais, linguísticos e políticos.

Se não existe Arte, então é coerente supormos que não é possível ter uma “Teoria da Arte” (ou qualquer coisa similar). Mas isto não nos impede de que a partir do materialismo dialético sermos  capazes de analisar, entender e transformar as práticas culturais em contextos históricos concretos.

\subsection{Arte}

Eu gostaria que este texto não fosse longo, mas vejo necessidade de abordarmos ainda mais um aspecto referente à cultura. Especificamente a arte. Uma vez que eu mencionei anteriormente sobre a arte ser sempre política, e que quando falamos sobre cultura é impossível não pensarmos em arte, é importante abordarmos rapidamente este assunto. De fato, a relação entre arte e política tem sido um dos grandes motivos pelo qual tenho estudado o marxismo.

Uma das questões que sempre se coloca diante de mim é “o que é arte?” (ou em sua variação “o que é literatura?”). Os autores que eu tenho lido (de Terry Eagleton a Ernst Gombrich, passando por Antoine Compagnon) tendem a concordar que não existe nenhum critério objetivo e universal sobre o que é arte (ou literatura). A resposta sempre depende para quem você pergunta (e frequentemente a classe de quem vai responder). Definir arte é sempre uma escolha, e não algo a ser descoberto na natureza.

\begin{quote}
A definição de um termo como literatura não oferecerá mais que o conjunto das circunstâncias que os usuários de uma língua aceitam empregar este termo…É uma sociedade que, pelo uso que faz dos textos, decide se certos textos são literários fora de seus contextos originais... Uma definição de literatura é sempre uma preferência (um preconceito)

(O Demônio da Teoria)
\end{quote}

Tenho quase certeza que é Eagleton que diz que o que diferencia uma nota fiscal de literatura, é apenas a intenção de quem lê. Podemos pensar na mesma distinção entre literatura e filosofia, se queremos ler diálogos de Platão como filosofia ou literatura, é algo que em última análise, é uma decisão que tomamos, e não algo intrínseco à obra. Mas se não existe um critério universal, se não existe Arte, o que seriam as artes? Ou o conjunto de “coisas” que normalmente chamamos de arte (como cinema, literatura, etc)?

\begin{quote}
UMA COISA QUE realmente não existe é aquilo a que se dá o nome de Arte. Existem somente artistas.

(História da Arte)
\end{quote}

Aqui eu preciso fazer uma proposta que eu reconheço ser altamente contestável e polêmica. Mas eu acredito que é uma proposta que tem pelo menos o mérito de suscitar um debate interessante. Se nos apoiarmos em Marx, o trabalho é um dos grandes diferenciais do homem em relação aos outros animais, e o trabalho é também uma atividade intrinsecamente social desde seus primórdios. Avançando com Bakhtin, obtemos a ideia de que a vida social humana é uma esfera mediada por signos (e não há signos neutros). Então eu proponho pensarmos nas diferentes formas de arte como formas específicas de socialização mediada via signos.

O termo ‘arte’ (ou qualquer nome específico de determinada forma de arte) designa apenas classificações históricas e sociais, usadas para classificar certas formas de comunicação em determinados contextos. Eu estou assumindo implicitamente que quando um artista cria uma obra de arte, sempre há a proposta de uma comunicação, ainda que o receptor da sua comunicação seja a si próprio. Este talvez seja o ponto mais polêmico, mas não acredito ser totalmente ignorante.

Me parece razoável supor que sempre um artista cria algo através do seu trabalho, e aqui estamos considerando arte apenas aquilo que é produto do trabalho humano, ele visa comunicar algo. Minha concepção de obra de arte aqui, é então, o produto de um trabalho humano, no qual o artista realiza uma operação sobre natureza visando transformá-la e expressar a si mesmo através do próprio trabalho. Lembrando que tanto o consumidor da sua obra de arte pode ser si próprio, quanto o artista não precisa estar plenamente ciente de o que deseja comunicar. Um artista pode acreditar estar escrevendo uma história neutra, ainda que não seja possível existir uma história neutra.

O trabalho artístico, sob esta definição, é sempre teleológico, isto é, orientado por um fim previamente posto, no qual o ser humano antecipa idealmente a obra e organiza sua prática de modo intencional, ainda que o resultado final não seja um simples reflexo dessa intenção inicial. Assim, a diferença entre a declamação de uma poesia e um diálogo cotidiano não diz respeito a uma diferença fundamental nas atividades, são ambas atividades sociais mediadas por signos onde uma mensagem é enviada para alguém, a diferença está em como os signos que são mobilizados são vistos e classificados pela sociedade.

Uma mesma carta, em um contexto pode ser apenas uma carta, e em outro contexto, pode ser literatura. Não sou capaz de conceber nenhuma produção artística que esteja ausente de mensagem, até mesmo porque todo signo possui um significado, e consequentemente transmite alguma mensagem, e por sua vez, é impossível consumirmos uma obra de arte sem que esta interação seja mediada por signos.

Sendo assim, a obra de arte adquire uma importância política intrínseca. Pois ela é, na sua essência, a transmissão de alguma mensagem que não é neutra. De fato, esta é uma mensagem que tende a ser construída de forma mais elaborada para ser mais interessante e persuasiva que um diálogo cotidiano. Muitas vezes a mensagem que é transmitida sequer é algo óbvio como uma típica comunicação no cotidiano exige, o que pode conferir uma falsa aparência de neutralidade à obra de arte. Mas é aqui também que podemos perceber a dupla função da arte de divertir e ensinar, que se não estou errado, já foi discutido tanto por Aristóteles como Brecht.

Não livre de polêmicas, mas podemos concluir então, que se não existe Arte, existem artes. E estas artes são diferentes formas de interação mediadas via signos, isto é, diferentes formas de transmitir uma mensagem que em dado local e época recebe o título de arte. Este título de arte tende a estar relacionado à um planejamento prévio que tipicamente confere uma maior complexidade nos signos mobilizados para transmitir a mensagem, e também possui uma dupla função de divertir e ensinar, características que muitas vezes podem dar a falsa sensação de que há arte sem mensagem, ou seja, sem política.

Podemos entender um pouco mais sobre como se manifesta esse caráter pedagógico da arte. Sobre isto, a principal leitura que eu tenho em mente é “ Teatro do Oprimido e outras poéticas políticas” de Augusto Boal. No campo artístico, assim como Marx “inverteu” Hegel, Brecht faz uma operação análoga na área da poética.

\begin{quote}
…para Hegel o personagem é inteiramente livre…; para Brecht (e para Marx) o personagem é objeto de forças sociais… Sobretudo, Hegel insiste em um ponto fundamental que marcará sua profunda diferença com a poética marxista de Brecht; “a ação não parece nascer de circunstâncias exteriores mas sim da vontade interior e dos caracteres dos personagens’’,

(Augusto Boal em Teatro do Oprimido e outras poéticas políticas)
\end{quote}

É interessante notar como as visões destes artistas de como a arte deveria ser divergem baseado na sua concepção idealista ou materialista do mundo. Para Hegel, a única restrição que deveria existir nos personagens é interna, por isso ele vai defender personagens que estejam livres das preocupações materiais, personagens que não tenham restrições econômicas, que podem mais facilmente tornar sua vontade em realidade.

Já em Brecht, existe uma preocupação em mostrar a realidade social e as restrições que esta realidade impõe ao personagem. Ele não é livre, sua vontade interior não é soberana, este personagem ocupa uma posição central na poética para Brecht. Ainda que alguém em posição de poder seja representado no teatro, ele se encontra com determinadas restrições devido ao cargo que ocupa. Boal dá de exemplo uma peça onde o Papa como pessoa demonstra boa vontade com Galileu, mas esta mesma pessoa, como alguém na posição de Papa o condena. Para Brecht é necessário mostrar essa restrição social que se impõe à vontade interior dos personagens.

Ou seja, enquanto Hegel busca uma poética que busca apresentar uma história no qual as ações nascem das vontades interiores, Brecht propõe uma poética onde as restrições da vida social estão explícitas e consistem precisamente na origem de muitas das ações.

Estas diferentes posturas surgem pois seus teóricos possuem diferentes objetivos políticos com a arte. Para Hegel, ao final do espetáculo o equilíbrio deveria ser restabelecido, e ela deveria estimular uma coerência entre os valores defendidos pela sociedade atual e a plateia, assim como Platão, purificando os pensamentos “errados” da plateia através do consumo da obra.

Brecht, de certa forma,  também visa atingir uma conformidade entre a sociedade e os sentimentos da plateia, mas não através do expurgo dos sentimentos da plateia que vão contra ao sistema socioeconômico estabelecido, mas pelo contrário, motivando a plateia a alterar a sociedade, para que então esta nova sociedade esteja de acordo com seus sentimentos. Neste sentido Hegel promove a restauração do equilíbrio (conservador) através da modificação dos sentimentos da plateia, e Brecht a promove a transformação (revolucionário), através da modificação da Sociedade.

Para isso, Hegel propõe histórias em que as ações dos personagens têm causa interna e se externalizam no mundo (idealismo); já Brecht propõe histórias nas quais as ações possuem causas sociais que moldam o homem e o transformam contra suas vontades internas (materialismo).

\begin{quote}
Os primeiros desejam uma quieta sonolência ao final do espetáculo: Brecht deseja que o espetáculo teatral seja o início da ação, o equilíbrio deve ser buscado transformando-se a sociedade e não purgando o indivíduo dos seus justos reclamos e de suas necessidades.

(Augusto Boal em Teatro do Oprimido e outras poéticas políticas)
\end{quote}

É por isso também, que o público alvo do Brecht, não é uma elite aristocrática, mas o povo.

\begin{quote}
Nossas representações da vida social devem estar destinadas aos técnicos fluviais, aos cuidadores das árvores, aos construtores de veículos e aos revolucionários. Nós os convidamos para que venham aos nossos teatros e lhes pedimos que não se esqueçam de suas ocupações (alegres ocupações), para que nos seja possível entregar o mundo e nossa visão do mundo às suas mentes e aos seus corações, para que eles modifiquem O MUNDO AO SEU critério

(Augusto Boal em Teatro do Oprimido e outras poéticas políticas)
\end{quote}

Acho interessante observar que a arte sob este uso político, possui um teor de normatividade. Porém essa normatividade é pragmática, operacional. Ela cumpre um programa político. A arte é feita de determinada forma e com determinado conteúdo, para que determinados fins políticos sejam alcançados.

Podemos entender melhor o uso político da arte através de um exemplo, acompanhando a discussão sobre o teatro grego desenvolvida também por Boal no mesmo livro. Primeiro precisamos destacar que a tragédia grega era patrocinada pelo estado diretamente ou indiretamente através dos mecenas, então não é ao acaso que ela sirva de ferramenta política para a manutenção do poder pela classe dominante.

Para entender como isso funciona, precisamos entender que a arte era concebida por Aristóteles como tendo um papel de correção. Para Aristóteles, as coisas tenderiam sempre à perfeição, mas não significa que conseguiriam alcançar esta perfeição. A arte e a ciência surgem então para recolocar o desenvolvimento humano “nos trilhos”.

\begin{quote}
… para isso serve a arte e serve a ciência: para, “recriando o princípio criador " das coisas criadas, corrigir a natureza naquilo em que haja fracassado.

(Augusto Boal em Teatro do Oprimido e outras poéticas políticas).
\end{quote}

E aqui, a gente precisa ter em mente que o bem supremo, para Aristóteles, é a busca pela felicidade. E que a felicidade não é outra coisa, senão obedecer às leis.

\begin{quote}
“A Tragédia imita as ações da alma racional do homem, suas paixões tomadas hábitos, em busca da felicidade, que consiste no comportamento virtuoso, que é aquele que se afasta dos extremos possíveis em cada situação dada concreta» cujo bem supremo é a Justiça, cuja expressão máxima é a Constituição!”… Em última instância, a felicidade consiste em obedecer às leis!

(Augusto Boal em Teatro do Oprimido e outras poéticas políticas)
\end{quote}

Temos então uma produção teatral que visa corrigir o desenvolvimento daqueles que não se encontram inclinados a seguir a constituição, ou seja, que desafiam as normas sociais e leis correntes.

\begin{quote}
O homem, como parte da natureza, tem certos fins em vista: a saúde, a vida gregária no Estado, a felicidade, a virtude, a justiça, etc. Quando falha na consecução desses objetivos, intervém a arte da Tragédia. Esta correção das ações do homem, do cidadão, chama-se “catarse”.

(Augusto Boal em Teatro do Oprimido e outras poéticas políticas)
\end{quote}

E isso é feito através da “catarse”, que é exatamente esse processo de destruir uma “impureza” no consumidor do teatro. Ou seja, o teatro trágico descrito por Aristóteles é intrinsecamente repressivo, pois ela visa corrigir o homem que não segue as leis como ideal de busca da felicidade.

\begin{quote}
A impureza que o processo trágico vai destruir é pois algo que atenta contra as leis.

(Augusto Boal em Teatro do Oprimido e outras poéticas políticas)
\end{quote}

Como se busca alcançar esse efeito é bastante interessante e educativo. Apesar deste modelo clássico ser antigo, ele continua em uso e também deu origem a algumas variações modernas, mas todos estes formatos compartilham desse objetivo comum de “destruir a impureza”.

O formato clássico de Aristóteles funciona então, resumidamente do seguinte modo: primeiro somos apresentados a um personagem que é dotado de muitas qualidades (conforme as ideias dominantes da sociedade) e apenas um defeito, preferencialmente um defeito que é o responsável por levar o protagonista na boa situação que ele se encontra no começo da história. Nesta fase, somos levados a sentir empatia pelo personagem, a ideia é que alguém que possui tão poucas falhas, não merece um fim tão trágico como o para o qual a história se encaminha.

Em um segundo momento, o sentimento que deve surgir na plateia é o terror. A plateia deve agora, perceber que a falha no caráter do protagonista é o que o levará a sua ruína, e não só isso, mas que é uma falha que ele próprio também é capaz de possuir. Sentimos então o terror por reconhecermos esta falha em nós mesmos.

E por fim, a única coisa que resta ao nosso protagonista é o desfecho trágico. Afinal, ele deve ser punido pela sua falha, pelo vício no qual carrega. E nós então, como plateia, somos educados a combater este vício em nós mesmos. Temos então, um jogo de empatia e terror mediado por signos que visa combater determinados sentimentos na plateia. Evidentemente, como um teatro patrocinado pelo estado, quem define o sentimento a ser combatido, é o estado, ou seja, a classe dominante.

\begin{quote}
Podemos agora concluir que, quando o homem falha nas suas ações, no seu comportamento virtuoso em busca da felicidade, através da virtude máxima que é a obediência às leis, a arte da Tragédia intervém para corrigir essa falha… O herói trágico surge quando o Estado começa a utilizar o teatro para fins políticos de coerção do povo.

(Augusto Boal em Teatro do Oprimido e outras poéticas políticas)
\end{quote}

Podemos notar que a empatia cumpre um papel central na dinâmica do teatro trágico de Aristóteles. É por isso, que Boal considera a empatia como uma “arma terrível”:

\begin{quote}
A empatia tem que ser entendida como a arma terrível que realmente é. A empatia é a arma mais perigosa de todo o arsenal do teatro de artes afins (cinema e TV).

(Augusto Boal em Teatro do Oprimido e outras poéticas políticas)
\end{quote}

Gostaria de recomendar o trecho inteiro de “Empatia ou osmosis?”, porém para poupar o leitor, tentarei resumir a ideia. O mecanismo no qual a arte opera com a empatia, consiste em justapor duas pessoas (uma fictícia e uma real) e fazer com que uma dessas pessoas ofereça a outra seu poder de decisão. Isto é, a pessoa real abdica em favor do personagem fictício seu poder de decisão.

Esta justaposição de dois universos (real e fictício) faz com que o espectador vivencie a ficção e incorpore elementos da ficção. Porém, o espectador que é um homem real e vive, assume então, como realidade o que se apresenta como ficção na obra de arte. É isto que Boal chama de: \textit{osmosis estética}.

\begin{quote}
Mas existe aqui algo monstruoso: o homem, quando ele­ge, elege em uma situação real, vital, elege em sua própria vida; o personagem quando elege, ( e por isso, quando induz o homem a eleger), elege em uma situação fictícia, irreal, des­provida de toda densidade de fatos, matizes e complicações que a vida oferece. Isto faz com que o homem, real, eleja segundo situações e critérios irreais.

(Augusto Boal em Teatro do Oprimido e outras poéticas políticas)
\end{quote}

Como consequência disto, o homem real em uma situação real toma decisões na sua própria vida influenciado por um personagem que elege (e com isso induz o homem real a fazer a mesma escolha) em uma situação fictícia, irreal e desprovida de toda densidade da vida real. Assim, um homem real faz suas decisões segundo situações e critérios irreais.

Concluindo então, podemos dizer que:

\begin{itemize}
    \item Não há critério universal/objetivo para definir arte;
    \item Toda vida social é mediada por signos;
    \item Não há signos neutros;
    \item A arte é trabalho humano intencional que se insere na vida social.
\end{itemize}
Portanto:

\begin{itemize}
    \item Arte é uma forma específica de comunicação mediada por signos;
    \item Seu status como “arte” é determinado social/historicamente;
    \item Sendo comunicação, é sempre política (transmite mensagens que não são neutras).
\end{itemize}
Ou seja, “artes” são diferentes tecnologias de mediação simbólica que reforçam ou combatem determinados valores. Portanto, o que chamamos de “arte” não é uma categoria metafísica, mas um evento comunicativo socialmente consagrado. As diversas artes (do teatro à literatura) são linguagens materiais que as sociedades desenvolvem para dizer, com especial complexidade e intensidade, aquilo que julgam merecer ser reforçado ou combatido para além do discurso comum. São estratégias de significação complexas, mas no princípio e no fim, é sempre gente tentando falar com gente, apenas com melhores ferramentas.

\textit{Conceber a existência da Arte como algo distinto da comunicação mediada por signos, é despolitizar a arte.}

\subsection{Bakhtin}

Grande parte das minhas ideias sobre cultura e arte são fortemente baseadas nos trabalhos de Bakhtin. Sendo assim acho pertinente que eu realize pelo menos uma pequena exposição sobre as ideias deste pensador.

Bakhtin aprofunda a discussão de que é a matéria que cria a consciência, ele discute como isso acontece partindo da ideia de que a ideologia é um fenômeno social, criticando algumas abordagens que tratam ideologia como algo individual, algo que existe na mente de um indivíduo como um fenômeno psicológico. Podemos dizer que ele é contra um psicologismo da ideologia.

\begin{quote}
Todo signo é ideológico; a ideologia é um reflexo das estruturas sociais; assim, toda modificação da ideologia encadeia uma modificação da língua.

(Mikhail Bakhtin em marxismo e Filosofia da Linguagem)
\end{quote}

Precisamos entender o que é ideologia para Bakhtin:

\begin{quote}
Para o Círculo de Bakhtin, a palavra ideologia tem significado diferente daquele usado por parte da tradição marxista… a palavra ideologia é usada em geral para designar o universo dos produtos do “espírito humano”, aquilo que algumas vezes é chamado por outros autores de cultura imaterial ou produção espiritual… Assim, ideologia para o Círculo de Bakhtin, abrange um grande universo: a arte, a filosofia, a ciência, religião,ética, política. Todo produto ideológico parte de uma realidade (natural ou social), possui um significado e remete a algo que lhe é exterior, ou seja, é um signo.

(Ideologia e dialogismo: Mikhail Bakhtin, história e conceitos que cabem na sala de aula.)
\end{quote}

Bakhtin defende então que a ideologia existe de forma exterior a nós, ela está nos signos que utilizamos na vida social. Com os importantes detalhes de que esta significação do signo está sempre inserido em uma luta ideológica e que não existe objeto físico neutro, isto é, ele já entra na vida social sendo mediado por algum signo.

Uma foice e um martelo, o objeto físico seria só isso, um objeto material, em um mundo sem humanos. Mas ao interagirmos com os objetos e entre nós, construímos signos para mediar nossa interação. Um exemplo explícito consiste no fato de que o signo da foice e o martelo por um lado está relacionado aos respectivos objetos físicos, mas por outro, atualmente é socialmente entendido como simbolizando as ideias comunistas. Os signos que mediam nossa vida social nos fazem ir além do mero objeto físico em si, ele reflete outra coisa, transfere nosso pensamento para outro lugar. Você olha uma foice e um martelo, o objeto físico, mas o que ‘vê’ é o comunismo.

Bakhtin defende que nossa própria consciência emerge da vida social mediada por signos. Devemos lembrar também que o próprio signo não tem uma existência meramente abstrata, ele precisa ter uma existência material, ele precisa existir como algo que possa ser percebido e compartilhado, mesmo que seja um som ou um registro escrito, já que não temos telepatia. Isso faz com o meio material que estamos inserido se torne mais determinante para nossa consciência do que normalmente estamos acostumados a  pensar.

Esta posição combate diretamente à ideia de que precisamos apenas de “boas ideias” para convencer as pessoas. Também relaciona a forma e o conteúdo, de maneira indissociável. A forma com que algo se apresenta não é independente do seu conteúdo, toda interação do qual participamos é mediada por signos, e cada signo possui um significante em disputa. Quando emitimos uma mensagem, a forma não é independente do conteúdo para definir a mensagem que queremos propagar. Por exemplo, podemos dizer “uma mesma coisa” com diferentes palavras, ou até mesmo, uma mesma palavra pode ser pronunciada de diferentes formas e expressar diferentes mensagens.

Me parece que com isso caminhamos para ter condições de explicar como que o meio material onde estamos inseridos opera para definir as ideias dominantes de uma época de uma forma mais profunda do que normalmente é considerado. A classe dominante domina também os meios materiais por onde os signos se produzem e se transmitem. Com isso somos obrigados a politizar a vida social por completo. Porque agora a política e propagação ou combate de ideologias vão muito além de materiais de cunho explicitamente político, mas é algo que se faz presente em toda vida social.

\begin{quote}
… a “ideologia do cotidiano”, que se exprime na vida corrente, é o cadinho onde se formam e se renovam as ideologias constituídas… A ideologia do cotidiano constitui o domínio da palavra interior e exterior desordenada e não fixada num sistema, que acompanha cada um dos nossos atos ou gestos e cada um dos nossos estados de consciência. Considerando a natureza sociológica da estrutura da expressão e da atividade mental, podemos dizer que a ideologia do cotidiano corresponde, no essencial, àquilo que se designa, na literatura marxista, sob o nome de “psicologia social”.

(Mikhail Bakhtin em marxismo e Filosofia da Linguagem)
\end{quote}

Bakhtin vai defender também que a palavra é um instrumento da consciência e é devido a esse papel que a “\textit{… palavra funciona como elemento essencial que acompanha toda criação ideológica, seja ela qual for}.” Juntamos a isso o fato de que a consciência tem o poder de abordar cada signo, independente da sua natureza, de modo verbalmente, então a palavra se torna “\textit{…o objeto fundamental do estudo das ideologias}”. Por isso então:

\begin{quote}
A única maneira de fazer com que o método sociológico marxista dê conta de todas as profundidades e de todas as sutilezas das estruturas ideológicas “imanentes” consiste em partir da filosofia da linguagem concebida como filosofia do signo ideológico.

(Mikhail Bakhtin em marxismo e Filosofia da Linguagem)
\end{quote}

Tenho que uma das coisas mais geniais do Bakhtin é “diluir” a fronteira entre a atividade mental interior e a comunicação como atividade exterior. Ambos são tratados como um mesmo processo, uma mesma atividade, interpretando a atividade mental interior como uma “conversa interna”. A diferença entre ambos é mais de grau que de qualidade, a atividade mental tende a ser mais confusa e fragmentada, ela ganha clareza à medida que passa a se transformar em uma comunicação externa, quando organizamos nossas ideias como uma auto explicação, ela já se aproxima desta dimensão de exteriorização da atividade mental. Isso também parece corroborar minha própria experiência pessoal de escrever para organizar os pensamentos.

Essa forma de entender a consciência é bastante radical quando se une ao fato que os signos utilizados na comunicação tem uma origem social. Isso imprime uma natureza social na nossa própria atividade mental interior, isso nos proíbe que estejamos, a qualquer momento, totalmente isolados das influências culturais da sociedade. E a natureza social dos signos, também vincula eles a uma realidade material no qual essa realidade social é construída.

Nossa existência então, como animais sociais se dá tanto em uma dimensão material onde a vida material se desenrola, quanto em uma dimensão simbólica (ou como Bakhtin chama: dimensão ideológica) onde a vida social acontece, mas ambas dimensões relacionam-se entre si de forma dialética. Dessa forma conectamos a questão da infraestrutura e superestrutura e chegamos à questão sobre como as ideias dominantes de uma época são as ideias da classe dominante. Reduzimos pouco a pouco a abstração na discussão sobre a cultura.

Claro que ainda há sempre algo a mais para ser explicado, mas aqui partimos de duas categorias (no sentido mais vulgar do termo) mais gerais como infra e superestrutura e agora temos uma questão mais específica e objetiva de responder, a abstração foi reduzida. Queremos entender agora como os signos pertencentes a dimensão ideológica afetam e são afetados pela dimensão material, uma vez que já fomos convencidos da importância dos signos em nossa atividade social.

Podemos pensar em um exemplo concreto. Lembrando da posição central que a palavra ocupa, é fácil perceber como o surgimento de novos signos dependem das condições materiais no qual os indivíduos se encontram, não teremos a palavra “telefone” se não existir algo ser chamado de telefone. Não é ao acaso que a bíblia nunca nos advertiu sobre os riscos do tik tok.

E à medida que as palavras são os mesmos signos que mobilizamos tanto para comunicação como para a atividade mental, assim são capazes de influenciar e moldar nossa concepção de mundo. Afetando não apenas como o percebemos, mas também como nos posicionamos diante deste mundo. Estamos em uma situação melhor de entender agora como os signos influenciam o modo no qual agimos no mundo, o que em última instância, nos leva a alterar o próprio mundo material no qual estamos inseridos e serve de base para a dimensão ideológica.

Você precisa que o modo de produção capitalista exista para que existam signos relacionados ao capitalismo, sem esta condição satisfeita não podem existir signos que carregam ideias anticapitalistas. Mas a partir do momento que o capitalismo existe, estes signos sob determinadas condições, podem desempenhar um papel relevante para a formação de uma concepção de mundo anticapitalista que induza as pessoas a agir pela superação do próprio capitalismo no qual estes mesmos signos se originaram.

Entre a infra e a superestrutura, entre a dimensão material e a ideológica, temos a palavra como um signo de posição especial que permeia tanto nossa comunicação externa com o mundo quanto nossa atividade mental em nosso próprio mundo interior. Tornando assim impossível que nossa concepção de mundo permaneça inalterada diante da realidade material e social mediada que se desenrola diante de nós. E é a partir disso que entendo que podemos refletir sobre a importância da arte e da cultura de forma verdadeira materialista.

Podemos dizer que:

\begin{itemize}
    \item não existe pensamento sem signos;
    \item não existe vida social sem signos;
    \item não existe signo que não seja social;
    \item logo, a consciência individual é sempre atravessada pela vida social.
\end{itemize}
Ou seja, nossa atividade externa (comunicação) e interna (atividade mental) e consequentemente nossa própria apreensão da realidade não se dá apenas por palavras, mas de forma mais geral, por signos. E não há elementos na nossa organização social e comunicação desprovida de signos.

Isso se conecta à uma antiga reflexão que tenho feito sobre o método de aprendizado (passivo) estético através dos sentidos, em comparação com um método de aprendizado (ativo) racional. O segundo seria o tradicional sentar e refletir sobre algo, o primeiro é aquela apreensão que temos diretamente pelos sentidos ao entrarmos em contato com o mundo. 

Por exemplo, podemos ter um aprendizado racional ao estudarmos sobre a história dos EUA de forma intencional. Mas podemos ter também um aprendizado estético não intencional que nos induz a construir uma opinião sobre os EUA baseado em filmes  e séries. O elemento em comum que une os dois processos de aprendizado é o uso de signos. Pode-se argumentar que o método de aprendizado racional faz uso da palavra como principal signo, e o estético tende a trabalhar com um conjunto maior de signos, porém ambos aprendizados são mediados por signos.

Assim sendo, me parece natural concluir que uma teoria estética integra de forma mais ampla, uma teoria do conhecimento.

\begin{quote}
Como toda concepção do mundo, o marxismo possui a sua própria teoria estética, que integra, de modo geral, a sua teoria do conhecimento.

(Leandro Konder em Os Marxistas e a Arte)
\end{quote}

\subsection{Economia}

A terceira componente então que vai interligar ambos os temas é a economia marxista. Afinal, é preciso investigar os aspectos econômicos da sociedade capitalista para saber como superá-lo. Assim como é natural supor que é preciso possuir uma teoria adequada para gerenciar a produção e distribuição de bens no comunismo.

Podemos notar como na verdade o materialismo dialético, o socialismo científico e a economia marxista não podem ser entendidos como disciplinas isoladas. É preciso mobilizar todos os conhecimentos para uma efetiva transformação da realidade. Por exemplo, quanto a primeira fase do comunismo, Marx nos fornece alguns detalhes da organização socioeconômica que ele visava para esta etapa histórica:

\begin{quote}
O que o produtor deu à sociedade constitui sua cota individual de trabalho. Assim, por exemplo, a jornada social de trabalho compõe-se da soma das horas de trabalho Individual; o tempo Individual de trabalho de cada produtor em separado é a parte da jornada social de trabalho com que ele contribui, é sua participação nela. A sociedade entrega-lhe um bônus consignando que prestou tal ou qual quantidade de trabalho … e com este bônus ele retira dos depósitos sociais de meios de consumo a parte equivalente à quantidade de trabalho que prestou. A mesma quantidade de trabalho que deu à sociedade sob uma forma, recebe-a desta sob uma outra forma diferente…

...Variaram a forma e o conteúdo, porque sob as novas condições ninguém pode dar senão seu trabalho, e porque, de outra parte, agora nada pode passar a ser propriedade do indivíduo, fora dos meios individuais de consumo.

(Marx em Crítica ao programa de Gotha)
\end{quote}

Agora para entendermos melhor como se daria a construção desta sociedade, sobre como este bônus deve operar e como ele difere do dinheiro, acabamos tendo que nos dedicarmos ao estudo da economia marxista.

\begin{quote}
Quando Marx afirma que os certificados de trabalho não valem mais do que um ingresso de teatro, podemos inferir certas implicações:
\begin{enumerate}
\item Os certificados não circulam; eles só podem ser trocados diretamente por bens de consumo.

\item  … Somente a pessoa que realizou o trabalho poderia usá-los.

\item Eles seriam cancelados após um único uso… A loja, como uma organização comunitária, não precisa comprar mercadorias, apenas as recebe, portanto, seu único interesse nos comprovantes de trabalho é para fins de registro.

\item Eles não serviriam como reserva de valor. Eles poderiam ter uma data de validade…
\end{enumerate}
(Paul Cockshott em Towards a New Socialism)
\end{quote}

Evidentemente, a forma com que esta ciência será abordada, irá depender de como entendemos o materialismo dialético, e isto impactará por sua vez no desenvolvimento do socialismo científico. Por exemplo, a nossa interpretação de o que é a teoria do valor-trabalho apresentada por Marx, depende de como entendemos o materialismo dialético.

Ou seja, apesar de tradicionalmente discutirmos que o marxismo tem três componentes, estes três componentes não podem ser entendidos de maneira isolada. Faço esta discussão com o propósito de contextualizar a defesa de que quando vamos estudar economia política, temos hoje ferramentas disponíveis mais interessantes e úteis do que as que Marx dispôs em sua época, de forma que é natural adotarmos novas abordagens. Em outras palavras, reafirmo a ideia expressa anteriormente de que a dialética hegeliana não é mais a principal ferramenta para realizar esta investigação.

Sendo assim, defendo que uma economia marxista que se afasta de Hegel ainda pode ser marxista, desde que mantenha os pressupostos do materialismo dialético expostos anteriormente. Seguindo esta linha de raciocínio, só para citar alguns exemplos, Farjoun e Machover vão desenvolver uma formulação probabilística da economia política e da teoria do valor ao longo de suas principais obras: “Laws of Chaos: A Probabilistic Approach to Political Economy” e “How Labor Powers the Global Economy”.

Paul Cockshott vai escrever “Classical Econophysics”. Esta é uma contribuição para a continuação da tradição marxista que se inicia com o trabalho de Marx, mas que não abre mão do uso de ferramentas científicas modernas.

\begin{quote}
Após uma análise das trocas básicas de mercadorias, passamos a examinar o dinheiro e o que Marx denominou ‘a forma do valor’. A obra de Marx permanece interessante, pois tentou aplicar uma análise formal à troca e ao dinheiro, utilizando as ferramentas então disponíveis: a lógica hegeliana. Tentamos reanalisar o aspecto formal do valor de troca utilizando uma tríade de conceitos extraídos de teorias modernas de sistemas formais:
\begin{enumerate}
\item o conceito de espaço métrico,

\item a noção de simetria e

\item a ideia de assinatura.
\end{enumerate}
(Paul Cockshott em Classical Econophysics)
\end{quote}

Ian Wright vai também contribuir com abordagens probabilísticas e estatísticas que investigam temas como a arquitetura social do capitalismo (“The Social Architecture of Capitalism”) e a lei do valor(“The Emergence of the Law of Value in a Dynamic Simple Commodity Economy”), etc.

Este não é um espaço para me aprofundar em nenhum tema, apenas uma exposição geral sobre minha posição em diferentes temas, e principalmente de que não só podemos, mas que devemos nos apropriar das ferramentas modernas mais eficientes que a ciência produziu. De certa forma, eu entendo por economia marxista simplesmente uma abordagem da economia que respeita os pressupostos do materialismo dialético. De modo geral podemos falar em:

\begin{itemize}
    \item Economia Marxista: quando realizamos uma investigação econômica coerente com o materialismo dialético e nasce da crítica à economia política inglesa.
    \item Socialismo científico: quando realizamos uma investigação (e prática) acerca da superação do capitalismo coerente com o materialismo dialético e nasce da crítica ao socialismo utópico francês.
    \item Materialismo Histórico: quando realizamos uma investigação acerca da realidade histórica e social da humanidade em particular, e evidentemente, coerente com o materialismo dialético.
\end{itemize}

\subsection{Inovação}

O argumento de que o socialismo seria menos produtivo que o capitalismo é frequentemente mobilizado por críticos do socialismo. Embora eu não veja razões para que um modo de produção socialista seja intrinsecamente menos produtivo, afinal, ele pode se apropriar de todas as tecnologias existentes e utilizá-las de forma ainda mais racional, livre das restrições impostas pela busca do lucro, há um problema real que merece atenção: o ritmo de desenvolvimento das forças produtivas.

Diversos autores marxistas reconhecem que um dos diferenciais históricos do capitalismo em relação aos modos de produção anteriores foi o incentivo permanente à maximização da produtividade. Trata-se de um mérito histórico específico: nenhum modo de produção anterior estimulou de forma tão contínua o desenvolvimento das forças produtivas. Esse incentivo, no entanto, não é voluntário, mas coercitivo. O capitalista é compelido a priorizar o lucro sob pena de falência, e essa lógica o obriga a investir continuamente em produtividade.

Nos modos de produção pré-capitalistas, a exploração operava predominantemente por vias extra econômicas. O aumento da exploração ocorria por meios como a ampliação de tributos ou a conquista de novas terras, isto é, a intensificação da coerção política direta. Mesmo quando havia pressão por maior produção, essa pressão tendia a se estabilizar após certo ponto. No capitalismo, ao contrário, a coerção econômica é permanente e recai justamente sobre aqueles que controlam os meios de produção, forçando o reinvestimento contínuo.

Além disso, a ameaça constante do desemprego cria também para o trabalhador um incentivo coercitivo à elevação da produtividade. Diferentemente do servo, que não podia ser facilmente expulso da terra, o trabalhador assalariado pode ser substituído, o que intensifica a disciplina do trabalho.

O problema que se coloca é que, no socialismo, essa coerção econômica capitalista é suprimida. Isso não implica menor produtividade em termos absolutos, mas pode implicar um ritmo menor de inovação tecnológica, caso não existam mecanismos substitutivos capazes de sustentar um desenvolvimento acelerado das forças produtivas.

A experiência histórica parece confirmar essa dificuldade. Nem a URSS nem a China conseguiram eliminar a coerção. A URSS recorreu amplamente a formas de coerção extra econômica; a China, por sua vez, reintegrou com eficiência mecanismos de coerção econômica baseados no mercado, mantendo a lei do valor em funcionamento. Isso permite caracterizar a China menos como um modo de produção socialista e mais como uma forma específica de ditadura do proletariado operando sobre um modo de produção capitalista..

Outras alternativas, como a generalização de cooperativas, oficialmente defendida pela Coreia Popular, também não eliminam a coerção econômica, já que a lei do valor persiste nas relações entre unidades produtivas. Além disso, tais experiências não demonstraram capacidade de sustentar um ritmo de desenvolvimento comparável ao da China, especialmente no plano militar.

Esse ponto é decisivo. Historicamente, modos de produção menos coercitivos, como o comunismo primitivo das sociedades originárias, foram incapazes de competir militarmente com sociedades capitalistas, apesar de sua organização social menos opressiva. O capitalismo demonstrou superioridade não apenas econômica, mas militar, e derrotou outros modos de viver por essa via.

Diante disso, a questão central não é se o socialismo pode existir sem coerção, mas que tipo de coerção é aceitável e funcional. Uma hipótese plausível é que a resposta esteja no campo cultural. Assim como o capitalismo se sustenta em uma cultura hegemônica que naturaliza a coerção econômica, um socialismo viável exigiria uma cultura hegemônica própria, capaz de incentivar o desenvolvimento das forças produtivas e legitimar coerções secundárias — como a vinculação entre contribuição produtiva e acesso a bens, nos termos do bônus de trabalho mencionado por Marx.

Nesse sentido, a coerção não desaparece: ela é transformada. Primeiro em coerção social, mediada pela cultura hegemônica; depois em coerção institucional, legitimada por essa mesma cultura. 

\subsection{Valor}

Este é sem dúvida um dos capítulos mais importantes e que talvez eu deveria ter escrito enquanto estava estudando economia, isto me permitiria trazer a questão de forma melhor elaborada e com mais citações. De toda forma, quero hoje pelo menos lançar a discussão sobre o assunto, e que se for o caso, que seja aperfeiçoada no futuro.

Eu quero começar discutindo um pouco sobre o que é a lei do valor, pois eu vejo ela como o núcleo central do capitalismo, e a própria superação do capitalismo passa pela superação da lei do valor. A forma com que eu avalio se um governo é socialista ou não, mais do que avaliando a situação no qual se encontra, é avaliando como ele muda no tempo e com que intenção. Se este governo caminha no sentido de buscar a superação da lei do valor ou não.

Evidentemente para isso ser possível, precisamos começar apresentando o que é a lei do valor. Primeiro, eu gostaria de discutir em que situações o valor corresponde ao preço.

Para que os preços a que as mercadorias são trocadas correspondam aproximadamente aos seus valores, nada mais é necessário do que 1) que a troca das várias mercadorias deixe de ser puramente acidental ou apenas ocasional; 2) no que diz respeito à troca direta de mercadorias, que essas mercadorias sejam produzidas em ambos os lados em quantidades aproximadamente suficientes para atender às necessidades mútuas, algo aprendido com a experiência mútua no comércio e, portanto, uma consequência natural do comércio contínuo; e 3) no que diz respeito à venda, que não haja nenhum monopólio natural ou artificial que permita a qualquer uma das partes contratantes vender mercadorias acima de seu valor ou obrigá-las a vender abaixo do valor. Por monopólio acidental entendemos um monopólio que um comprador ou vendedor adquire por meio de uma situação acidental de oferta ou demanda.

\begin{quote}
A suposição de que as mercadorias das várias esferas de produção são vendidas pelo seu valor implica, naturalmente, que o seu valor é o centro de gravidade em torno do qual os seus preços flutuam, e que as suas subidas e descidas contínuas tendem a se equalizar.

(Marx em O Capital Vol 3)
\end{quote}

Esta forma do Marx abordar a questão, conduz diretamente a uma interpretação da lei do valor.  A lei do valor nos diz que, em equilíbrio, o preço médio de uma mercadoria, a expressão monetária desta lei, tende a se correlacionar com seu valor, expresso quantitativamente através do número médio de horas necessárias para sua produção. Ou seja, se a sociedade produzir uma determinada mercadoria em quantidade suficiente para atender à demanda, o preço médio dessa mercadoria estará relacionado ao tempo médio necessário para sua produção. 

Comparando com outras mercadorias em equilíbrio, quanto maior o tempo de produção, maior o preço. Evidentemente, uma unidade individual de uma mercadoria pode ser vendida acima ou abaixo desse preço. E, como se trata de uma média, também é necessário que a mercadoria seja produzida em grande escala para que a lei do valor se manifeste dada a sua natureza estatística.

Evidentemente, não quero dizer que o valor em si seja estatístico, mas sim que a maneira como o valor se manifesta nos preços só pode ser compreendida como uma tendência estatística: ele se torna visível nas médias, em toda a produção em grande escala e ao longo do tempo. Este é um tema longe de ser consensual.

Portanto, a lei do valor não nos diz nada sobre o preço pelo qual a camisa autografada por Michael Jackson foi vendida em um leilão, pois este trata-se de um bem único, não de uma mercadoria produzida em grande escala. Assim,é importante notar que o preço serve como um sinal da relação entre demanda e oferta. Por exemplo, se uma mercadoria não está sendo produzida em quantidade suficiente para atender à demanda, isso causará uma distorção de preço, fazendo com que ela seja vendida muito acima de seu preço médio de equilíbrio. Por outro lado, a superprodução levará à venda abaixo da média.

É nítido então que o marxismo não ignora a distorção nos preços causado pela relação entre oferta e demanda de uma mercadoria, apenas a insere em um contexto mais amplo. É aqui também que surge as maiores críticas ao comunismo, pois defensores do capitalismo alegam que o livre-mercado, através da variação dos preços, é a melhor ferramenta que temos para processar as informações sobre que mercadorias as sociedade demanda e que recursos temos para produzir essas mercadorias, ou seja, para gerenciar a produção de bens.

Evidentemente como comunista somos levados a discordar. Por um lado, a demanda não representa o interesse genuíno da sociedade  por um determinado bem, apenas a demanda da fração da sociedade que tem condições de pagar por este bem. E por outro lado, os avanços tecnológicos que foram realizados pelo próprio capitalismo nos últimos 2 séculos tornam esta questão um tanto ultrapassada. 

Em um cenário de baixa tecnologia, o livre mercado é uma ferramenta rudimentar mas eficiente para estimar a demanda. Se eu produzo uma mercadoria e não tenho outros recursos disponíveis, posso observar a variação no preço dos meus concorrentes para estimar a demanda no mercado, ou até mesmo ir ajustando empiricamente o preço das minhas mercadorias conforme a relação entre oferta e demanda que encontro. Porém, é fato conhecido que hoje em dia há diversos recursos que podem ser utilizados para entender a demanda real por uma dada mercadoria na sociedade. Seria inocência acreditar que as maiores empresas capitalistas dependem apenas de observar a relação entre demanda e oferta para organizar sua produção e precificar seus produtos.

Outra maneira de entender a lei do valor é pensar nos produtos digitais. Para que o preço médio de uma mercadoria tenda para o seu valor, conforme descrito por Marx, é necessária a concorrência. Quando a concorrência é restrita, como no caso do monopólio, a expressão monetária da lei do valor não desaparece, mas fica distorcida. Um jogo digital, uma vez produzido, pode ser copiado infinitamente a custo zero. Naturalmente, se a empresa A produz o jogo X e o vende a qualquer preço inicial, a empresa B poderia copiá-lo e vendê-lo mais barato. Eu mesmo poderia copiá-lo e distribuí-lo gratuitamente. Nesse cenário, em que a lei do valor poderia operar sem barreiras, o preço tenderia a zero, consistente com o fato de que o tempo médio de produção por cópia tenderia a zero. No entanto, isso não acontece na realidade. Longe de negar a lei do valor, o que acontece é que as empresas de jogos impedem legalmente esse cenário por meio da propriedade intelectual. São criados então obstáculos que impedem outras empresas de vender o jogo, o que garante um monopólio e distorce a lei do valor.

Além disso, isso parece ser uma boa ilustração da afirmação de que o desenvolvimento das forças produtivas leva a uma mudança nas relações de produção, uma vez que elas se tornam uma restrição. Foi durante o capitalismo que desenvolvemos as forças produtivas a um nível em que podemos desenvolver softwares complexos, alcançando níveis sem precedentes de automação e integração de dados na história da humanidade. Estamos discutindo questões como a Indústria 4.0 e até mesmo a 5.0. Hoje em dia, podemos distribuir esses produtos de software infinitamente a um custo praticamente nulo. A distribuição pode ser gratuita e aberta, permitindo que todos não apenas tenham acesso às ferramentas mais avançadas, mas também contribuam para o desenvolvimento futuro. Mas isso não interessa à classe que detém o poder político; ela usa sua posição privilegiada nas relações de produção para criar restrições burocráticas que impedem o desenvolvimento e o uso dessas ferramentas e, consequentemente, afetam negativamente a produção.

A discussão sobre a automação do processo de produção e sua relação com a teoria do valor não é novidade.

\begin{quote}
Uma terceira e última prova da correção da teoria do valor-trabalho é a prova por redução ao absurdo. Além disso, é a mais elegante e “moderna” das provas.

Imagine por um momento uma sociedade na qual o trabalho humano vivo desapareceu completamente, ou seja, uma sociedade na qual toda a produção foi 100\% automatizada... Mas imaginemos que esse desenvolvimento foi levado ao extremo e o trabalho humano foi completamente eliminado de todas as formas de produção e serviços. O valor pode continuar a existir nessas condições? Pode haver uma sociedade em que ninguém tenha renda, mas as mercadorias continuem a ter valor e a ser vendidas? Obviamente, tal situação seria absurda. Uma enorme quantidade de produtos seria produzida sem que essa produção gerasse qualquer renda, uma vez que nenhum ser humano estaria envolvido nessa produção. Mas alguém iria querer “vender” esses produtos para os quais não havia mais compradores!

É óbvio que a distribuição de produtos em tal sociedade não seria mais efetuada na forma de venda de mercadorias e, na verdade, a venda se tornaria ainda mais absurda devido à abundância produzida pela automação geral. Em outras palavras, uma sociedade na qual o trabalho humano fosse totalmente eliminado da produção, no sentido mais geral do termo, incluindo os serviços, seria uma sociedade na qual o valor de troca também teria sido eliminado. Isso prova a validade da teoria, pois, no momento em que o trabalho humano desaparece da produção, o valor também desaparece com ele. 

(Introdução à Teoria Econômica Marxista)
\end{quote}

E eu acrescentaria que, juntamente com a lei do valor, a própria troca de mercadorias também desapareceu. A única maneira de continuar desenvolvendo forças produtivas livres dessas restrições é através da própria superação das relações de produção: a revolução. Isto é um tópico especialmente importante nos dias atuais com o desenvolvimento da Inteligência Artificial e da automatização. 

É também compartilhando a partir desta linha de raciocínio que Coggiola em “Teoria Econômica Marxista: uma Introdução” declara que a semente da superação do capitalismo está contida no próprio capitalismo. O aumento do capital constante (custos como instalações fixas, máquinas, terrenos, prédios, matérias primas, etc) em razão do capital variável (os custos com a mão de obra) devido aos constantes desenvolvimento tecnológicos empurraria a humanidade para uma situação propícia pela superação do capitalismo.

Erik Wright em “Como ser um anti-capitalista no século XXI” faz uma discussão parecida apontando sobre a precarização do trabalho e do salário de grande parte dos trabalhadores e como isso poderia forçar a instituição de uma renda básica universal como forma de conter a crise. Independente da estratégia adotada pela sociedade para lidar com as consequências do desenvolvimento da automação, é evidente que o desenvolvimento das forças produtivas no capitalismo caminham em uma direção de se chocar cada vez de forma mais violenta com as próprias relações de produção que definem o capitalismo, uma vez que o aumento de produtividade do capitalismo faz com que cada vez seja preciso menos trabalho para produzir as condições materiais necessárias para o desenvolvimento da sociedade.

Em uma sociedade de produção em abundância, o capitalismo não faz sentido. O capitalismo perde o seu grande “mérito” de ser o mais eficiente motor que a humanidade já construiu para maximizar a produtividade, uma vez que não há mais esta necessidade. Para aumentar sua própria sobrevivência, o capitalismo adota diferentes estratégias, como a expansão das suas fronteiras espaciais quando possível, e quando não mais possível, a destruição da base material através de guerras.

Até aqui, eu discuti apenas a manifestação monetária do valor dado pelo preço da mercadoria, ou de forma mais geral, a sua manifestação quantitativa dada pela quantidade de horas socialmente necessárias na produção da mercadoria. Mas o valor não se reduz a isso.

\begin{quote}
O valor de uso e o valor não são, como sustenta Bohm-Bawerk, duas diferentes propriedades das coisas. O contraste entre ambos é provocado pelo contraste entre o método das ciências naturais, que trata a mercadoria como uma coisa, e o método sociológico, que trata de relações sociais “fundidas com coisas". "O valor de uso expressa uma relação natural entre uma coisa e um homem, a existência de coisas para o homem. Mas o valor de troca representa a existência social das coisas"...

(Isaak Rubin em Teoria marxista do valor)
\end{quote}

Algo que não havia ainda sido comentado é o caráter social do valor. Isto é, o valor de uma mercadoria não é uma propriedade intrínseca da mercadoria. Não é como o peso de um bem material, que é uma propriedade física desse material. Mas uma mercadoria só adquire valor dentro de um contexto social adequado. Até mesmo sua expressão quantitativa como a quantidade socialmente de horas necessárias, exige considerar a sociedade no qual a produção está inserida e conecta os produtores entre si, uma vez que o valor não é resultado da produção isolada deste produtor. Também conecta o produtor a possíveis consumidores, afinal a relação entre oferta e demanda não vai afetar apenas a expressão monetária do valor, mas a própria existência do mesmo.

Se você plantar uma laranjeira apenas para consumo pessoal, não há valor de troca nesta laranja. Se ninguém ver uso nesta laranja, mesmo que você tenha produzido para troca, não há valor de uso. Em ambos os casos, o valor cessa de existir. 

É preciso ter isto em mente de forma bastante clara, que o que determina o capitalismo são as relações sociais de produção, que são em si relações sociais entre indivíduos. Como visto, o próprio valor de uma mercadoria é fruto de uma relação social.

Em uma economia socialista podemos utilizar como métrica para organizar a produção a quantidade média de horas necessárias para que uma dada mercadoria seja produzida, isto é proposto por Cockshott em “Towards a new socialism”, mas isto não é mais o valor expresso na teoria do valor marxista, pois isto diz respeito especificamente à uma relação que existe no modo de produção capitalista.

Esta natureza social do capitalismo é constantemente ocultada pela cultura hegemônica capitalista. Esquecemos que a sociedade se organiza através da relação entre pessoas, e passamos a pensar que a sociedade se manifesta pela relação entre pessoas e coisas, ou muitas vezes, até mesmo apenas pela relação entre coisas. Aquilo que tem uma natureza social, passa a ser vista como se tivesse uma existência autônoma e é naturalizado como algo imutável e eterno, desde o valor de uma mercadoria, até a própria existência do capitalismo.

Esta é uma discussão sobre algo que constantemente ganha o nome de fetiche. Eu pensei em escrever uma sessão separada, mas acredito que está intrinsecamente relacionado com a teoria do valor, assim como que eu não tenho conhecimento suficiente que justifique uma sessão própria, então manterei aqui.

O próprio Marx em seus escritos mais tardios, como O Capital, muito pouco usou este termo, então para evitar ser crucificado pelos especialistas no assunto emitirei este rótulo e me concentrarei em tentar discutir o fenômeno que eu entendo que está por trás disto. 

Quando pensamos no nosso exemplo da laranja, pudemos observar que uma simples produção visando a troca no mercado conecta todos produtores e consumidores entre si diretamente, tanto para a existência do valor quanto para a determinação do valor seja na sua expressão monetária em preço na sua sua expressão quantitativa em horas socialmente necessárias. Podemos ainda pensar que de forma direta ou indireta o produtor também se relaciona aos fornecedores de insumos necessários para a produção de laranja e de outras frutas, assim como a demanda por outras frutas no mercado. Podemos seguir este raciocínio de forma indefinida, não é difícil perceber que em última análise toda a sociedade está conectada.

Esta relação tende a ser ocultada no capitalismo e a relação entre pessoas nos parece uma relação entre coisas. Que uma troca entre um saco de 1kg laranja por uma moeda de R\$1, é uma relação entre as moedas e as laranjas apenas. Que 1kg de laranja vale 1 real pelas propriedades físicas da laranja e das moedas. Mas isso oculta toda natureza social do sistema de produção e as relações sociais envolvidas neste processo.

Mas isso não é um mero acaso que ocorre no capitalismo, isto decorre da forma com que o capitalismo organiza a sociedade. Pois ainda que seja determinado pelas relações sociais, a posição que as pessoas ocupam nesta arquitetura social do capitalismo é determinada pelas coisas que as pessoas possuem. Uma pessoa se torna capitalista por possuir capital. É também esta característica da natureza do capitalismo que é explorada pela cultura hegemônica que visa difundir entre os trabalhadores a perspectiva de que eles próprios, independente da sua origem, possam vir a ser capitalistas. Pois afinal agora, diferente de outros modos de produção que muitas vezes a produção que ocupavam na sociedade era definida predominantemente por relações sanguíneas, agora sua posição é definida pelos bens em sua posse. Ou seja, qualquer um, independente de sua origem, que passe a possuir capital, passa a ocupar a posição de capitalista.  Isto é, se por um lado há uma objetificação das relações sociais, há uma humanização das coisas.

Evidentemente na prática, a posse das coisas e a mobilidade social é largamente herdada de nossos antepassados, uma vez que como discutido anteriormente, o capitalismo é, antes de tudo, uma forma de organizar as relações sociais de produção. E isso nos permite conectar toda esta discussão da dimensão econômica da sociedade com uma discussão da sua dimensão cultural.

Rubin discutindo a forma social do capitalismo é obrigado a falar do estabelecimento de uma nova cultura, ainda que não intencionalmente.

Esta objetivação, ou "reificação" das relações de produção entre as pessoas sob a forma social de coisas, dá ao sistema econômico maior durabilidade, estabilidade e regularidade. O resultado é a "cristalização" das relações de produção entre as pessoas.

\begin{quote}
Quando o tipo de produção dominante era a produção artesanal, na qual o objetivo era a "manutenção" do artesão, este ainda se considerava um "mestre artesão" e considerava seu rendimento como a fonte de sua "manutenção", mesmo quando expandia sua empresa e já tinha se tornado, em essência, um capitalista que vivia do trabalho assalariado de seus operários. Ele ainda não considerava seu rendimento como o "lucro" do capital, nem seus meios de produção como "capital".

(Isaak Rubin em Teoria marxista do valor)
\end{quote}

Isto é, tanto da objetificação das relações sociais quanto da humanização das coisas não emergem de uma única ou poucas trocas isoladas. De fato, o mercado e a prática de trocas de mercadorias, inclusive com o uso de uma mercadoria comum ou dinheiro existiu desde muito antes do estabelecimento do modo de produção capitalista. Estas práticas precisam ocorrer com certa frequência e intensidade, para que só então ocorra de fato uma “reificação” ou “cristalização” das relações de produção. É apenas neste cenário que se estabelece então uma nova cultura hegemônica onde novos signos, valores e práticas passam a dominar a sociedade, consequentemente organizando também um novo modo de vida.

O que estou dizendo é que a cultura não vem depois e nem antes do estabelecimento de um novo modo de produção, ela vem junto. Não é uma questão simples, você pode impor violentamente as relações sociais correspondentes a um novo modo de produção enquanto a cultura hegemônica ainda não se transformou, mas essa incompatibilidade entre a cultura hegemônica e estas relações cria um conflito que força que um ou outro se modifique.

O estabelecimento de um novo modo de produção a longo prazo só é possível com o estabelecimento de uma nova cultura hegemônica correspondente. Isso decorre do fato de que não existe uma relação social que ocorra fora da arquitetura social da sociedade, não existe relação social fora da dimensão cultural. Assim sendo novas relações sociais carregam consigo uma nova cultura. Esta cultura é determinada pelo plano econômico, mas assume uma relação dialética com este. Novas formas sociais são ao mesmo tempo novas formas de vida.

Concluindo então, entendo que a cristalização de novas relações sociais de produção implica simultaneamente a constituição de uma cultura hegemônica prática, na medida em que tais relações só conseguem existir enquanto modos socialmente reconhecidos, reiterados e vividos de agir, perceber e se orientar no mundo por pelo menos uma fração significativa da população que detém o monopólio da força. Se lido corretamente, o que Rubin descreve por um lado é uma descrição materialista da formação de uma nova cultura hegemônica determinada pela base material da sociedade, e por outro lado, da importância da dimensão cultural no estabelecimento a longo prazo deste novo modo de produção.

\subsection{Religião}

A religião talvez seja a manifestação cultural mais lembrada quando se discute Marxismo, assim sendo, quero fazer uma rápida discussão sobre a relação entre marxismo e religião, pois este foi talvez o primeiro tópico que me fez realmente estudar um pouco mais de marxismo, e ainda é um tema que costuma trazer bastante controvérsia. 

Partirmos da ideia de que o materialismo dialético é uma concepção de mundo, precisamos entender então que toda religião também é uma própria concepção de mundo. Partindo da definição de filosofia como uma unidade entre concepção de mundo e método de investigação, Spirkin vai nos dizer que a história do pensamento filosófico é uma história de disputa entre diferentes concepções de mundo. Assim fica nítido a origem do conflito entre marxismo e religião.

\begin{quote}
Toda a história do pensamento filosófico é, de fato, uma disputa entre várias visões de mundo…

(Dialectical Materialism)
\end{quote}

É evidente que duas concepções de mundo que entram em conflito não podem estar corretas ao mesmo tempo. Logo é uma consequência natural que exista uma disputa de natureza filosófica tanto entre diferentes filosofias, quanto entre diferentes religiões, e também entre diferentes filosofias e religiões.. É o próprio Marx, nos Manuscritos Econômico-Filosóficos vai dizer:

\begin{quote}
A grande realização de Feuerbach é … ter mostrado a filosofia nada mais ser do que a religião trazida para o pensamento e desenvolvida por este, devendo ser igualmente condenada como outra forma e modo de existência da alienação humana;

(Marx nos Manuscritos Econômico-Filosóficos)
\end{quote}

Porém devemos agora desenvolver uma outra discussão. O que é religião? Essa é uma discussão bastante profunda e eu não tenho condições de concluir. É um fato histórico que esta discussão assume normalmente como referência as maiores religiões do mundo, e em particular as abraâmicas. Mais particularmente ainda o cristianismo. Hegel (Lições sobre a filosofia da religião), Feuerbach (A Essência do Cristianismo) e Engels (Contribuição Para a História do Cristianismo Primitivo), por exemplo, centraram suas discussões de filosofia da religião no cristianismo.

A compreensão exata de como funciona o materialismo é complexa, pois não se trata de um reducionismo simplista do tipo “só existe matéria”, mas sim da ideia de que a realidade fundamental é constituída por processos materiais regrados em constante transformação. Esses processos são responsáveis pela formação tanto do mundo físico quanto da consciência. A percepção e a consciência emergem das interações materiais, sem existir independentemente da matéria. A matéria, nesse sentido, não é uma entidade estática ou sólida, mas um conjunto dinâmico de processos em desenvolvimento contínuo. É essa base material, com sua organização dinâmica e regrada, que permite o conhecimento objetivo do mundo.

É evidente que esta concepção de mundo conflita diretamente com as visões mais tradicionais de religião que estamos acostumados. Nosso senso comum é acostumado a utilizar o catolicismo como uma régua no qual comparamos diferentes religiões para definirmos o que é ou não religião. Ressaltamos aqui a importância do fato de que para discutir a relação do marxismo com a religião, primeiro precisamos definir o que é marxismo - ao qual em grande parte este trabalho se dedicou - e depois definir o que é religião.

Porém, até onde meu conhecimento se limita, por “religiões”, na crítica marxista à religião, temos frequentemente o modelo cristão de religião em mente. Logo, quanto mais distante estivermos deste modelo de religiosidade, provavelmente maior a probabilidade da crítica não ser precisa.

Na minha experiência pessoal, eu percebo que tradicionalmente os marxistas quando criticam a religião, concentram suas críticas em três elementos que estão presentes na forma mais popular de cristianismo:

\begin{enumerate}
    \item A existência de um deus pessoal que transcende a matéria e existe independente dela. Isto é, uma consciência incorpórea com intenção e capacidade de criar, destruir e interferir na matéria existente de forma totalmente livre;
    \item A existência de alma, espíritos, fantasmas e afins concebidos como consciências que existem de forma independente à matéria.
    \item Um começo e um fim do universo. 
\end{enumerate}
Os dois primeiros pontos dizem mais (mas não exclusivamente) respeito a concepção materialista do mundo, pois ambos pontos exigem que seja concebido a possibilidade da consciência existir de forma independente e/ou anterior da matéria, além de conceber uma consciência (deus) com capacidade de interferir no mundo de forma independente de qualquer regra que regule o movimento da matéria. 

O terceiro ponto eu entendo se relacionar mais (mas não exclusivamente) à dialética. O universo não poderia ter início nem fim, pois a ideia da criação ou do fim do universo, representaria um momento antes e um tempo do próprio universo, e neste momento então não haveria mais mudanças. O que de certa forma, conflita com o que está no coração da dialética. Evidentemente tem uma discussão profunda sobre o tempo surgir com o Big Bang, portanto levanta-se o argumento de que não existiria um antes, mas não cabe aqui entrar em detalhes, apenas apontar um ponto de potencial conflito.

Mas agora eu quero levantar uma polêmica e separar a discussão sobre o desenvolvimento da discussão sobre a criação do mundo, e mais que isso, hierarquizar as questões. Certamente eu discordo do esforço de reduzir o materialismo dialético a apenas questões históricas e sociais, defendo que o materialismo dialético é uma concepção de mundo que inclui também o mundo natural. Porém talvez a questão principal de que o materialismo no marxismo busca discutir, não seja exatamente sobre a matéria ser anterior à consciência. Historicamente a questão tem sido colocada dessa forma, mas talvez este não seja o ponto realmente relevante. 

O que realmente nos interessa é que o mundo material hoje seja independente de qualquer consciência. Isto é, o materialismo descreve um mundo que evolui seguindo leis naturais e impessoais. É essa impessoalidade das leis que regem o mundo que nos permite entendê-lo e ter uma certa previsibilidade sobre sua evolução, é deste fato que podemos obter uma capacidade de transformar o mundo. Certamente a existência de um Deus intervencionista rompe com tudo, pois ficamos à mercê da vontade divina, e claro, o cristianismo tradicionalmente incorpora ambas ideias: um deus criador e intervencionista, mas ambos não se apresentam necessariamente sempre juntos em qualquer tradição, e na minha opinião, merece uma discussão mais cuidadosa.

Levando meu ceticismo ao extremo, se você defender a posição de que uma consciência criou o mundo e nunca mais interferiu, surge um problema sobre de onde veio essa consciência? Como ela pode existir sem que exista um mundo? Recuando um pouco mais, se você defender que uma “força superior” criou o mundo, pode ser que tenha um problema com a dialética porque pressupomos um mundo que está sempre se transformando, então como pode existir um tempo anterior ao mundo ser criado?  Este parece ser um cenário onde nada muda.

Mas de fato, esta posição não tem mais problemas que um ateu falar que o mundo surgiu com no Big Bang. Normalmente a resposta é no sentido de que o próprio tempo surgiu no Big Bang, de forma que não faz sentido falar sobre como era o mundo 'antes'. Verdade seja dita, defender que o 'mundo sempre existiu e está sempre se transformando' também não é um fato incontestável com evidências finais. Esta posição, também levanta perguntas, por exemplo, como pode algo 'existir pra sempre'? 

Por isso me parece que desde que se assuma que não existe nenhuma entidade interferindo no desenvolvimento do mundo, isto é, se assuma o materialismo dialético como princípio sobre como o mundo evolui, talvez 'o mito da criação' não seja uma questão tão relevante para a concepção de mundo comunista. Talvez esta questão não seja mais importante do que discussões como sobre o pós-vida,  que não é uma questão relevante. Ou seja, o primeiro ponto sobre uma divindade capaz de interferir no desenvolvimento do mundo, é indiscutivelmente central para concepção do mundo comunista, já o terceiro ponto sobre a criação do mundo, talvez já ocupe uma posição secundária. Já o segundo ponto, de certa forma é apenas uma variação do primeiro considerando  que estes “espíritos” estejam livres de agir de acordo com as leis naturais.

Então, o que acontece se concebermos uma religião onde por definição não há fenômenos sobrenaturais e onde todos os fenômenos são compreendidos como parte da natureza? Onde o universo não tem começo nem fim. Ainda é uma religião? E se os espíritos forem concebidos de forma que que não existam independente da natureza e da matéria? Seja por dependerem da existência da matéria que conhecemos, ou até talvez por serem constituídos de um outro tipo de matéria ainda a ser descoberta. Estas suposições podem soar estranhas para alguns leitores, mas nada impede que sejam assumidas como verdadeiras por alguém, e se nos depararmos com esta pessoa, o que iremos dizer a ela?  Ainda é religião? É compatível com o marxismo?

\begin{quote}
A própria tentativa de definir a religião, de encontrar alguma essência ou conjunto de qualidades distintivo ou possivelmente único que distinga o religioso do resto da vida humana, é principalmente uma preocupação ocidental. A tentativa é uma consequência natural da disposição especulativa, intelectualista e científica do Ocidente. É também o produto do modo religioso ocidental dominante, o que é chamado de tradição judaico-cristã ou, mais precisamente, a herança teísta do judaísmo, do cristianismo e do islamismo. A forma teísta de crença nesta tradição, mesmo quando rebaixada culturalmente, é formativa da dicotômica visão ocidental da religião. Ou seja, a estrutura básica do teísmo é essencialmente uma distinção entre uma divindade transcendente e tudo o mais, entre o criador e sua criação, entre Deus e o homem.

(Enciclopédia de Religiões MacMillan)
\end{quote}

Eu não trago respostas. Minha intenção é apontar que a questão é mais complexa do que é tradicionalmente pensado. Como não somos capazes de encontrar uma definição precisa de religião, eu percebo que nos resta são duas opções: cabe então aos marxistas criar uma definição da religião que será criticada. Assim podemos delimitar melhor a que manifestações religiosas nossa crítica se refere. 

Ou então, simplesmente tratamos cada religião como uma tradição independente e avaliamos então seus credos ponto a ponto. Assim podemos discutir especificamente onde a concepção de mundo de determinada religião converge ou diverge da concepção de mundo comunista. Esta prática é independente de qual tipo de tradição estamos avaliando. Seja religião, filosofia, ou qualquer outro nome que queiramos dar.

Em resumo, a minha leitura do tópico é que se não estamos em condições de afirmar que o materialismo dialético se opõe necessariamente a toda religião (pois não definimos o que é uma religião) podemos dizer que ele limita quais conceitos religiosos são compatíveis com o marxismo. Dessa forma o que nos resta é discutir que elementos de determinadas religiões entram em conflito com o marxismo, analisando cada elemento de forma particular. Lembrando que este é um trabalho extremamente exaustivo pois mesmo dentro do que tradicionalmente é concebido como uma única religião, encontramos uma grande variedade de concepções de mundo diferentes.

\subsection{Valores}

Originalmente esta seção era uma transcrição/resumo de uma seção do livro "Como ser anticapitalista no século XXI" de Erik Wright. Mas eu vou tentar resumir e ainda mais pois acredito que este é um debate muito importante que precisa ser feito.

A esta altura, já deve estar óbvio o quão importante a cultura é para mim. Podemos lembrar que estou definindo a cultura como um conjunto de signos, valores e práticas, que se relacionam por reforçar uma dada concepção de mundo, ou seja, uma forma de produzir e reproduzir a vida social. Evidentemente não há fronteiras precisas sobre onde começa e onde termina uma cultura, poderíamos dizer que tanto o liberalismo quanto o conservadorismo promoveriam uma cultura capitalista, assim em certo nível poderíamos agrupar com sub-culturas dentro de uma mesma cultura, mas em outro nível, se apresentam como duas culturas diferentes.

Dito isto, eu não acredito que seja necessário que especifiquemos exatamente as fronteiras de cada cultura, mas acredito que esta é uma definição bastante útil, e aqui eu quero destacar a importância dos \textbf{valores}. Pois as práticas de uma dada cultura estarão preferencialmente alinhadas com os valores que esta cultura promove, assim como os signos terão seus significados em um contato direto com os valores que esta cultura promove.

Como um exemplo banal, podemos pensar como o anarco-capitalismo promove como valor a individualidade. Este valor incentiva uma prática de combate a idéias coletivistas, o que também se reflete tanto nos signos escolhidos, quanto nos juízos que as pessoas que internalizaram essa cultura fazem de determinados signos. Por exemplo, talvez ‘egoísta’ não seja visto de forma tão negativo quanto seria em outra cultura, enquanto ‘coletivista’ se transforme em um insulto, talvez um médico que se recuse a ajudar alguém em situação de emergência não seja visto como uma atitude que demande punição, e o direito de um capitalista gastar seu dinheiro como quiser seja visto como sagrado, independente das consequências que ambas decisões causem ao coletivo.

Sendo assim, eu gosto da atitude do Erik ao propor três valores que ele considera central para uma crítica moral do capitalismo. Isto é, valores que devemos esperar encontrar em qualquer cultura anticapitalista que esteja sendo construída como uma oposição à cultura hegemônica capitalista. Vou chamar estes valores de: igualdade, democracia e comunidade.

Vamos tentar abordar cada um destes valores de forma independente, lembrando que originalmente Erik sempre traz os valores em pares, como “Igualdade/Justiça”, sendo assim eu não estou comprometido em reproduzir a visão original do autor, apenas utilizar seus escritos como auxílio para exibir minha própria visão do assunto, ainda que eu considere que estejamos bastante alinhados neste sentido.

\textbf{Igualdade}

\begin{quote}
Em uma sociedade justa, todas as pessoas teriam amplo e igual acesso aos meios materiais e sociais necessários para viver uma vida plena.

(Como ser um anti-capitalista no século XXI)
\end{quote}

O aspecto central então deste valor, é a garantia de igual acesso aos meios materiais e sociais. Isto é diferente de fornecer meramente apenas a oportunidade para que as pessoas possam ter acesso, este segundo caso se aproxima mais da defesa da (inexistente) meritocracia no capitalismo, onde há a ideia que tudo bem você nascer sem acesso à uma moradia digna se você tiver oportunidade de vir a ter acesso algum dia. Aqui a defesa é explicitamente sobre a igualdade no acesso aos meios materiais e sociais.

A mera oportunidade não basta, pois ao se basear na ideia da meritocracia e recompensa do esforço, nos apoiamos na ideia de que o resultado é responsabilidade nossa, fruto direto do nosso esforço. Porém é impossível que a gente consiga distinguir claramente os resultados no qual somos responsáveis pelos resultados que foram fruto do acaso, ou mesmo, quando olhando para apenas um resultado, conseguir avaliar em que medida resultado se deve ao nosso esforço ou ao acaso. Nesse caso, garantir apenas a oportunidade pode impedir que as pessoas tenham acesso aos meios materiais e sociais apenas por uma questão de sorte.

Agora precisamos entender o que são exatamente os meios materiais e sociais que estamos discutindo. Quanto aos meios materiais, são tradicionalmente mais fáceis de entender. Ainda tenha uma variação de sociedade para sociedade, tipicamente deve incluir coisas como alimentação, moradia, roupas, mobilidade, lazer, saúde, educação, etc…

Já por meios sociais temos uma questão mais complexa.

\begin{quote}
Mas eu incluiria no mínimo estes itens: atividades que tragam realização e sentido para as pessoas, de preferência conectadas ao que denominamos comumente de “trabalho”; intimidade e vínculo social; autonomia, no sentido de controle significativo sobre a própria vida; 

(Como ser um anti-capitalista no século XXI)
\end{quote}

Ou seja, enquanto os meios materiais dizem respeito principalmente aos acesso de cada indivíduo à bens materiais, os meios sociais dizem respeito sobre como o indivíduo se insere na vida social da comunidade. Neste sentido, estigmas sociais ligados à raça, gênero, sexualidade, aparência, religião, linguagem, etnicidade, etc devem ser combatidos. 

Isto decorre, pois mesmo se não impedirem diretamente o acesso aos meios materiais (mas muitas vezes as duas questões estão relacionadas), impedem que o indivíduo viva uma vida plena. E em uma sociedade justa, todos devem ter igual acesso tanto à condições materiais quanto sociais  para ter uma vida plena.

O que nos conduz a uma questão sobre o que significa uma vida plena. 

\begin{quote}
Uma vida plena é aquela na qual as capacidades e os talentos individuais se desenvolveram de tal forma que lhes é permitido buscar seus desejos, de modo que, num sentido mais amplo, conseguiram realizar tanto seu potencial quanto seus propósitos. 

(Como ser um anti-capitalista no século XXI)
\end{quote}

Ou seja, é quando o indivíduo é capaz de desenvolver seus potenciais e realizar seus desejos. Evidentemente, este valor de igualdade não é sobre todo mundo viver igual, mas garantir que todo mundo tenha igual acesso aos meios sociais e materiais que são necessários para ter uma vida plena.

Por fim, antes de encerrar a discussão sobre igualdade, podemos conectar esta discussão também à questão da sustentabilidade ambiental.

\begin{quote}
As gerações futuras devem ter, no mínimo, o mesmo acesso aos meios materiais e sociais necessários para obterem vidas plenas que a geração atual possui.

(Como ser anti-capitalista no século XXI)
\end{quote}
Acredito que nada mais é preciso acrescentar.

\textbf{Democracia}

\begin{quote}
Em uma sociedade totalmente democrática, todos teriam amplo e igual acesso aos meios necessários de participar na tomada de decisões substantivas sobre aquilo que afeta suas vidas.

(Como ser um anti-capitalista no século XXI)
\end{quote}

Assim como igualdade estava combinada com justiça, aqui democracia aparece também combinado com liberdade. E assim como antes a justiça se dá através da igualdade, aqui a liberdade se dá através da democracia.

Liberdade então está intrinsecamente conectada à ideia da democracia. Podemos chamar de liberdade a possibilidade que eu tome decisões sobre as questões que me afetam sem interferência externa, e democracia é quando as pessoas que serão afetadas pela minha decisão tomam parte da decisão tanto quanto eu. Isto é, democracia é sobre eu não poder tomar decisões que afetem o outro sem a participação deste outro.

Assim sendo, este valor é centrado na ideia de que as pessoas devem ter o máximo possível de controle sobre suas vidas. Em termos econômicos então, podemos dizer que é por exemplo, que a sociedade esteja em condições de decidir de forma coletiva o que será produzido e como esta produção será distribuída.

Podemos ver claramente como o capitalismo não obedece este valor, uma vez que no capitalismo são os capitalistas que detêm o poder econômico sobre o que será produzido. E mais que isso, a concentração de riqueza que conduz a concentração de poder política possibilita que os capitalistas tomem decisões sobre toda sorte de diferentes assuntos que afetam a sociedade como um todo.

Neste arranjo capitalista então, é fácil também notar como o valor da igualdade previamente discutido também falha no capitalismo. Uma vez que sob a desculpa da meritocracia (que falha desde o princípio, uma vez que nem mesmo a distribuição de oportunidades é igualitária no capitalismo) o capitalismo condiciona o acesso aos bens à posse de uma determinada quantidade de riqueza. O valor da igualdade se opõe diretamente à mercantilização da vida.

E é importante notar que esta concentração de riqueza não deforma apenas a distribuição dos bens materiais, mas também deforma a influência que a sociedade como um todo é capaz exercer na determinação do que será produzido, uma vez que a “lei” da oferta e procura de bens só levará em conta o desejo daqueles que têm dinheiro para comprar este bem, excluindo do processo o restante da sociedade.

Não é uma discussão nova a relação do capitalismo com problemas estruturais de natureza social, como o racismo e o machismo, e talvez de forma ainda mais explícita, também o  etarismo e o capacitismo. Estes últimos expõem radicalmente a lógica capitalista de mercantilização até mesmo dos nossos corpos, pois revela como a fração da sociedade que é vista como menos produtiva, perde sua utilidade diante do capitalismo e consequentemente até mesmo seu valor como ser humano, uma vez que falham em contribuir para lógica capitalista da maximização da produtividade.

Assim sendo, a forma com que o capitalismo se organiza viola de forma intrínseca tanto o valor de igualdade quanto o de democracia discutidos até aqui, mais adiante poderemos ver que de fato, o capitalismo não é compatível com nenhum dos três valores expostos.

Evidentemente, na prática, qualquer decisão e ação de uma pessoa afeta os outros, vivemos em sociedade. E ao mesmo tempo, é impossível que participemos de todas as decisões de todo mundo. Portanto é necessário definir um conjunto de regras que define uma fronteira socialmente aceita entre a vida pública e a vida privada, isto é, entre a democracia e a liberdade. 

A ideia é que para aquelas ações que ocorrem na esfera da vida privada, sejamos livres para fazer o que quisermos. Enquanto as ações que ocorrem na esfera da vida pública passam demandem a participação  de todos aqueles que são afetados direta ou indiretamente por nossas decisões. 

Evidentemente onde se situa esta fronteira entre público e privado não é algo dado, nem simples. É algo que deve ser criado por meio de um processo social, e sujeito a modificações tanto no espaço quanto no tempo, através do próprio processo democrático. Evidentemente, também, os valores de igualdade e democracia estão intrinsecamente relacionados, por exemplo, é preciso que as pessoas tenham os meios sociais necessários para que consigam participar do processo democrático de forma adequada. Isto é, apesar de estarmos falando em três valores centrais para construção de uma cultura anticapitalista, estes três valores não estão isolados, mas formam um todo coerente onde os valores reforçam uns aos outros.

\textbf{Comunidade}

\begin{quote}
Comunidade/solidariedade expressa o princípio pelo qual as pessoas devem cooperar umas com as outras não apenas por aquilo que recebem individualmente, mas por comprometimento real com o bem-estar dos outros e por um senso de obrigação moral de que isso é o certo a ser feito.

(Como ser anti-capitalista no século XXI)
\end{quote}

Seguindo os exemplos anteriores, aqui é dito que a comunidade se realiza através da solidariedade. A solidariedade aqui é entendida como uma cooperação que não é motivada exclusivamente por preocupações com nós mesmos,  mas por meio de uma combinação entre obrigações morais e preocupação com os outros. E por sua vez, falamos então de comunidade como valor quando essa cooperação ocorre em atividades mundanas que se dão no cotidiano no interior de uma comunidade.

Os valores de igualdade e democracia se realizam de forma mais eficaz em uma sociedade que estime o valor da comunidade. É mais fácil você garantir acessos iguais aos meios materiais e sociais, assim como garantir que todos afetados diretamente e indiretamente tenham voz em todas questões importantes que os afetam quando cultivamos conjuntamente valores de comunidade e solidariedade. Evidentemente este valor se aplica a qualquer unidade social em que as pessoas convivem e cooperam reciprocamente. 

Como nos casos anteriores, aqui também não deve ser surpresa que o capitalismo falhou em realizar este valor. Ou ainda indo além, o capitalismo não apenas falha em promover este valor, mas promove valores contrários. Se nós, como comunistas, queremos promover a comunidade e a solidariedade, o capitalismo precisa se apoiar necessariamente no individualismo. Se como comunistas exaltamos a cooperação, o capitalismo exige que a competição seja vista como algo moralmente superior, ou pelo menos, naturalizada.

Isto não é fruto do acaso. Como o comunismo visa produzir uma sociedade que atenda o bem estar da sociedade, é natural que ela se apoie em valores como igualdade, democracia e solidariedade, que promova um campo fértil para a cooperação. Já o capitalismo, como tem como meta a maximização da produção e do lucro, exige o fomento da competição. Este cenário de intensa competição fomenta o desenvolvimento de um cenário de atomização da vida social onde a vida comunitária é destruída.

Assim o capitalismo precisa se apoiar em valores como individualismo e a própria desigualdade. Isto é necessário uma vez que o capitalismo tem como pressuposto a própria desigualdade na divisão entre aqueles que detém os meios de produção e aqueles que tem que vender sua força de trabalho para sobreviver. Este arranjo social também só pode sobreviver se a própria noção de democracia for corrompida de forma que tenha por prioridade a defesa da propriedade privada acima de qualquer outra questão..

Este conjunto de valores mostram-se relevantes para que, entre outras coisas, possamos avaliar que tipo de cultura contra hegemônica queremos construir.

\subsection{Luta Indígena}

Lendo “A Terra dos Mil Povos” (Kaká Werá Jecupé) e “Futuro Ancestral” (Ailton Krenak), eu percebi que o Brasil possui uma contracultura anticapitalista com raízes materiais que antecedem ao próprio capitalismo e eu preciso falar disto. Eu defendo que a principal tarefa atual da militância comunista no Brasil é a disputa pela consciência política dos trabalhadores. E mais do que isso, essa disputa passa pela necessidade da defesa de uma contracultura anticapitalista.

E não precisamos inventar uma cultura nova. O que podemos e devemos é reconhecer uma luta que se desenvolve em nosso solo desde que o colonialismo chegou por aqui.

As culturas indígenas, em toda sua diversidade, já são culturas anti-hegemônicas que têm resistido e se desenvolvido ao longo de séculos de capitalismo. Com o grande diferencial de que estas culturas não são uma cultura artificial importada do estrangeiro sem raízes, são culturas ancestrais. Possuem uma longa história e uma sólida base material.

Nenhuma cultura pode se desenvolver de fato sem raízes. É por isso, que nós, como comunistas, devemos reconhecer e dar a devida importância ao fato de que a dimensão material das culturas dos povos originários possuírem raízes firmes em nossa terra. Não é sem razão que a luta indígena adota o seguinte lema: “\textit{A nossa luta é uma luta ancestral}”.

Assim sendo, não precisamos “partir do zero” para construir uma cultura anticapitalista no Brasil, e muito menos devemos importar uma cultura estrangeira. Devemos (se você não é indígena como eu) nos aliar à luta dos povos originários.

Já faz um tempo que eu carrego a ideia de que nosso imaginário foi colonizado. Defendo que uma das principais violências do capitalismo foi que nos tiraram a capacidade de sequer imaginar um mundo melhor. Esta ideia pode ser lida também em “Futuro Ancestral”, onde Ailton Krenak a expressa da seguinte forma:

\begin{quote}
As pessoas acham mais fácil acabar com o mundo do que acabar com o capitalismo.

(Ailton Krenak em Futuro Ancestral)
\end{quote}

Para lidar com este problema, Ailton Krenak propôs então que “\textit{temos que reflorestar nosso imaginário}”. Em outras palavras, é preciso combater esta concepção de mundo capitalista com uma outra concepção de mundo anticapitalista. Um imaginário reflorestado para combater um imaginário colonizado.

Isto ressoa com tudo que tenho refletido e discutido. A necessidade de fomentar uma contracultura para disputar a cultura hegemônica, passa pela necessidade de defender os povos originários.

De fato, todo marxista sempre fala de que o marxismo não é uma fórmula para ser mecanicamente copiada de um lugar para outro. Sabe que é uma teoria que se desenvolve de forma diferente em diferentes contextos. Mas me parece que ainda encontramos dificuldades para saber exatamente como esse marxismo tipicamente brasileiro deve ser desenvolvido. Hoje, eu tenho que isso passa necessariamente pela da luta índigena;

Os povos originários compõem um movimento de resistência ao capitalismo anterior ao próprio marxismo. Poderíamos pensar na questão de forma superficial, e concluir que deveríamos nos inspirar nesta luta. Mas eu acredito que a união deve ser mais profunda. Não basta apenas nos inspirar, devemos aprender e nos aliar a este movimento, devemos ouvir e dar voz. Me arriscando, mas tendo a pensar que um marxismo adaptado a nossa cultura é um marxismo que respeite e seja absorvido pela luta indígena.

Evidentemente a luta indígena é diversa. As nações indígenas são diversas. Não devemos esquecer disso em momento nenhum. Mas é preciso ser enfático que existe este movimento anticapitalista na luta indígena. Recorrendo novamente a Ailton Krenak, respeitado e importante porta-voz da causa indígena, ele também denuncia o capitalismo com a seguinte declaração:

\begin{quote}
O cancro do capitalismo só admite propriedade privada e é incompatível com qualquer outra perspectiva de uso coletivo da terra”.

(Ailton Krenak em Futuro Ancestral)
\end{quote}

O próprio nascimento da Aliança dos Povos da Floresta é descrito como uma insurgência para eliminar a figura do patrão. Estamos diante de um discurso anticapitalista por excelência, que se nega a fazer coro com qualquer discurso colonial. A lógica capitalista que faz com que as cidades sejam invadidas pela indústria e pela produção transformando a lógica da vida coletiva em vida privada é denunciada ao longo de todo “Futuro Ancestral”.

Um assunto que eu sempre busco discutir, que é a linguagem linguagem, também é abordado por Krenak, Krenak diz ver nela um fator determinante nas relações sociais. Se juntando a nossa denúncia contra a mercadorização da vida promovida pelo capitalismo, em “A Queda do Céu” (Davi Kopenawa Yanomami) os brancos são chamados de “povo da mercadoria”.

O que eu estou tentando realizar aqui é trazer o reconhecimento desta luta para o interior do movimento comunista. Para que o movimento comunista por sua vez busque a luta indígena. Se, como eu defendo, o primeiro campo de batalha contra o capitalismo é cultural, historicamente os povos originários estão entre os primeiros a liderar esta frente de batalha. E nós temos muito a aprender com eles.

Falando sobre o marxismo, especialmente em sua formulação mais tradicional, é conhecido que ele traz uma concepção de mudo própria. Mas os povos originários já trazem em si suas próprias concepções de mundo. Ainda que ambas sejam anti capitalistas, para além disso, até que ponto elas convergem?

Mesmo admitindo a existência de uma convergência quanto a postura anticapitalista não implica que exista uma convergência geral entre a luta indígena e o movimento marxista. Devo me aprofundar nisso futuramente, mas eu gostaria de esboçar uma discussão.

Vamos começar relembrando a típica organização do marxismo em três componentes: política (socialismo), econômica e filosófica. Podemos dizer que até aqui apresentamos, ainda que grosseiramente, uma convergência no eixo político. Pelo menos no que tange ao objetivo de superação do capitalismo

Vamos agora comentar os dois eixos restantes. Eles nos ajudarão a entender melhor as convergências e divergências entre a luta indígena e o movimento marxista. A questão principal é: Admitindo que ambos queiram superar o capitalismo, ambos querem construir a mesma sociedade pós-capitalista?

\textbf{Economia}: entre os dois eixos que restam, acredito que este é o mais simples, por isso opto por começar por ele. Tendo em mente que a economia comunista visa a superação do capitalismo, e que a luta indígena expressa por pensadores como Ailton Krenak é também anticapitalista, temos um sinal favorável de convergência neste aspecto.

A proposta de modo de produção comunista visa colocar diretamente nas mãos dos trabalhadores todo o poder político de gerenciar a produção e a distribuição destes bens produzidos. Assim, os povos originários, em uma sociedade comunista, por definição, devem possuir completa autonomia para decidir como gerir a produção, a distribuição e sua vida como um todo.

Acredito ser um fator determinante de convergência entre as lutas, a reorientação da finalidade da produção. A produção no capitalismo visa a maximização do lucro acima de qualquer coisa. Já o ideal comunista de produção é a maximização do bem estar dos membros da sociedade. Esta mudança de foco na produção me parece alinhar claramente as ideias comunistas à luta indígena.

Pode existir um conflito com uma visão marxista 'industrial' que já foi popular, mas que deve ser denunciada como inadequada. Esta concepção vulgar de marxismo se aproxima de um certo modo de desenvolvimentismo, vê como necessidade que a sociedade persiga um desenvolvimento desenfreado das suas indústrias acima de qualquer coisa.

Esta é uma concepção inspirada em uma espécie de visão de desenvolvimento linear e hierárquica da história da humanidade. Onde quanto mais e maiores as indústrias mais desenvolvida estaria a sociedade.

Apesar de ter gozado de certa popularidade, é preciso ter em mente que esta não é uma posição inerente ao comunismo. Ela vem sendo combatida desde o próprio Marx. Quando acusado de que estaria defendendo uma “inevitabilidade histórica” do desenvolvimento das formas de produção, Marx afirmou ter sido mal compreendido e que o curso de desenvolvimento descrito em “O Capital” não é uma afirmação sobre uma inevitabilidade histórica universal, mas apenas a descrição da gênese da produção capitalista no contexto dos países da Europa Ocidental.

\begin{quote}
A “inevitabilidade histórica” desse curso está, portanto, expressamente restrita aos países da Europa Ocidental.

(Karl Marx em carta para Zasulich de 08/03/1881)
\end{quote}

O que é inerente ao comunismo é uma defesa da total autonomia para que os trabalhadores tenham controle sobre como gerir sua vida. Uma marxista atual que se opõe a esta perspectiva de desenvolvimento inevitável do capitalismo é Ellen Wood. Criticando argumentos que naturalizam o capitalismo e a indústria como resultados inevitáveis da "natureza humana", ela também escreve:

\begin{quote}
Mais particularmente, esses argumentos tendem a reforçar a visão profundamente eurocêntrica de que a ausência do capitalismo, é, de algum modo, uma falha histórica (linha de pensamento que é especialmente contraproducente para os críticos do capitalismo).

(A origem do capitalismo)
\end{quote}

Ainda sobre questões econômicas, devemos citar a abolição do dinheiro. Característica também inerente do comunismo que combinado então com o controle da produção nas mãos dos trabalhadores, cria um cenário bastante flexível quanto a sociedade pós-capitalista. Contrariando essa concepção de exaltação do trabalho industrial e da urbanização desenfreada da vida, podemos recorrer novamente ao próprio Marx:

\begin{quote}
Logo que o trabalho começa a ser distribuído [no capitalismo], cada um passa a ter um campo de atividade exclusivo e determinado, que lhe é imposto e ao qual não pode escapar; o indivíduo é caçador, pescador, pastor ou crítico crítico, e assim deve permanecer se não quiser perder seu meio de vida – ao passo que, na sociedade comunista, onde cada um não tem um campo de atividade exclusivo, mas pode aperfeiçoar-se em todos os ramos que lhe agradam, a sociedade regula a produção geral e me confere, assim, a possibilidade de hoje fazer isto, amanhã aquilo, de caçar pela manhã, pescar à tarde, à noite dedicar-me à criação de gado, criticar após o jantar, exatamente de acordo com a minha vontade, sem que eu jamais me torne caçador, pescador, pastor ou crítico.

(Karl Marx em Ideologia Alemã)
\end{quote}

\textbf{Concepção de mundo}: Este é certamente o eixo onde o conflito surge diretamente. O marxismo historicamente tem se mostrado hostil à religião e se portado como um defensor do ateísmo. Ser também uma concepção de mundo é a principal característica do marxismo que tem se mostrado difícil de adequar a diferentes culturas não-europeias sem recorrer a um novo tipo de colonialismo. E é portanto, o tópico mais sensível e polêmico.

Há diversos exemplos de pensadores que não sendo europeus, apesar de verem o marxismo interessante sob certos aspectos, não conseguem se sentir representados pelo mesmo devido a aspectos da sua própria cultura e concepção de mundo.

Podemos citar por exemplo Fausto Reinaga na Bolívia ou Ali Shariati no Irã. Ambos vão descrever o marxismo como portador de uma concepção de mundo eurocêntrica, de forma que sua imposição nas culturas locais acarretaria em uma nova espécie de colonialismo.

Eu não estou aqui para contrariar nenhum deles, pelo contrário, o que estou propondo é que a solução para este impasse tenha por foco não a adequação da cultura dos povos originários à uma teoria europeia, mas que esta concepção de mundo muito restrita e simplista que historicamente tem sido defendida por alguns marxistas, seja uma abordagem inadequada tal qual a abordagem "desenvolvimentista" discutida anteriormente no plano econômico.

O fato do marxismo tradicionalmente ser visto como necessariamente ateu, me parece nascer de uma leitura simplista tanto do próprio materialismo defendido pelo marxismo, quanto de outras visões de mundo. É importante destacar que o materialismo no marxismo \textbf{não} significa algo tão simples como uma afirmação do tipo de que “só existe matéria”. Significa sobretudo o reconhecimento de que o mundo existe objetivamente e se organiza segundo leis impessoais e naturais. 

Evidentemente isto ainda é polêmico, mas se afasta da leitura mais simplista da questão, de uma simples defesa irrestrita de um ateísmo inconsequente. Possíveis conflitos entre o marxismo e uma concepção de mundo qualquer, deve ser analisada ponto a ponto para entender de forma mais criteriosa exatamente qual é este conflito e como podemos lidar com ele.

De toda forma, é preciso evitar uma leitura simplista de materialismo. E vale dizer, mesmo no marxismo mais tradicional, esta questão já figurava como uma questão secundária. Falando sobre o contexto cristão da Rússia, Lênin disse:

\begin{quote}
A unidade desta luta realmente revolucionária da classe oprimida pela criação do paraíso na terra é mais importante para nós do que a unidade de opiniões dos proletários sobre o paraíso no céu.

(Lênin em O Socialismo e a Religião)
\end{quote}

O que eu busco defender é a ideia de que o marxismo pode ser absorvido pela luta indígena sem adotar um tom paternalista ou condescendente. Meu objetivo com este texto, é incentivar o movimento marxista brasileiro de forma que ele consiga se transformar em uma ferramenta de transformação social empunhada também pelos povos originários.

Se eu for ser acusado de algo, espero que me acusem de deformar o marxismo, mas não por deformar a luta dos povos originários para caber em uma versão eurocêntrica da teoria marxista. Evidentemente eu não tenho uma resposta final. Isto só pode ser desenvolvido de forma coletiva.

O que trago nada mais é que uma simples proposta de reflexão sujeita ao escrutínio público. Ainda que eu não seja indígena, para concluir gostaria de me dar a liberdade de reproduzir algumas palavras famosas de um indígena de minha terra:

Esta terra tem dono!

\textit{Djekupé A Djú}
